\documentclass[12pt]{book}
\usepackage{epigraph} 
\usepackage{mathrsfs}
\usepackage[utf8]{inputenc}
\usepackage{amsmath}
\usepackage{amsfonts}
\usepackage{amssymb}
\usepackage{amsthm}
\usepackage{url}
\usepackage{hyperref}
\urlstyle{rm}
\usepackage[spanish]{babel}
\usepackage{booktabs}
\usepackage{subfigure}
%\usepackage{makeidx}
\usepackage{imakeidx}
%\usepackage[spanish,onelanguage]{algorithm2e}
%\usepackage[ruled,vlined]{algorithm2e} % ou \usepackage{algorithmic} selon ta préférence
\usepackage[spanish]{algorithm} % Charge le package algorithm avec support espagnol
\floatname{algorithm}{Algoritmo} % Change le nom de l'environnement
%

%\usepackage{algorithm}
\usepackage{algpseudocode}
\algrenewcommand\algorithmicrequire{\textbf{Requerido:}}
\algrenewcommand\algorithmicensure{\textbf{Garantizar:}}
% Redéfinir les mots-clés en espagnol
\algrenewcommand\algorithmicfor{\textbf{para}}
\algrenewcommand\algorithmicdo{\textbf{hacer}}
\algrenewcommand\algorithmicend{\textbf{fin}}
%\algrenewcommand\algorithmicendfor{\textbf{FinPara}}
\algrenewcommand\algorithmicreturn{\textbf{Devolver}}

%\usepackage{tocbibind}
\usepackage[nottoc]{tocbibind}

%\usepackage{footmisc}
%\usepackage{perpage}

% Format Curso
%\usepackage[letterpaper,top=3cm,bottom=3cm,left=3cm,right=3cm,marginparwidth=1.75cm]{geometry}
% Format DIRAC
\usepackage[paperwidth=17cm, paperheight=24cm, top=2.7cm, bottom=2.7cm, left=2.0cm, right=2.0cm]{geometry}
%\usepackage[paperwidth=17cm, paperheight=24cm, top=2.5cm, bottom=2.5cm, left=1.5cm, right=1.5cm]{geometry}
%\usepackage[paperwidth=17cm, paperheight=24cm, margin=1.5cm]{geometry}

\usepackage{titlesec}

% Personalizar el formato del capítulo
\titleformat{\chapter}[display]
    %{\normalfont\Large\bfseries}
    {\normalfont\LARGE\bfseries}
    % Texto plano (sin cursiva), grande y en negrita
    {\chaptertitlename\ \thechapter}
    {1ex}
    {} % Sin formato adicional antes del título
    [\vspace{1ex}\titlerule] % Línea horizontal después del título

% Ajustar el espaciado antes y después del título
\titlespacing*{\chapter}{0pt}{0pt}{20pt}

% Configurar el encabezado y pie de página
\usepackage{fancyhdr}
\pagestyle{fancy}
\fancyhf{} % Limpiar encabezados y pies de página existentes
\fancyhead[RE,LO]{\normalfont\leftmark} % Título del capítulo en el encabezado exterior (izquierdo en páginas pares, derecho en páginas impares)
\fancyhead[LE,RO]{\normalfont\thepage} % Número de página en el encabezado interior (derecho en páginas pares, izquierdo en páginas impares)
\renewcommand{\headrulewidth}{0.4pt} % Grosor de la línea de separación del encabezado
\renewcommand{\chaptermark}[1]{\markboth{\thechapter. #1}{}} % Incluir número del capítulo en el encabezado

% Désactiver l'en-tête sur la 2e page (table des matières)
\fancypagestyle{plain}{%
  \fancyhf{}%
  \fancyfoot[C]{\thepage}%
  \renewcommand{\headrulewidth}{0pt}%
}

\usepackage{tikz}
\usetikzlibrary{decorations.pathmorphing}
\usetikzlibrary{arrows.meta, shapes.geometric}

%\usepackage{listings}
\usepackage{xcolor}

\usepackage{listings}
%\renewcommand{\lstlistingname}{Listado} % Change "Listing" à "Listado"
\renewcommand{\lstlistlistingname}{Índice de listados} % Optionnel : pour la table des listados

% Configuration globale pour toutes les lstlistings
\renewcommand{\lstlistingname}{\textbf{\footnotesize Listado}} % Nom en gras
\lstset{
    captionpos=b, % Position de la légende (b = bas)
    basicstyle=\ttfamily\footnotesize, % Style du code
    caption=\textbf{\footnotesize \lstname} % Légende en gras et footnotesize
}

% Configuration supplémentaire (optionnelle)
\lstset{
    basicstyle=\ttfamily\footnotesize,
    breaklines=true,
    captionpos=b, % Position de la légende (b=bas, t=haut)
    frame=single,
    numbers=left,
    numberstyle=\tiny,
    stepnumber=1,
    showstringspaces=false,
    tabsize=2,
}

\lstset{
    basicstyle=\ttfamily,
    breaklines=true,
    inputencoding=utf8, 
    extendedchars=true, 
    literate=%
        {á}{{\'a}}1
        {é}{{\'e}}1
        {í}{{\'i}}1
        {ó}{{\'o}}1
        {ú}{{\'u}}1
        {ñ}{{\~n}}1
}

% Define a Python style
\lstdefinestyle{pythonstyle}{
    language=Python,
    basicstyle=\ttfamily\footnotesize,
    keywordstyle=\color{blue}\bfseries,
    commentstyle=\color{gray}\itshape,
    stringstyle=\color{red},
    numberstyle=\tiny\color{gray},
    numbers=left,
    stepnumber=1,
    breaklines=true,
    frame=single,
    backgroundcolor=\color{black!5},
    captionpos=b
}

\lstdefinestyle{juliastyle}{
    %language=Julia,
    basicstyle=\ttfamily\footnotesize,
    keywordstyle=\color{blue}\bfseries,
    commentstyle=\color{gray}\itshape,
    stringstyle=\color{red},
    numberstyle=\tiny\color{gray},
    numbers=left,
    stepnumber=1,
    breaklines=true,
    frame=single,
    backgroundcolor=\color{black!5},
    captionpos=b
}

\lstdefinestyle{matlabstyle}{
  language=Matlab,
  basicstyle=\ttfamily\footnotesize,
  keywordstyle=\color{blue},
  commentstyle=\color{green!50!black},
  stringstyle=\color{red},
  numbers=left,
  numberstyle=\tiny,
  stepnumber=1,
  breaklines=true,
  frame=single,
  captionpos=b
}

%\newtheorem{theorem}{Theorem}
%\newtheorem{cor}{Corollary}
%\newtheorem{lemma}{lemma}
%\newtheorem{proposition}{Proposition}
%\newtheorem{definition}{Definition}

\newtheorem{theorem}{Teorema}
\newtheorem{cor}{Corolario}
\newtheorem{lemma}{Lema}
\newtheorem{proposition}{Proposición}
\newtheorem{definition}{Definición}
\newtheorem{remark}{Observación}
\newtheorem{ex}{Ejemplo}

\DeclareMathOperator{\envolt}{\mathrm{conv}}

\newcommand{\calF}{\mathcal{F}}
\newcommand{\calG}{\mathcal{G}}
\newcommand{\dcaption}[1]{\caption{\footnotesize \textbf{#1}}}

% Pour changer le titre de la table des matières
\addto\captionsspanish{%
  \renewcommand{\contentsname}{Tabla de contenido}%
}

\usepackage{caption}

\captionsetup[figure]{
    labelfont={bf},      % Met le label (ex: "Figure 1:") en gras
    textfont={bf},       % Met le texte de la légende en gras
    font=footnotesize    % Met la légende en taille footnotesize
}

\captionsetup[table]{
    labelfont={bf},      % Met le label (ex: "Figure 1:") en gras
    textfont={bf},       % Met le texte de la légende en gras
    font=footnotesize    % Met la légende en taille footnotesize
}

%\MakePerPage{footnote}
\usepackage{chngcntr} % Pour contrôler les compteurs
\counterwithout{footnote}{chapter} % Empêche la réinitialisation des notes de bas de page à chaque chapitre

% Cambiar el título de la bibliografía
\addto\captionsspanish{\renewcommand{\bibname}{Referencias bibliográficas}}

\title{Análisis convexo}
%ODEs and PDEs}
\author{Franz Chouly}
\date{En Construcci\'on (\today)}

% Define el índice y cambia el título
\makeindex[title=Índice temático] % Cambia "Índice" por "truc much
%\makeindex

\newcommand{\sob}{\mathcal{H}}
\newcommand{\bes}{\mathcal{B}}
\newcommand{\R}{\mathbb{R}}
\newcommand{\Z}{\mathbb{Z}}
\newcommand{\dist}{\text{dist}}

% --- Comando para productos internos ---
\DeclareMathOperator{\inner}{\langle \cdot, \cdot \rangle}

% --- Comando para normas ---
\DeclareMathOperator{\norma}{\Vert \cdot \Vert}


\begin{document}

\maketitle

\thispagestyle{empty}

% Table des matières : on force le style "plain" pour cette page
\thispagestyle{plain}
% Retour au style "fancy" pour les pages suivantes
\pagestyle{fancy}
\tableofcontents

\chapter*{Agradecimientos}
\addcontentsline{toc}{chapter}{Agradecimientos}
\markboth{Agradecimientos}{Agradecimientos}

Agradezco al CMAT. Tambi\'en por sus consejos:
Dominique Attali,
Juan Pablo Borthagaray,
Lucas Curci,
Fran\c{c}oise Dal'Bo,
Jean-Baptiste Hiriart Urruty,
Mauricio Velasco.
%Quiero agradecer, primero \dots
% Conseils
%%%%%%%%%%
% NON
% Patrick, Yves, Initiation VI
% Dalmasso, Le Marechal,
% Peter Sturm
% JB Durand
% Hedy Attouch
% Mircea Sofonea
% Control, Cherif and co

%\chapter*{Pr\'ologo}
%\addcontentsline{toc}{chapter}{Pr\'ologo}
%\markboth{Pr\'ologo}{Pr\'ologo}

%\medskip

%\begin{flushright}
%{{Franz Chouly}} %, Montevideo, el 6 de Agosto de 2025}}
%\end{flushright}

\chapter{Introducción}

El \textrm{análisis convexo} es una rama fundamental de las matemáticas que estudia funciones y conjuntos convexos, ofreciendo herramientas teóricas y prácticas para resolver problemas complejos en ciencia, ingeniería y economía. Este libro aborda desde los fundamentos matemáticos del análisis convexo hasta sus aplicaciones más avanzadas, incluyendo desigualdades variacionales, optimización convexa, aprendizaje automático e inteligencia artificial, geometría algorítmica y teoría de juegos.

Se presentan las definiciones y propiedades esenciales de conjuntos y funciones convexas, junto con conceptos clave como el subgradiente, la cáscara convexa y la dualidad. Estos fundamentos son la base para entender algoritmos de optimización y técnicas numéricas que se desarrollan en varios capítulos.
%Se incluyen ejemplos ilustrativos, fórmulas matemáticas precisas y fragmentos de código en Python para facilitar la implementación de los métodos descritos.

%El libro está diseñado para estudiantes de grado y posgrado. A lo largo del texto, se enfatiza la relación entre el análisis convexo y otras áreas de las matemáticas, como el análisis funcional, la teoría de la medida y la geometría computacional, mostrando cómo estas disciplinas se complementan para abordar desafíos modernos.

%El enfoque pedagógico del libro facilita la comprensión de conceptos abstractos mediante ejemplos paso a paso, ejercicios propuestos y aplicaciones en contextos reales. 
%Al final de cada capítulo, se proporcionan referencias bibliográficas clave para profundizar en los temas tratados, así como enlaces a recursos adicionales en línea, como repositorios de código y conjuntos de datos.

% Dire sur quoi se base le cours
% References cles
Para el desarrollo de este trabajo, nos apoyaremos fundamentalmente en las obras de G. Allaire~\cite{allaire2007},
de S. Boyd and L. Vandenberghe~\cite{BoydVandenberghe2004},
de J.-B. Hiriart-Urruty and C. Lemar\'{e}chal~\cite{hiriart2001}
y D. G. Luenberger~\cite{luenberger1969}. 
%complementadas con las contribuciones de [bidule] y [chose].
% BREZIS
% Rockafellar
%\cite{Rockafellar1970}
% ETC
A partir de ahora, utilizaremos la notación $V$ para designar un espacio de Hilbert (espacio vectorial normado completo). 
Notaremos $(\cdot,\cdot)_V$ el producto interno de $V$ y $\| \cdot \|_V$ la norma asociada.
Cuando $V$ sea de dimensión finita, y cuando se precisa, lo denotaremos por $V_n$, donde $n$ indica la dimensión de $V$.

\chapter{Descenso de gradiente}
El descenso de gradiente es un algoritmo iterativo de optimización de primer orden utilizado para minimizar funciones. Su idea central es moverse en la dirección opuesta al gradiente de la función en el punto actual, lo que garantiza una disminución local del valor de la función. Este método es ampliamente aplicado en aprendizaje automático y optimización numérica debido a su simplicidad y eficacia. La elección del tamaño de paso, o tasa de aprendizaje, es crucial para la convergencia del algoritmo. Sin embargo, su rendimiento puede verse afectado por la elección de parámetros y la geometría de la función. Seguimos~\cite{allaire2007}.

\section{Aplicaciones}

Antes de describir el algoritmo y hacer su análisis, presentemos ya algunas aplicaciones donde es relevante.

\subsection{Aprendizaje autom\'atico} 
%Función Fuertemente Convexa: Regresión Lineal con Regularización L2 (Ridge)}
%\subsection{Contexto}
En aprendizaje automático, la regresión Ridge busca minimizar el error cuadrático entre las predicciones de un modelo lineal y los datos reales, añadiendo un término de regularización %L2 
para evitar el sobreajuste.

%\subsection{Funcional}
Dada una matriz de datos $ \mathbf{X} \in \mathbb{R}^{n \times d} $, un vector de etiquetas $ \mathbf{y} \in \mathbb{R}^n $, y un vector de pesos $ \mathbf{w} \in \mathbb{R}^d $, la funcional a minimizar es:
$$
\mathcal{J}(\mathbf{w}) := \frac{1}{2n} \| \mathbf{X}\mathbf{w} - \mathbf{y}\|_V^2 + \varepsilon \|\mathbf{w}\|_V^2
$$
donde $ \varepsilon > 0 $ es el parámetro de regularización, y
\[
\|  \mathbf{y}\|_V^2 := y_1^2 + \cdots + y_n^2, \quad
\|\mathbf{w}\|_V^2 := w_1^2 + \cdots + w_d^2. 
\]
Aquí tenemos $V = V_d = \mathbb{R}^d$, con la norma Euclidiana. %[introducir antes].
% VOIR Ridge 
% Explication et reference biblio chic
% Pourquoi ne resoud on pas directement le systeme normal ?
% AI et recherche
% Commencer a incorporer exos
% BIBLIO PLUS
% Ekeland Temam
% Bertsekas
% YVES
% Rockafellar Wets

\subsection{Optimización de forma} %: Área de un Triángulo bajo Restricción de Perímetro}
%\subsection{Contexto}

Ahora presentamos un problema simplificado de optimización de forma.
Se buscan las coordenadas de los vértices de un triángulo que maximice su área, sujetas a una restricción de perímetro fijo. 
%
%\subsection{Funcional}
Para un triángulo con vértices 
\[ \mathbf{p}_1 = (x_1, y_1), \quad \mathbf{p}_2 = (x_2, y_2), \quad \mathbf{p}_3 = (x_3, y_3),\] y un perímetro $ \mathcal{P} $ fijo:
\begin{itemize}
    \item \textbf{Área} (a maximizar):
    \[
    \mathcal{A}(\mathbf{p}_1, \mathbf{p}_2, \mathbf{p}_3) := \frac{1}{2} |  (\mathbf{p}_2 - \mathbf{p}_1) \times (\mathbf{p}_3 - \mathbf{p}_1) |, 
    \]
    donde $\times$ es el producto vectorial en dos dimensiones.
    \item \textbf{Restricción}:
    \[
    \mathcal{R}(
    \mathbf{p}_1 ,
    \mathbf{p}_2 ,
    \mathbf{p}_3 ) = 0
    \]
    con 
    \[
    \mathcal{R}(
    \mathbf{p}_1 ,
    \mathbf{p}_2 ,
    \mathbf{p}_3 ) = 
    \| \mathbf{p}_2 - \mathbf{p}_1 \|_{\mathrm{R}^2}
    + 
    \|
    \mathbf{p}_3 - \mathbf{p}_2 
    \|_{\mathrm{R}^2}
    + \| \mathbf{p}_1 - \mathbf{p}_3 
    \|_{\mathrm{R}^2}
    - \mathcal{P}. \]
\end{itemize}
%
%\subsection{Penalización}
Para transformar el problema restringido en uno no restringido, se usa, por ejemplo:
$$
\mathcal{J}(\mathbf{p}_1, \mathbf{p}_2, \mathbf{p}_3) = -\mathcal{A}(\mathbf{p}_1, \mathbf{p}_2, \mathbf{p}_3) + 
\frac{1}{\varepsilon} \left (
    \mathcal{R}(
    \mathbf{p}_1 ,
    \mathbf{p}_2 ,
    \mathbf{p}_3 ) \right )^2
% \left( \mathbf{p}_2 - \mathbf{p}_1 + \mathbf{p}_3 - \mathbf{p}_2 + \mathbf{p}_1 - \mathbf{p}_3 - P \right)^2
$$
a minimizar en $V_6 = \mathbb{R}^6$, y
donde $ \varepsilon > 0 $ es un parámetro de penalización.

\section{Gradiente a paso constante}

Sea $V = V_n = \mathbb{R}^n$ y
\[
\mathcal{J} : V \rightarrow \mathbb{R}
\]
una funci\'on $\mathscr{C}^1(V;\mathbb{R})$. Notamos $\nabla \mathcal{J}$ su gradiente:
\[
\nabla \mathcal{J} :
\mathbb{R}^n \rightarrow \mathbb{R}^n.
\]
El objetivo es encontrar un mínimo %global 
de esta función. % utilizando el método de descenso de gradiente.

\subsection{Definición}
El método de {descenso de gradiente a paso constante} se define como:
\begin{equation}
{x}_{k+1} = {x}_k - \alpha \nabla \mathcal{J}({x}_k)    
\end{equation}
con $k\geq0$, a partir de $x_0 \in V$ y
donde:
\begin{itemize}
    \item ${x}_k$ es la posición actual en la iteración $k$,
    \item $\alpha$ es el {tamaño del paso} (constante),
    \item $\nabla \mathcal{J}({x}_k)$ es el gradiente de la función $\mathcal{J}$ evaluado en ${x}_k$.
\end{itemize}

\subsection{Python}

\begin{lstlisting}[style=pythonstyle, caption={Gradiente a paso constante.}, label={lst:g1}]
import numpy as np

def J(x):
    return x*2 + 2*x + 1

def dJ(x):
    return 2*x + 2

def gradient_descent(alpha, x0, max_iter=1000):
    x = x0
    for _ in range(max_iter):
        grad = dJ(x)
        x = x - alpha * grad
    return x

# Ejemplo de uso

alpha = 0.1
x0 = 5.0
x_min = gradient_descent(alpha, x0)
print(f"Mínimo encontrado en x = {x_min:.4f}")
\end{lstlisting}

\subsection{Análisis}

Supongamos, para simplificar, que $\mathcal{J}$ cumple las siguientes condiciones:
\begin{enumerate}
    \item 
%1. **Existencia de mínimo global:**
   Existe un único mínimo estricto global de $\mathcal{J}$.

\item %2. **Regularidad:**
   $\mathcal{J}$ es de clase $\mathscr{C}^2(V;\mathbb{R})$ (continuamente diferenciable en $V$).

\item
%3. **Convexidad fuerte:**
   Existe una constante $c_{\mathcal{J}} > 0$ tal que, para todo $x, y \in V$:
   \begin{equation}
       \label{strongconvexity1}
   ( \nabla \mathcal{J}(x) - \nabla \mathcal{J}(y), \, x - y )_V \geq c_{\mathcal{J}} \|x - y\|_V^2.
    \end{equation}
 \end{enumerate}
De hecho, veremos luego que 2 y 3 implican 1.
Entonces, tenemos lo siguiente:

\begin{theorem}
Sea \( \mathcal{J}: V \to \mathbb{R} \) una función de clase %\( \mathscr{C}^2(\mathbb{R}^d;\mathbb{R}) \) 
que cumple con las condiciones 1, 2 y 3 anteriores.
Sea \( x^\star \) el único mínimo estricto global de \( \mathbb{R} \). %Sea \( B \) un entorno de \( x^\star \). 
Entonces, si %\( x_0 \in B \) y 
\( \alpha > 0 \) es suficientemente pequeño, el método de gradiente descendente:
\[
x_{k+1} = x_k - \alpha \nabla \mathcal{J}(x_k),
\]
converge hacia \( x^\star \).
\end{theorem}

\begin{proof}
Sea $x_0 \in V$ y $B_0 \subset V$ una bola cerrada y acotada que contiene $x^\star$ y $x_0$, centrada en $x^\star$. Supongamos $x_k \in B_0$.
Ahora, consideremos la iteración del gradiente descendente:
\[
x_{k+1} = x_k - \alpha \nabla \mathcal{J}(x_k).
\]
Queremos analizar \( \|x_{k+1} - x^\star\|^2 \). Desarrollamos:
\[
\|x_{k+1} - x^\star\|_V^2 = \|x_k - \alpha \nabla \mathcal{J}(x_k) - x^\star\|_V^2 %= \|x_k - x^\star\|^2 - 2\alpha (x_k - x^\star)^T \nabla J(x_k) + \alpha^2 \|\nabla J(x_k)\|^2.
\]
y luego
\[
\|x_{k+1} - x^\star\|_V^2  %\|x_k - \alpha \nabla J(x_k) - x^\star\|^2 
= \|x_k - x^\star\|_V^2 - 2\alpha (x_k - x^\star ,  \nabla \mathcal{J}(x_k) )_V + \alpha^2 \|\nabla \mathcal{J}(x_k)\|_V^2.
\]
Veamos primero:
\[
(x_k - x^\star ,  \nabla \mathcal{J}(x_k) )_V
=
(x_k - x^\star ,  \nabla \mathcal{J}(x_k) - \nabla \mathcal{J}(x^\star) )_V
\geq
c_{\mathcal{J}}
\|x_{k+1} - x^\star\|_V^2.
\]
Hemos utilizado $\nabla \mathcal{J}(x^\star) = 0$ y \eqref{strongconvexity1}.
Esto implica
\[
- 2 \alpha
(x_k - x^\star ,  \nabla \mathcal{J}(x_k) )_V
\leq 
-2 \alpha
c_{\mathcal{J}}
\|x_{k+1} - x^\star\|_V^2.
\]
%
%\medskip
%
Además, como 
$\mathcal{J}$ es de clase $\mathscr{C}^2(B;\mathbb{R}^d)$, en $B_0$, acotada y cerrada, entonces
\( \nabla \mathcal{J} \) es Lipschitz en \( B_0 \). %(porque \( J \) es \( C^2 \))
Existe \( L > 0 \) tal que:
\[
\|\nabla J(x_k)\|_V \leq L \|x_k - x^\star\|_V.
\]
Por lo tanto:
\[
\alpha^2 \|\nabla J(x_k)\|^2_V \leq \alpha^2 L^2 \|x_k - x^\star\|^2_V.
\]
%
%\medskip
%
Sustituyendo todo en la desigualdad:
\[
\|x_{k+1} - x^\star\|_V^2 \leq \|x_k - x^\star\|_V^2 - 2\alpha c_{\mathcal{J}} \|x_k - x^\star\|_V^2 +  \alpha^2 L^2 \|x_k - x^\star\|_V^2.
\]
Simplificando:
\[
\|x_{k+1} - x^\star\|^2_V \leq \left(1 - 2\alpha 
c_{\mathcal{J}} + \alpha^2 L^2\right) \|x_k - x^\star\|^2_V.
\]
%
Elegimos \( \alpha \) tal que:
\[
1 - 2\alpha c_{\mathcal{J}}  + \alpha^2 L^2 < 1.
\]
Esto se cumple si:
\[
-2\alpha c_{\mathcal{J}} +  \alpha^2 L^2 < 0 \quad \Rightarrow \quad 2 \alpha  c_{\mathcal{J}} > \alpha^2 L^2.
\]
La desigualdad se satisface si:
\[
\alpha < \frac{2 c_{\mathcal{J}}}{L^2}.
\]
Por lo tanto, para \( \alpha \) suficientemente pequeño, existe \( \rho \in (0,1) \) tal que:
\[
\|x_{k+1} - x^\star\|^2_V \leq \rho \|x_k - x^\star\|^2_V.
\]
Verificamos que $x_{k+1}$ queda en $B_0$.
Esto implica que la secuencia \( \{x_k\} \) converge linealmente hacia \( x^\star \).
\end{proof}

\section{Otros métodos de minimización}

Mencionamos brevemente otros métodos de minimización que se basan únicamente en el cálculo del gradiente, como el \emph{método del gradiente con paso variable} y el \emph{método del gradiente conjugado}. Para m\'as detalles, referirse a \cite{allaire2007,ggk}.

\subsection{Gradiente a paso variable}

En el método de \emph{descenso de gradiente a paso variable}, el tamaño del paso $\alpha_k$ puede cambiar en cada iteración. Una estrategia común es usar una \emph{búsqueda de línea} (line search) para elegir $\alpha_k$ que minimice $\mathcal{J}({x}_k - \alpha_k \nabla \mathcal{J}({x}_k))$.
%
La actualización se escribe como:
%
$$
{x}_{k+1} = {x}_k - \alpha_k \nabla \mathcal{J}({x}_k).
$$

\subsection{Método del gradiente conjugado}

En el método del \emph{gradiente conjugado}, se optimiza la minimización de funciones cuadráticas de la forma
$
\mathcal{J}({x}) = \frac{1}{2} {x}^T A {x} - \mathbf{b}^T {x} + c,
$
donde $A$ es una matriz simétrica y definida positiva. Este método acelera la convergencia del descenso de gradiente al utilizar direcciones conjugadas con respecto a $A$.
%
La actualización se escribe como:
$$
{x}_{k+1} = {x}_k + \alpha_k {d}_k,
$$
donde ${d}_k$ es la dirección conjugada y $\alpha_k$ se elige mediante una búsqueda de línea o fórmula analítica~\cite[Chapter 11]{ggk}.

\subsection{Método BiCG}

El método \emph{BiCG} (Biconjugate Gradient) está adaptado para resolver sistemas lineales no simétricos de la forma $A{x} = {b}$~\cite[S. 11.7.8]{ggk}, donde $A$ es una matriz cuadrada no necesariamente simétrica. Es una extensión del gradiente conjugado para matrices no simétricas.
%
La actualización se basa en residuos biconjugados y se escribe como:
$$
{x}_{k+1} = {x}_k + \alpha_k {p}_k,
$$
donde ${p}_k$ es la dirección de búsqueda y $\alpha_k$ se calcula para minimizar el error en un subespacio de Krylov.

\section{Para profundizar}

Para un análisis más detallado sobre los métodos de descenso de gradiente y sus variantes, se puede consultar 
\cite[Capítulo~11]{ggk}
y
\cite[Capítulos~8--9]{allaire2007}.

%\addcontentsline{toc}{chapter}{Bibliograf\'ia}
%\nocite{*}
%{\footnotesize
%\bibliographystyle{siam}
%\bibliography{biblio}
%}

%\addcontentsline{toc}{chapter}{\'Indice alfab\'etico}
%{\scriptsize \printindex}

%\end{document} 


\chapter{Funciones convexas}
Las funciones convexas juegan un papel fundamental en optimización, ya que garantizan que cualquier mínimo local es también un mínimo global. Una función $\mathcal{J}: V \to \mathbb{R}$ es convexa si el segmento de línea que une cualquier par de puntos en su gráfico siempre queda por encima del gráfico. Esta propiedad simplifica significativamente el análisis y la resolución de problemas de optimización. Además, las funciones convexas tienen derivadas monótonas, lo que facilita la aplicación de métodos numéricos. Su estudio es esencial para entender la eficiencia de algoritmos como el descenso de gradiente. %Seguimos~\cite{allaire2007}.

%\section{Introducción}

\section{Contexto}

Consideramos el problema de minimización de una función objetivo $\mathcal{J}$ sobre un conjunto $K \subseteq V$, donde $V$ es un espacio de Hilbert. El objetivo es encontrar el \textbf{ínfimo} de $J$ en $K$ y, en caso de existir, un punto donde este ínfimo se alcance.
%
El problema se formula como:
\begin{equation}
    \label{pb}
    \inf_{K} \mathcal{J}.
\end{equation}

\subsection{Definiciones} %Fundamentales}

\begin{definition}[Ínfimo]
Dado un conjunto $K \subseteq V$ y una función $\mathcal{J}: K \to \mathbb{R}$, el \emph{ínfimo} de $J$ sobre $K$ se define como:
\[
\inf_{v \in K} 
\mathcal{J}(v) 
= \inf \{ \mathcal{J}(v) 
\mid v \in K \}.
\]
Es el mayor de los cotas inferiores de $\mathcal{J}$ en $K$. No necesariamente existe un punto $v \in K$ tal que $J(v) = \inf_{v \in K} \mathcal{J}(v)$.
\end{definition}

\begin{definition}[Punto de mínimo local]
Un punto $u^\star \in K$ es un \emph{mínimo local} de $\mathcal{J}$ en $K$ si existe $\delta > 0$ tal que:
\[
\mathcal{J}(u^\star) \leq \mathcal{J}(v) \quad \forall v \in K \text{ con } \|v - u^\star\| < \delta.
\]
\end{definition}

\begin{definition}[Punto de mínimo global]
Un punto $u^\star \in K$ es un \emph{mínimo global} de $\mathcal{J}$ en $K$ si:
\[
\mathcal{J}(u^\star) \leq \mathcal{J}(v) \quad \forall v \in K.
\]
\end{definition}

\begin{definition}[Sucesión minimizante]
Una sucesión $\{u_n\}_{n \in \mathbb{N}} \subseteq K$ se dice \textbf{minimizante} para $\mathcal{J}$ si:
\[
\lim_{n \to \infty} \mathcal{J}(u_n) = \inf_{u \in K} \mathcal{J}(u).
\]
\end{definition}

\subsection{Algunos ejemplos}

\begin{ex}
Consideremos $V = \mathbb{R}$, $K = \mathbb{R}$, y $\mathcal{J}(v) = v^2$. 

\medskip

Entonces:
\[
\inf_{v \in \mathbb{R}} \mathcal{J}(v) = 0,
\]
y el mínimo global se alcanza en $u^\star = 0$. La sucesión $u_n = {1}/{n}$ es minimizante, ya que:
\[
\mathcal{J}(u_n) = \left(\frac{1}{n}\right)^2 \to 0 = \inf_{v \in \mathbb{R}} J(v).
\]
\end{ex}

\begin{ex}
Sea $K = [0, +\infty)$ y $\mathcal{J}(v) = e^{-v}$.

\medskip

Entonces:
\[
\inf_{v \in K} \mathcal{J}(v) = 0,
\]
pero no existe $u^\star \in K$ tal que $J(u^\star) = 0$. Sin embargo, la sucesión $u_n = n$ es minimizante, pues:
\[
\mathcal{J}(u_n) = e^{-n} \to 0.
\]
\end{ex}

\subsection{Existencia %del ínfimo y 
de sucesiones minimizantes}

La proposición siguiente se establece directamente a partir de la definición del ínfimo.

\begin{proposition}
Para cualquier función $\mathcal{J}: K \to \mathbb{R}$, %acotada inferiormente, existe el ínfimo $\inf_{u \in K} J(u)$ y 
siempre existe una {sucesión minimizante} $\{u_n\}_{n \in \mathbb{N}} \subseteq K$ tal que:
\[
\lim_{n \to \infty} \mathcal{J}(u_n) = \inf_{v \in K} \mathcal{J}(v).
\]
\end{proposition}

\section{Existencia de mínimos globales}

Retomamos este resultado general y simple de \cite{allaire2007}.

\begin{theorem}
Sea $V$ un espacio vectorial normado de \emph{dimensión finita}, $K \subseteq V$ un conjunto \emph{cerrado y no vacío}, y $\mathcal{J}: K \to \mathbb{R}$ una función \emph{continua} sobre $K$. Supongamos además que:
\begin{equation}
\label{eq:coercividad}
\text{Si } \{u_n\} \subseteq K \text{ es tal que } \|u_n\| \to +\infty, \text{ entonces } \mathcal{J}(u_n) \to +\infty.
\end{equation}
Entonces, la función $\mathcal{J}:$ admite al menos un \emph{punto de mínimo global} en $K$.
\end{theorem}

\begin{proof} Es un baile en cinco pasos.

\medskip

\noindent \emph{Paso 1.} \emph{Existencia del ínfimo.}
   Como $\mathcal{J}$ está acotada inferiormente (por continuidad y por la condición  %$\|u_n\|_V \to +\infty \implies \mathcal{J}:(u_n) \to +\infty$
   \eqref{eq:coercividad}), existe:
   \[
   m = \inf_{v \in K} \mathcal{J}:(v) \in \mathbb{R}.
   \]

\medskip

\noindent \emph{Paso 2.}
\emph{Existencia de una sucesión minimizante.}
   Por %definición de ínfimo, para 
   la Proposici\'on anterior, para
   todo $n \in \mathbb{N}$, existe $u_n \in K$ tal que:
   %$
   %J(u_n) \leq m + \frac{1}{n}.
   %$
   %Así, $\{u_n\}$ es una sucesión minimizante, ya que:
   \[
   \lim_{n \to \infty} \mathcal{J}(u_n) = m.
   \]

\medskip

\noindent \emph{Paso 3.}
\emph{Acotación de la sucesión minimizante.}
   Supongamos que $\{u_n\}$ no estuviera acotada en $V$. Entonces, existiría una subsucesión $\{u_{n_k}\}$ tal que $\|u_{n_k}\|_V \to +\infty$. Por hipótesis~\eqref{eq:coercividad}, esto implicaría que:
   \[
   \mathcal{J}(u_{n_k}) \to +\infty,
   \]
   lo cual contradice que $\mathcal{J}(u_{n_k}) \to m$. Por lo tanto, $\{u_n\}$ está acotada en $V$.

\medskip

\noindent \emph{Paso 4.}
\emph{Compacidad en dimensión finita.}
   En un espacio de dimensión finita, toda sucesión acotada tiene una subsucesión convergente. Sea $\{u_{n_k}\}$ una subsucesión de $\{u_n\}$ tal que $u_{n_k} \to u^\star \in V$. Como $K$ es cerrado, $u^\star \in K$.

\medskip

\noindent \emph{Paso 5.}
\emph{Continuidad de $\mathcal{J}$.}
   Por la continuidad de $\mathcal{J}$, se tiene:
   $
   \mathcal{J}(u^\star) = \lim_{k \to \infty} \mathcal{J}(u_{n_k}) = m.
   $
   Por lo tanto, $u^\mathcal{J}$ es un punto de mínimo global de $J$ en $K$.
\end{proof}

El teorema anterior no se extiende en dimensión infinita~\cite{allaire2007}. Tampoco asegura unicidad. Esto motiva la introducción de los siguientes conceptos.

%\section{Relación con Funciones Convexas}

%\begin{remark}
%Si además $J$ es una \textbf{función convexa} y $K$ es un conjunto convexo, entonces:
%\begin{itemize}
%    \item Todo punto de mínimo local es también un \textbf{mínimo global}.
%    \item Si $J$ es estrictamente convexa, el mínimo global es \textbf{único}.
%\end{itemize}
%\end{remark}

%---
%\section{Conclusión}
%Bajo las hipótesis del teorema, la existencia de mínimos globales está garantizada. Además, las sucesiones minimizantes siempre existen y, en el caso de espacios de dimensión finita con funciones continuas y coercivas, convergen a un mínimo global.

\section{Conjuntos y funciones convexas} %(dimensión infinita, sobre un espacio de Hilbert $V$)}

%\subsection{Definiciones}
%Sea $V$ un espacio de Hilbert sobre $\mathbb{R}$. 

\begin{definition}
%\begin{itemize}
%    \item 
Un conjunto $K \subset V$ es \emph{convexo} si para todo $v, w \in K$ y todo $\lambda \in [0, 1]$, se cumple:
\begin{equation}
    \label{def:convexo}
    \lambda v + (1 - \lambda) w \in K.
\end{equation}
%    \item 
\end{definition}

%\medskip

\begin{definition}
Una función $\mathcal{J}: V \to \mathbb{R}$ 
%\cup {+\infty}$ 
es \emph{convexa} si para todo $v, w \in V$ y todo $\lambda \in [0, 1]$, se cumple:
\begin{equation}
    \label{def:convexa}
\mathcal{J}(\lambda v + (1 - \lambda) w) \leq \lambda \mathcal{J}(v) + (1 - \lambda) \mathcal{J}(w).
\end{equation}
Es convexa estricta, si, cuando
$\lambda \in (0, 1)$, se cumple:
\begin{equation}
    \label{def:convexaestricta}
\mathcal{J}(\lambda v + (1 - \lambda) w) < \lambda \mathcal{J}(v) + (1 - \lambda) \mathcal{J}(w).
\end{equation}
\end{definition}

\begin{definition}
Una función $\mathcal{J}: V \to \mathbb{R}$ 
%\cup {+\infty}$ 
es \emph{convexa} fuerte, si existe $c_\mathcal{J} > 0$ tal que para todo $v,w \in V$, se cumple:
\begin{equation}
    \label{def:convexafuerte}
\mathcal{J}
\left( \frac12 v + \frac12 w \right) \leq 
 \frac12 \left (\mathcal{J}(v)+\mathcal{J}(w)\right) - \frac{c_\mathcal{J}}{8} \| v-w \|_V^2.
\end{equation}
\end{definition}

\begin{ex}
%\subsection{Ejemplo}
Sea $V = L^2([0,1])$ el espacio de Hilbert de funciones cuadrado-integrables en $[0,1]$.
La función \[
\mathcal{J}(v) := \int_0^1 |v(x)|^2  \mathrm{d}x\] es convexa en $V$.
\end{ex}

\subsection{Obstáculo}
Recordamos aquí la formulación del problema de obstáculo, 
descrito con mayor detalle en, por ejemplo, %\cite{ciarlet2002,glowinski-84,gustafsson2024adaptive}.
\cite{gustafsson2024adaptive}.
%---
%### Configuración general
%
Sea \(\Omega \subset \mathbb{R}^{n}\), con \(n \geq 1\), un dominio abierto, acotado, poliedro, simplemente conexo y con frontera de Lipschitz. Sea \(\kappa\) un parámetro físico (por ejemplo, la rigidez de una membrana).

%---
%### 
%Formas bilineal y lineal
Sea $H^1_0(\Omega)$ el espacio de Sobolev usual~\cite{allaire2007}.
Definimos la forma bilineal:
\[
a : H^1_0 (\Omega) \times H^1_0 (\Omega) \to \mathbb{R}
\]
dada por:
\[
a(u,v) := \int_{\Omega} \kappa \, \nabla u(x) \cdot \nabla v(x) \, \mathrm{d}x, \quad \forall \, u,v \in H^1_0 (\Omega).
\]
%
Asimismo, definimos la forma lineal:
\[
L(v) := \int_{\Omega} f \, v \, dx, \quad \forall \, v \in H^1_0 (\Omega),
\]
donde \(f \in L^2(\Omega)\) es un término fuente.

%### Función de obstáculo

Consideramos una función de obstáculo dada \(\Psi : \Omega \to \mathbb{R}\), que cumple \(\Psi \in H^2(\Omega)\) y \(\Psi \geq 0\) en \(\partial \Omega\). %, como en \cite{ciarlet2002}.
%*Observación: La regularidad del término fuente y de la función de obstáculo puede relajarse con mayor tecnicismo.*
%### Problema de minimización

El problema de minimización asociado consiste en encontrar \(u \in K\) tal que:
$
\min_{v \in K} \mathcal{J}(v),
$
donde \(\mathcal{J}\) es la funcional cuadrática definida como:
\[
\mathcal{J}(v) := \frac{1}{2} a(v,v) - L(v),
\]
y el conjunto convexo \(K\) se define como:
\[
K := \{ v \in H^1_0(\Omega) \mid v \geq \Psi \text{ c.p.p. en } \Omega \}.
\]

Ahora sea la funcional (general) cuadrática:
$
\mathcal{J}(v) = \frac{1}{2} a(v,v) - L(v),
$
donde \(a\) es una forma bilineal elíptica (es decir, existe \(c_J > 0\) tal que \(a(v,v) \geq c_J \|v\|_V^2\) para todo \(v \in V(=H^1_0(\Omega))\)).
%
%---
%### 
%**Demostración de la fuerte convexidad**

Para demostrar que \(\mathcal{J}\) es fuertemente convexa, verificamos la siguiente desigualdad para todo \(u, v \in H^1_0(\Omega)\):

\[
\mathcal{J}\left(\frac{u+v}{2}\right) + \frac{c_J}{8} \| u - v \|^2 \leq \frac{1}{2} \mathcal{J}(u) + \frac{1}{2} \mathcal{J}(v).
\]
%
%---
%%#### **Cálculo paso a paso**
%
%1. **
Desarrollo de \(\mathcal{J}\left(\frac{u+v}{2}\right)\)**:
   \[
   \mathcal{J}\left(\frac{u+v}{2}\right) = \frac{1}{2} a\left(\frac{u+v}{2}, \frac{u+v}{2}\right) - L\left(\frac{u+v}{2}\right).
   \]
   Usando la bilinealidad y la simetría de \(a\):
   \[
   a\left(\frac{u+v}{2}, \frac{u+v}{2}\right) = \frac{1}{4} a(u+v, u+v) = \frac{1}{4} \left( a(u,u) + 2a(u,v) + a(v,v) \right).
   \]
   Por lo tanto:
   \[
   \mathcal{J}\left(\frac{u+v}{2}\right) = \frac{1}{8} \left( a(u,u) + 2a(u,v) + a(v,v) \right) - \frac{1}{2} L(u+v).
   \]
Combinación convexa de \(\mathcal{J}(u)\) y \(\mathcal{J}(v)\):
   \[
   \frac{1}{2} \mathcal{J}(u) + \frac{1}{2} \mathcal{J}(v) = \frac{1}{2} \left( \frac{1}{2} a(u,u) - L(u) \right) + \frac{1}{2} \left( \frac{1}{2} a(v,v) - L(v) \right) %= \frac{1}{4} \left( a(u,u) + a(v,v) \right) - \frac{1}{2} L(u+v).
   \]
  y luego
     \[
   \frac{1}{2} \mathcal{J}(u) + \frac{1}{2} \mathcal{J}(v) %= \frac{1}{2} \left( \frac{1}{2} a(u,u) - L(u) \right) + \frac{1}{2} \left( \frac{1}{2} a(v,v) - L(v) \right) 
   = \frac{1}{4} \left( a(u,u) + a(v,v) \right) - \frac{1}{2} L(u+v).
   \]
%
%3. **
Diferencia entre la combinación convexa y \(\mathcal{J}\left(\frac{u+v}{2}\right)\):
   \[
   \frac{1}{2} \mathcal{J}(u) + \frac{1}{2} \mathcal{J}(v) - \mathcal{J}\left(\frac{u+v}{2}\right) = \frac{1}{8} \left( a(u,u) - 2a(u,v) + a(v,v) \right) %= \frac{1}{8} a(u-v, u-v).
   \]
 y
    \[
   \frac{1}{2} \mathcal{J}(u) + \frac{1}{2} \mathcal{J}(v) - \mathcal{J}\left(\frac{u+v}{2}\right) %= \frac{1}{8} \left( a(u,u) - 2a(u,v) + a(v,v) \right) 
   = \frac{1}{8} a(u-v, u-v).
   \]
   Por la elipticidad de \(a\), sabemos que:
   \[
   a(u-v, u-v) \geq c_J \|u-v\|^2_V.
   \]
   Por lo tanto:
   \[
   \frac{1}{2} \mathcal{J}(u) + \frac{1}{2} \mathcal{J}(v) - \mathcal{J}\left(\frac{u+v}{2}\right) \geq \frac{c_J}{8} \|u-v\|^2.
   \]
   Reordenando:
   \[
   \mathcal{J}\left(\frac{u+v}{2}\right) \leq \frac{1}{2} \mathcal{J}(u) + \frac{1}{2} \mathcal{J}(v) - \frac{c_J}{8} \|u-v\|^2.
   \]
%#### **Conclusión**
La funcional \(\mathcal{J}\) es fuertemente convexa con constante \({c_J}\), ya que cumple:
\[
\mathcal{J}\left(\frac{u+v}{2}\right) + \frac{c_J}{8} \|u-v\|^2 \leq \frac{1}{2} \mathcal{J}(u) + \frac{1}{2} \mathcal{J}(v).
\]
%
%
%\section{Existencia de mínimos}
%
%\begin{theorem}
%Si $V$ es un espacio de Hilbert y $f: V \to \mathbb{R} \cup {+\infty}$ es una función:
%\begin{itemize}
%    \item \textbf{convexa},
%    \item \textbf{semi-continua inferiormente} (para la topología fuerte de $V$),
%    \item \textbf{coerciva} (es decir, $f(v) \to +\infty$ cuando $|v|_V \to +\infty$),
%\end{itemize}
%entonces $f$ alcanza su mínimo en $V$.
%\end{theorem}

%\subsection{Ejemplo}
%Sea $V = H^1_0(\Omega)$ (espacio de Sobolev) y $f(u) = \frac{1}{2} \int_\Omega |\nabla u|^2 , dx - \int_\Omega fu , dx$, donde $f \in L^2(\Omega)$.
%Bajo ciertas condiciones, $f$ es convexa, continua y coerciva, por lo que tiene un mínimo en $V$.

%\title{Teoremas sobre Funciones Convexas y Mínimos}
% Entornos personalizados
%\newtheorem{theorem}{Teorema}[section]
%\newtheorem{cor}{Corolario}[section]
%\newtheorem{lemma}{Lema}[section]
%\newtheorem{proposition}{Proposición}[section]
%\newtheorem{definition}{Definición}[section]
%\newtheorem{remark}{Observación}[section]
%\newtheorem{ex}{Ejemplo}[section]

\section{Funciones convexas y mínimos}

Los siguientes resultados son adaptados de \cite{allaire2007}.

\subsection{Mínimos globales}

\begin{theorem}
Si $\mathcal{J}$ es convexa, entonces todo mínimo local es un mínimo global. Además, el conjunto de mínimos de $\mathcal{J}$ es convexo.
\end{theorem}

\begin{proof}
Sea $u^\star$ un mínimo local de $\mathcal{J}$, es decir, existe $\delta > 0$ tal que:
\[
\mathcal{J}(u^\star) \leq \mathcal{J}(u) \quad \forall u \in K \text{ con } \|u - u^*\|_V < \delta.
\]
Queremos demostrar que $u^\star$ es un mínimo global, es decir, que:
\[
\mathcal{J}(u^\star) \leq \mathcal{J}(u) \quad \forall u \in K.
\]
%
Sea $u \in K$ arbitrario. Consideremos la función $f: [0,1] \to \mathbb{R}$ definida por:
\[
f(t) = \mathcal{J}((1-t)u^\star + t u).
\]
Como $\mathcal{J}$ es convexa, se tiene:
\[
f(t) = \mathcal{J}((1-t)u^\star + t u) \leq (1-t)\mathcal{J}(u^\star) + t \mathcal{J}(u).
\]
%
Supongamos, por contradicción, que existe $u \in K$ tal que $\mathcal{J}(u) < \mathcal{J}(u^\star)$. Entonces, para todo $t \in (0,1)$:
\[
f(t) \leq (1-t)\mathcal{J}(u^\star) + t \mathcal{J}(u) < (1-t)\mathcal{J}(u^\star) + t \mathcal{J}(u^\star) = \mathcal{J}(u^\star).
\]
Pero como $u^*$ es un mínimo local, existe $\delta > 0$ tal que para todo $t$ suficientemente pequeño (por ejemplo, $t < \delta / \|u - u^*\|_V$), se cumple:
\[
f(t) \geq \mathcal{J}(u^\star).
\]
Esto es una contradicción, ya que $f(t) < \mathcal{J}(u^\star)$ y $f(t) \geq \mathcal{J}(u^*)$ no pueden ser ciertos simultáneamente. Por lo tanto, $\mathcal{J}(u^\star) \leq \mathcal{J}(u)$ para todo $u \in K$, y $u^\star$ es un mínimo global.

\medskip 

Para demostrar que el conjunto de mínimos es convexo, sean $u_1, u_2$ dos mínimos globales de $\mathcal{J}$ y $t \in [0,1]$. Entonces:
\[
\mathcal{J}(t u_1 + (1-t) u_2) \leq t \mathcal{J}(u_1) + (1-t) \mathcal{J}(u_2) %= t \inf_{u \in K} \mathcal{J}(u) + (1-t) \inf_{u \in K} \mathcal{J}(u) = \inf_{u \in K} \mathcal{J}(u).
\]
y
\[
%\mathcal{J}(t u_1 + (1-t) u_2) \leq 
t \mathcal{J}(u_1) + (1-t) \mathcal{J}(u_2) = t \inf_{v \in K} \mathcal{J}(v) + (1-t) \inf_{v \in K} \mathcal{J}(v) = \inf_{v \in K} \mathcal{J}(v).
\]
Como $\inf_{v \in K} \mathcal{J}(v)$ es el valor mínimo, se tiene:
\[
\mathcal{J}(t u_1 + (1-t) u_2) = \inf_{v \in K} \mathcal{J}(v),
\]
lo que implica que $t u_1 + (1-t) u_2$ es también un mínimo global. Por lo tanto, el conjunto de mínimos es convexo.
\end{proof}

\begin{theorem}
Si $\mathcal{J}$ es estrictamente convexa, entonces existe a lo sumo un mínimo global.
\end{theorem}

\begin{proof}
Supongamos, por contradicción, que existen dos mínimos globales distintos $u_1$ y $u_2$ tales que:
\[
\mathcal{J}(u_1) = \mathcal{J}(u_2) = \inf_{v \in K} \mathcal{J}(v).
\]
Como $\mathcal{J}$ es estrictamente convexa, para todo $t \in (0,1)$ se cumple:
\[
\mathcal{J}(t u_1 + (1-t) u_2) < t \mathcal{J}(u_1) + (1-t) \mathcal{J}(u_2) 
%= t \inf_{u \in K} \mathcal{J}(u) + (1-t) \inf_{u \in K} \mathcal{J}(u) = \inf_{u \in K} \mathcal{J}(u).
\]
y
\[
%\mathcal{J}(t u_1 + (1-t) u_2) < 
t \mathcal{J}(u_1) + (1-t) \mathcal{J}(u_2) = t \inf_{v \in K} \mathcal{J}(v) + (1-t) \inf_{v \in K} \mathcal{J}(v) = \inf_{v \in K} \mathcal{J}(v).
\]
Esto implica que:
\[
\mathcal{J}(t u_1 + (1-t) u_2) < \inf_{v \in K} \mathcal{J}(v),
\]
lo cual contradice la definición de ínfimo. Por lo tanto, no pueden existir dos mínimos globales distintos, y el mínimo global, si existe, es único.
\end{proof}

\subsection{Caso fuertemente convexo}

Tenemos lo siguiente.

\begin{theorem}
    Sea \(K\) un conjunto no vacío, cerrado y convexo en un espacio de Hilbert \(V\) y sea \(\mathcal{J}\) una función continua y fuertemente convexa en \(V\). Entonces \(\mathcal{J}\) admite un único minimizador global \(u\) en \(K\). Además, para todo \(v \in K\), se cumple:
    \[
    \| v - u \|_V^2 \leq \frac{4}{c_J} \left( \mathcal{J}(v) - \mathcal{J}(u) \right),
    \]
    y toda subsucesión minimizante de \(\mathcal{J}\) converge hacia \(u\).
\end{theorem}

\begin{proof}
    %Sea \(\{u_n\}\) una sucesión minimizante para \(\mathcal{J}\) en \(K\) (existe por definición del ínfimo y dado que \(K\) es no vacío). Como \(\mathcal{J}\) es continua y \(K\) es cerrado, el ínfimo \(I = \inf_{v \in K} \mathcal{J}(v)\) es finito.
    
    \emph{Paso 1: Existencia de una sucesión minimizante y acotación.}
    Dado que \(\mathcal{J}\) está acotada inferiormente, por la Proposición~\ref{convexafhahnbanach}, existe \(I = \inf_{v \in K} \mathcal{J}(v) > -\infty\). 
    Sea \(\{u_n\}\) una sucesión minimizante para \(\mathcal{J}\) en \(K\) (existe por definición del ínfimo y dado que \(K\) es no vacío).
    Por definición de sucesión minimizante, se tiene \[\lim_{n \to \infty} \mathcal{J}(u_n) = I.\]

    \emph{Paso 2: Uso de la fuerte convexidad.}
    Aplicamos la propiedad de fuerte convexidad de \(\mathcal{J}\) con constante \(c_J > 0\). Para \(n, m \in \mathbb{N}\) arbitrarios, se cumple:
    \[
    \mathcal{J}\left(\frac{u_n + u_m}{2}\right) \leq \frac{1}{2} \mathcal{J}(u_n) + \frac{1}{2} \mathcal{J}(u_m) - \frac{c_J}{8} \| u_n - u_m \|^2.
    \]
    Restando \(I\) a ambos lados y usando que \(\mathcal{J}\left(\frac{u_n + u_m}{2}\right) \geq I\) (pues \(I\) es el ínfimo), obtenemos:
    \[
    \frac{c_J}{8} \| u_n - u_m \|^2 \leq \frac{1}{2} \mathcal{J}(u_n) + \frac{1}{2} \mathcal{J}(u_m) - I.
    \]
    Como \(\mathcal{J}(u_n) \to I\) y \(\mathcal{J}(u_m) \to I\), el lado derecho tiende a cero cuando \(n, m \to \infty\). Por lo tanto, \(\{u_n\}\) es una sucesión de Cauchy en \(K\).

\medskip

    \emph{Paso 3: Convergencia y minimizador global.}
    Dado que \(K\) es cerrado y \(V\) es completo, \(\{u_n\}\) converge a un elemento \(u \in K\). Por la continuidad de \(\mathcal{J}\):
    \[
    \mathcal{J}(u) = \lim_{n \to \infty} \mathcal{J}(u_n) = I.
    \]
    Así, \(u\) es un minimizador global de \(\mathcal{J}\) en \(K\).

\medskip

    \emph{Paso 4: Unicidad del minimizador.}
    Supongamos que existe otro minimizador \(\tilde{u} \in K\). Aplicando la fuerte convexidad:
    \[
    I = \mathcal{J}(u) \leq \mathcal{J}\left(\frac{u + \tilde{u}}{2}\right) \leq \frac{1}{2} \mathcal{J}(u) + \frac{1}{2} \mathcal{J}(\tilde{u}) - \frac{c_J}{8} \| u - \tilde{u} \|^2.
    \]
    Como \(\mathcal{J}(u) = \mathcal{J}(\tilde{u}) = I\), esto implica:
    \[
    0 \leq -\frac{c_J}{8} \| u - \tilde{u} \|^2.
    \]
    Dado que \(c_J > 0\), se deduce que \(\| u - \tilde{u} \|^2 = 0\), es decir, \(u = \tilde{u}\). Por lo tanto, el minimizador es único.

\medskip

    \emph{Paso 5: Desigualdad para el minimizador.}
    Sea \(v \in K\). Como \(K\) es convexo, \({(u + v)}/{2} \in K\). Aplicando la fuerte convexidad:
    \[
    \mathcal{J}\left(\frac{u + v}{2}\right) \leq \frac{1}{2} \mathcal{J}(u) + \frac{1}{2} \mathcal{J}(v) - \frac{c_J}{8} \| u - v \|^2.
    \]
    Dado que \(u\) es minimizador, \(\mathcal{J}\left(\frac{u + v}{2}\right) \geq \mathcal{J}(u)\). Sustituyendo:
    \[
    \mathcal{J}(u) \leq \frac{1}{2} \mathcal{J}(u) + \frac{1}{2} \mathcal{J}(v) - \frac{c_J}{8} \| u - v \|^2.
    \]
    Simplificando:
    \[
    \frac{c_J}{8} \| u - v \|^2 \leq \frac{1}{2} \left( \mathcal{J}(v) - \mathcal{J}(u) \right).
    \]
    Multiplicando por 4:
    \[
    \| u - v \|^2 \leq \frac{4}{c_J} \left( \mathcal{J}(v) - \mathcal{J}(u) \right),
    \]
    lo que concluye la demostración.
\end{proof}

Aquí el enunciado del Teorema de Hahn-Banach (versión geométrica: separación de un convexo cerrado por un hiperplano) en un espacio de Hilbert, junto con una prueba simplificada que utiliza el teorema de proyección en un convexo cerrado. 

\begin{theorem}[Teorema de Hahn-Banach (Separación)]
    Sea \(V\) un espacio de Hilbert y \(K \subset V\) un conjunto no vacío, cerrado y convexo. Si \[v_0 \in V \setminus K,\] entonces existe un hiperplano cerrado que separa estrictamente \(v_0\) de \(K\). Es decir, existen \(z \in V\) y \(\beta \in \mathbb{R}\) tales que:
    \[
    (z,v_0)_V
    < \beta < 
    (z,w)_V, \quad \forall w \in K.
    \]
\end{theorem}

\begin{proof} Es una valse en tres pasos.

\medskip

    \emph{Paso 1: Proyección sobre el convexo cerrado.}
    Dado que \(K\) es cerrado y convexo, por el teorema de proyección en un convexo cerrado, existe un único \(w_0 \in K\) tal que, para todo \(w \in K\):
    \[ ( v_0 - w_0 , w - w_0 )_V \leq 0.
%    \[
%    \|v_0 - u\| = \min_{v \in K} \|v_0 - u\|.
    \]
    Este \(w_0\) es la proyección ortogonal de \(v_0\) sobre \(K\).

\medskip

    \emph{Paso 2: Construcción del hiperplano separador}.
    Definimos 
    \[
    m := \frac{w_0 + v_0}{2}, \quad
    %\]
    %y 
    %\[
    z := w_0 - m = \frac{w_0 - v_0}{2}.
    \]
    Como $v_0 \notin K$, $z$ es no nulo.
    Ahora,
    consideramos el hiperplano:
    \[
    \mathcal{H} := \{ v \in V \mid ( z , v )_V = \beta \},
    \]
    donde \(\beta = ( z , m )_V\).

\medskip

    \emph{Paso 3: Separación estricta.}
    %
    Primero tenemos
    \[
    (z, m - v_0)_V 
    = \frac14( w_0 - v_0 , w_0 - v_0 )
    > 0
    \]
    y luego
    \[
    (z , v_0 ) < ( z , m )_V = \beta.
    \]
    Tambi\'en ******
    \[
    (z,w)_V 
    =
    \frac12(w_0 - v_0, w )_V
    =
    \frac12
    \left (
    (w_0 - v_0, w - w_0 )_V
    + ( w_0 - v_0 , w_0 )_V \right ).
    \]
    y, como para todo \(w \in K\):
    \[ ( v_0 - w_0 , w - w_0 )_V \leq 0,\]
    entonces
    \[
    (z,w)_V \geq
    \frac12
    ( w_0 - v_0 , w_0 )_V.
    \]
    Acotamos
    \[
    (z,w)_V - \beta 
    \geq
        \frac12
    ( w_0 - v_0 , w_0 )_V
    -(z,m)_V
    =
        \frac12
    ( w_0 - v_0 , w_0 )_V
    - \frac14 ( w_0 - v_0 , w_0 + v_0 )_V.
    \]
    y
    \[
    (z,w)_V - \beta 
    \geq
        \frac12
    ( w_0 - v_0 , w_0 )_V
    -(z,m)_V
    =
        \frac12
    ( w_0 - v_0 , w_0 )_V
    - \frac14 ( w_0 - v_0 , w_0 + v_0 )_V.
    \]
    \[
    (z,w)_V - \beta 
    \geq
    =
    ( w_0 - v_0 ,\frac12 w_0 
    - \frac14w_0 - \frac14v_0)_V > 0.
    \]
\end{proof}

---

---

\begin{proposition}[Proposición (9.2.5, Allaire)]
Si $\mathcal{J}$ es convexa y continua sobre un conjunto $K$ cerrado, convexo y no vacío, entonces existen $\delta \in \mathbb{R}$ y $L \in V'$ tales que para todo $v \in K$:
\[
\mathcal{J}(v) \geq L(v) + \delta.
\]
\end{proposition}

\begin{proof}
Como $\mathcal{J}$ es convexa y continua, su epigrafo:
\[
\text{epi}(\mathcal{J}) = \{(v, \mu) \in V \times \mathbb{R} \mid \mathcal{J}(v) \leq \mu\}
\]
es un conjunto convexo y cerrado en $V \times \mathbb{R}$. Además, como $K$ es no vacío, existe al menos un punto $(v_0, \mathcal{J}(v_0)) \in \text{epi}(\mathcal{J})$.
%
Consideremos el conjunto:
\[
C = \left \{(v, \mu) \in V \times \mathbb{R} \mid \mu < \inf_{v \in K} \mathcal{J}(v) \right \}.
\]
Este conjunto es convexo y abierto, y no interseca con $\text{epi}(\mathcal{J})$. Por el \emph{teorema de separación de Hahn-Banach}, existen un funcional lineal $F \in (V \times \mathbb{R})'$ y $\alpha \in \mathbb{R}$ tales que:
\[
F(v, \mu) \leq \alpha \quad \forall (v, \mu) \in \text{epi}(\mathcal{J}),
\]
\[
F(v, \mu) > \alpha \quad \forall (v, \mu) \in C.
\]
%
Podemos descomponer $F$ como $F(v, \mu) = L(v) + \beta \mu$, donde $L \in V'$ y $\beta \in \mathbb{R}$. De las desigualdades anteriores, se deduce que $\beta > 0$ (de lo contrario, no se cumpliría la segunda desigualdad para $\mu \to -\infty$). Sin pérdida de generalidad, podemos normalizar $F$ para que $\beta = 1$. Entonces:
\[
L(v) + \mu \leq \alpha \quad \forall (v, \mu) \in \text{epi}(\mathcal{J}).
\]
En particular, para todo $v \in V$ y $\mu = \mathcal{J}(v)$, se tiene:
\[
L(v) + \mathcal{J}(v) \leq \alpha.
\]
Reordenando, obtenemos:
\[
\mathcal{J}(v) \geq -L(v) + \alpha.
\]
Definiendo $\delta = -\alpha$ y reemplazando $L$ por $-L$, se obtiene:
\[
\mathcal{J}(v) \geq L(v) + \delta.
\]
\end{proof}

\begin{proposition} \label{convexafhahnbanach}
Además, si $\mathcal{J}$ es fuertemente convexa, entonces existen $\gamma > 0$ y $\delta \in \mathbb{R}$ tales que para todo $v \in V$:
\[
\mathcal{J}(v) \geq \gamma \|v\|^2 + \delta.
\]
\end{proposition}

\begin{proof}
Como $\mathcal{J}$ es fuertemente convexa, existe $\gamma > 0$ tal que la función $\mathcal{J}(v) - \frac{\gamma}{2} \|v\|^2$ es convexa. Aplicando la proposición anterior a esta función, existen $\delta' \in \mathbb{R}$ y $L \in V'$ tales que:
\[
\mathcal{J}(v) - \frac{\gamma}{2} \|v\|^2 \geq L(v) + \delta'.
\]
Reordenando, obtenemos:
\[
\mathcal{J}(v) \geq \frac{\gamma}{2} \|v\|^2 + L(v) + \delta'.
\]
Como $\frac{\gamma}{2} \|v\|^2 + L(v) \geq \frac{\gamma}{4} \|v\|^2 - \frac{1}{\gamma} \|L\|^2$ (por la desigualdad de Young), podemos ajustar $\gamma$ y $\delta$ para obtener:
\[
\mathcal{J}(v) \geq \gamma \|v\|^2 + \delta,
\]
donde $\delta = \delta' - \frac{1}{\gamma} \|L\|^2$.
\end{proof}

%\addcontentsline{toc}{chapter}{Bibliograf\'ia}
%\nocite{*}
%{\footnotesize
%\bibliographystyle{siam}
%\bibliography{biblio}
%}

%\addcontentsline{toc}{chapter}{\'Indice alfab\'etico}
%{\scriptsize \printindex}

%\end{document} 

\newpage

\subsection{Caso estrictamente convexo}

%### **Ejemplos en dimensión finita**

1. **Función cuadrática degenerada**
   Consideremos la función \(f: \mathbb{R}^2 \to \mathbb{R}\) definida por:
   \[
   f(x, y) = x^2 + y^4.
   \]
   - **Estrictamente convexa**: Su matriz hessiana es
     \[
     H_f(x, y) = \begin{pmatrix} 2 & 0 \\ 0 & 12y^2 \end{pmatrix},
     \]
     que es positiva definida para todo \((x, y) \neq (0, 0)\) (pues \(12y^2 > 0\) si \(y \neq 0\) y \(2 > 0\)).
     Además, \(f\) cumple la desigualdad de convexidad estricta:
     \[
     f(\lambda x + (1-\lambda) y) < \lambda f(x) + (1-\lambda) f(y), \quad \forall \lambda \in (0, 1), \, x \neq y.
     \]
   - **No fuertemente convexa**: No existe \(\alpha > 0\) tal que \(f(x, y) - \frac{\alpha}{2} \| (x, y) \|^2\) sea convexa. En particular, la segunda derivada en la dirección \(y\) tiende a \(0\) cuando \(y \to 0\), lo que impide la existencia de una constante global \(\alpha > 0\) que acote inferiormente el hessiano.

---

2. **Función exponencial**
   Consideremos la función \(g: \mathbb{R} \to \mathbb{R}\) definida por:
   \[
   g(x) = e^x.
   \]
   - **Estrictamente convexa**: Su segunda derivada es \(g''(x) = e^x > 0\) para todo \(x \in \mathbb{R}\), lo que garantiza convexidad estricta.
   - **No fuertemente convexa**: No existe \(\alpha > 0\) tal que \(g''(x) \geq \alpha\) para todo \(x \in \mathbb{R}\), ya que \(g''(x) = e^x \to 0\) cuando \(x \to -\infty\).

---

---
---

%### **Ejemplos en dimensión infinita (vinculados a EDP elípticas)**

1. **Funcional asociado al \(p\)-Laplaciano con \(1 < p < 2\)**
   Sea \(\Omega \subset \mathbb{R}^n\) un dominio acotado y consideremos el funcional \(J: W_0^{1,p}(\Omega) \to \mathbb{R}\) definido por:
   \[
   J(u) = \frac{1}{p} \int_\Omega |\nabla u|^p \, dx.
   \]
   - **Estrictamente convexo**: El operador \(p\)-Laplaciano \(-\Delta_p u = -\text{div}(|\nabla u|^{p-2} \nabla u)\) es estrictamente monótono para \(1 < p < \infty\), lo que implica que \(J\) es estrictamente convexo.
   - **No fuertemente convexo**: Para \(1 < p < 2\), el funcional \(J\) **no es fuertemente convexo** en \(W_0^{1,p}(\Omega)\). Esto se debe a que la derivada de \(|\nabla u|^p\) no está acotada inferiormente por una constante positiva uniforme en todo el espacio. En particular, la "fuerte convexidad" falla porque la segunda derivada del término \(|\nabla u|^p\) no es uniformemente elíptica.

---

2. **Funcional de energía asociado a una EDP elíptica no lineal**
   Consideremos el funcional \(E: H_0^1(\Omega) \to \mathbb{R}\) definido por:
   \[
   E(u) = \int_\Omega \left( \sqrt{1 + |\nabla u|^2} - 1 \right) dx - \int_\Omega f u \, dx,
   \]
   donde \(f \in L^2(\Omega)\).
   - **Estrictamente convexo**: El integrando \(\sqrt{1 + |\nabla u|^2}\) es estrictamente convexo en \(\nabla u\), ya que su hessiana (con respecto a \(\nabla u\)) es positiva definida. Esto garantiza que \(E\) es estrictamente convexo.
   - **No fuertemente convexo**: La segunda derivada del término \(\sqrt{1 + |\nabla u|^2}\) es:
     \[
     \frac{\partial^2}{\partial (\nabla u)^2} \sqrt{1 + |\nabla u|^2} = \frac{I}{1 + |\nabla u|^2} - \frac{\nabla u \otimes \nabla u}{(1 + |\nabla u|^2)^{3/2}},
     \]
     que no está acotada inferiormente por una constante positiva uniforme (por ejemplo, cuando \(|\nabla u| \to \infty\), la norma de esta matriz tiende a \(0\)). Por lo tanto, \(E\) no es fuertemente convexo.

---

\newpage

\section{El caso convexo}

**1. Ejemplos de funciones convexas pero no estrictamente convexas (relativos a EDP elípticas)**

**Ejemplo 1: Funcional asociado al operador \(p\)-Laplaciano con \(p = 1\)**
Consideremos el funcional \(J: BV(\Omega) \to \mathbb{R}\) definido por:
\[
J(u) = \int_\Omega |\nabla u| \, dx,
\]
donde \(BV(\Omega)\) es el espacio de funciones de variación acotada.
- **Convexo**: El funcional \(J\) es convexo porque \(|\nabla u|\) es una norma (y las normas son convexas).
- **No estrictamente convexo**: No cumple la desigualdad estricta de convexidad. Por ejemplo, si \(u\) y \(v\) son funciones con \(\nabla u = -\nabla v\), entonces:
  \[
  J\left(\frac{u + v}{2}\right) = \int_\Omega \left|\nabla \left(\frac{u + v}{2}\right)\right| \, dx = \int_\Omega \left|\frac{\nabla u + \nabla v}{2}\right| \, dx = \int_\Omega \left|\frac{\nabla u - \nabla u}{2}\right| \, dx = 0,
  \]
  mientras que:
  \[
  \frac{1}{2} J(u) + \frac{1}{2} J(v) = \int_\Omega |\nabla u| \, dx > 0.
  \]
  Por lo tanto, \(J\left(\frac{u + v}{2}\right) = \frac{1}{2} J(u) + \frac{1}{2} J(v)\) (no hay desigualdad estricta).

---
**Ejemplo 2: Funcional de energía con término lineal**
Consideremos el funcional \(E: H_0^1(\Omega) \to \mathbb{R}\) definido por:
\[
E(u) = \int_\Omega |\Delta u| \, dx - \int_\Omega f u \, dx,
\]
donde \(f \in L^2(\Omega)\).
- **Convexo**: El término \(\int_\Omega |\Delta u| \, dx\) es convexo (porque es una norma), y el término lineal \(-\int_\Omega f u \, dx\) es convexo (y cóncavo a la vez).
- **No estrictamente convexo**: El término \(\int_\Omega |\Delta u| \, dx\) no es estrictamente convexo, ya que su derivada segunda es nula en casi todo punto (no hay fuerte convexidad).

---

---

---

---

**2. Definición: Convergencia débil en un espacio de Hilbert**

Sea \(V\) un espacio de Hilbert. Una secuencia \(\{u_n\} \subset V\) **converge débilmente** a \(u \in V\) si:
\[
\lim_{n \to \infty} \langle u_n, v \rangle_V = \langle u, v \rangle_V, \quad \forall v \in V.
\]
Se denota \(u_n \rightharpoonup u\).

---

---

---

**3. Proposición: Convergencia débil en Hilbert separable**

\begin{proposition}
    Sea \(V\) un espacio de Hilbert **separable**. Entonces, una secuencia \(\{u_n\} \subset V\) converge débilmente a \(u \in V\) si y solo si:
    \[
    \lim_{n \to \infty} \langle u_n, e_i \rangle_V = \langle u, e_i \rangle_V, \quad \forall i \in \mathbb{N},
    \]
    donde \(\{e_i\}_{i \in \mathbb{N}}\) es una **base hilbertiana** (base ortonormal completa) de \(V\).
\end{proposition}

---
**Demostración:**

- **Implicación directa (\(\Rightarrow\)):**
  Si \(u_n \rightharpoonup u\), entonces por definición de convergencia débil:
  \[
  \lim_{n \to \infty} \langle u_n, v \rangle_V = \langle u, v \rangle_V, \quad \forall v \in V.
  \]
  En particular, esto vale para \(v = e_i\) (pues \(e_i \in V\)).

- **Implicación inversa (\(\Leftarrow\)):**
  Supongamos que \(\lim_{n \to \infty} \langle u_n, e_i \rangle_V = \langle u, e_i \rangle_V\) para todo \(i \in \mathbb{N}\). Sea \(v \in V\). Como \(\{e_i\}\) es una base hilbertiana, podemos escribir:
  \[
  v = \sum_{i=1}^\infty \langle v, e_i \rangle_V e_i,
  \]
  con convergencia en la norma de \(V\). Entonces:
  \[
  \langle u_n, v \rangle_V = \left\langle u_n, \sum_{i=1}^\infty \langle v, e_i \rangle_V e_i \right\rangle_V = \sum_{i=1}^\infty \langle v, e_i \rangle_V \langle u_n, e_i \rangle_V.
  \]
  Por hipótesis, \(\langle u_n, e_i \rangle_V \to \langle u, e_i \rangle_V\) para cada \(i\). Además, como \(\{u_n\}\) es acotada (pues converge débilmente), por el **Teorema de Lebesgue dominado**, podemos intercambiar el límite y la suma:
  \[
  \lim_{n \to \infty} \langle u_n, v \rangle_V = \sum_{i=1}^\infty \langle v, e_i \rangle_V \langle u, e_i \rangle_V = \langle u, v \rangle_V.
  \]
  Por lo tanto, \(u_n \rightharpoonup u\).

---

---

---

---

    **Ejemplo:**
    En \(\ell^2\), la secuencia \(u_n = (1, 0, 0, \dots)\) (con \(1\) en la posición \(n\) y \(0\) en las demás) satisface:
    - \(u_n \rightharpoonup 0\) (pues \(\langle u_n, v \rangle = v_n \to 0\) para todo \(v \in \ell^2\)).
    - \(\|u_n\| = 1\) para todo \(n\), por lo que **no converge fuertemente** a \(0\).


\noindent Bressan: ejemplo con sen al cuadrado y un medio.
---

 **5. Lema: De toda secuencia acotada se puede extraer una subsucesión que converge débilmente**

\begin{lemma}
    Sea \(V\) un espacio de Hilbert y \(\{u_n\} \subset V\) una secuencia **acotada** (es decir, \(\|u_n\| \leq M\) para todo \(n\) y algún \(M > 0\)). Entonces existe una subsucesión \(\{u_{n_k}\}\) que converge débilmente a algún \(u \in V\).
\end{lemma}

---
**Demostración:**

- **Base hilbertiana:** Como \(V\) es separable, existe una base hilbertiana \(\{e_i\}_{i \in \mathbb{N}}\).
- **Acotación de las componentes:** Para cada \(i \in \mathbb{N}\), la secuencia \(\{\langle u_n, e_i \rangle\}\) está acotada en \(\mathbb{R}\) (pues \(|\langle u_n, e_i \rangle| \leq \|u_n\| \|e_i\| \leq M\)).
- **Subsucesión diagonal:** Por el **Teorema de Bolzano-Weierstrass**, para cada \(i\), existe una subsucesión \(\{u_n^{(i)}\}\) tal que \(\{\langle u_n^{(i)}, e_i \rangle\}\) converge. Usando un argumento de **subsucesión diagonal**, podemos construir una subsucesión \(\{u_{n_k}\}\) tal que:
  \[
  \lim_{k \to \infty} \langle u_{n_k}, e_i \rangle \text{ existe para todo } i \in \mathbb{N}.
  \]
- **Convergencia débil:** Por la **Proposición 3**, esto implica que \(\{u_{n_k}\}\) converge débilmente a algún \(u \in V\).

---

---

---

 **6. Lema: Un convexo cerrado \(K\) es exactamente la intersección de los semiplanos cerrados que lo contienen**

\begin{lemma}
    Sea \(V\) un espacio de Hilbert y \(K \subset V\) un conjunto **cerrado y convexo**. Entonces:
    \[
    K = \bigcap_{\substack{H \supset K \\ H \text{ semiplano cerrado}}} H.
    \]
\end{lemma}

---
**Demostración (por contradicción usando Hahn-Banach):**

- **Inclusión trivial:** Claramente, \(K \subset \bigcap_{\substack{H \supset K \\ H \text{ semiplano cerrado}}} H\), ya que todo semiplano cerrado que contiene a \(K\) incluye a \(K\).
- **Inclusión inversa:** Supongamos que existe \(x \in \bigcap_{\substack{H \supset K \\ H \text{ semiplano cerrado}}} H\) pero \(x \notin K\). Como \(K\) es cerrado y convexo, por el **Teorema de Hahn-Banach** (versión geométrica), existe un hiperplano cerrado que separa estrictamente \(x\) de \(K\). Es decir, existen \(a \in V\) y \(\beta \in \mathbb{R}\) tales que:
  \[
  \langle y, a \rangle \leq \beta < \langle x, a \rangle, \quad \forall y \in K.
  \]
  Esto implica que el semiplano cerrado \(H = \{ y \in V \mid \langle y, a \rangle \leq \beta \}\) contiene a \(K\) pero **no contiene a \(x\)**, lo que contradice la hipótesis de que \(x\) está en la intersección de todos los semiplanos cerrados que contienen a \(K\). Por lo tanto, \(x \in K\).

---

---

---

 **7. Lema: Un convexo cerrado \(K\) es cerrado para la convergencia débil**

\begin{lemma}
    Sea \(V\) un espacio de Hilbert y \(K \subset V\) un conjunto **cerrado y convexo**. Entonces \(K\) es **cerrado para la convergencia débil**, es decir:
    Si \(\{u_n\} \subset K\) y \(u_n \rightharpoonup u\), entonces \(u \in K\).
\end{lemma}

---
**Demostración:**

- Por el **Lema 6**, \(K\) es la intersección de todos los semiplanos cerrados que lo contienen. Sea \(H = \{ y \in V \mid \langle y, a \rangle \leq \beta \}\) uno de estos semiplanos. Como \(u_n \in K \subset H\) para todo \(n\), se tiene:
  \[
  \langle u_n, a \rangle \leq \beta.
  \]
  Tomando límite cuando \(n \to \infty\) (por la convergencia débil):
  \[
  \langle u, a \rangle \leq \beta.
  \]
  Por lo tanto, \(u \in H\) para todo semiplano cerrado \(H\) que contiene a \(K\). Por el **Lema 6**, esto implica que \(u \in K\).

---

---

---
---

 **8. Lema: Si \(J\) es convexa y semi-continua inferiormente (s.c.i.) sobre \(K\), entonces \(J\) es s.c.i. para la convergencia débil**

\begin{lemma}
    Sea \(V\) un espacio de Hilbert, \(K \subset V\) un conjunto **cerrado y convexo**, y \(J: K \to \mathbb{R}\) una función **convexa y semi-continua inferiormente (s.c.i.)** para la convergencia fuerte. Entonces \(J\) es también **s.c.i. para la convergencia débil**.
\end{lemma}

---
**Demostración:**

- Sea \(\{u_n\} \subset K\) una secuencia que converge débilmente a \(u \in K\) (por el **Lema 7**, \(u \in K\)).
- Como \(J\) es convexa y s.c.i. para la convergencia fuerte, por el **Teorema de Mazur**, existe una combinación convexa de los \(u_n\) que converge **fuertemente** a \(u\). Es decir, para cada \(k \in \mathbb{N}\), existen \(\lambda_i^{(k)} \geq 0\) con \(\sum_{i=1}^{N_k} \lambda_i^{(k)} = 1\) tales que:
  \[
  v_k = \sum_{i=1}^{N_k} \lambda_i^{(k)} u_i \quad \text{y} \quad v_k \to u \text{ fuertemente}.
  \]
- Por la convexidad de \(J\):
  \[
  J(v_k) \leq \sum_{i=1}^{N_k} \lambda_i^{(k)} J(u_i).
  \]
- Tomando límite inferior cuando \(k \to \infty\) y usando que \(J\) es s.c.i. para la convergencia fuerte:
  \[
  J(u) \leq \liminf_{k \to \infty} J(v_k) \leq \liminf_{k \to \infty} \sum_{i=1}^{N_k} \lambda_i^{(k)} J(u_i).
  \]
  Como \(\{u_n\}\) converge débilmente a \(u\), y \(J\) es convexa, se puede mostrar que:
  \[
  J(u) \leq \liminf_{n \to \infty} J(u_n).
  \]
  Por lo tanto, \(J\) es s.c.i. para la convergencia débil.

---

---
---
---

---
 **9. Existencia de un mínimo: Caso convexo, \(J\) continua, coerciva y convexa sobre \(K\)**

\begin{theorem}
    Sea \(V\) un espacio de Hilbert, \(K \subset V\) un conjunto **cerrado, convexo y no vacío**, y \(J: K \to \mathbb{R}\) una función **continua, coerciva** (es decir, \(\lim_{\|u\| \to \infty} J(u) = +\infty\)) y **convexa**. Entonces \(J\) admite un **mínimo global** en \(K\).
\end{theorem}

---
**Prueba (pasos detallados):**

1. **Existencia de una secuencia minimizante:**
   Sea \(I = \inf_{u \in K} J(u)\). Por definición del ínfimo, existe una secuencia \(\{u_n\} \subset K\) tal que:
   \[
   J(u_n) \to I.
   \]

2. **La secuencia minimizante es acotada (por contradicción):**
   Supongamos que \(\{u_n\}\) no es acotada. Entonces existe una subsucesión \(\{u_{n_k}\}\) tal que \(\|u_{n_k}\| \to \infty\). Como \(J\) es coerciva:
   \[
   J(u_{n_k}) \to +\infty,
   \]
   lo que contradice que \(J(u_{n_k}) \to I\). Por lo tanto, \(\{u_n\}\) es acotada.

3. **Existencia de una subsucesión que converge débilmente:**
   Por el **Lema 5**, existe una subsucesión \(\{u_{n_k}\}\) que converge débilmente a algún \(u \in V\). Además, por el **Lema 7**, como \(K\) es cerrado y convexo, \(u \in K\).

4. **\(J(u)\) es el mínimo:**
   Por el **Lema 8**, \(J\) es s.c.i. para la convergencia débil. Por lo tanto:
   \[
   J(u) \leq \liminf_{k \to \infty} J(u_{n_k}) = I.
   \]
   Como \(I\) es el ínfimo, se tiene \(J(u) \geq I\). Por lo tanto, \(J(u) = I\), y \(u\) es un **minimizador global** de \(J\) en \(K\).

---

\newpage

%\section{Funciones diferenciables}

%\subsection{Definición} 
%Sea $f: V \to \mathbb{R}$ una función. Decimos que $f$ es \textbf{diferenciable en el sentido de Fréchet} en $x \in V$ si existe un funcional lineal continuo $f'(x) \in V'$ (el dual de $V$) tal que:
%
%$$
%\lim_{h \to 0} \frac{|f(x + h) - f(x) - \langle f'(x), h \rangle|}{|h|_V} = 0.
%$$
%
%Si $f$ es diferenciable en todo $V$, decimos que $f$ es \textbf{diferenciable en $V$}.
%
%\subsection{Ejemplo}
%Sea $V = H^1_0(\Omega)$ y $f(u) = \frac{1}{2} \int_\Omega |\nabla u|^2 , dx$.
%La derivada de Fréchet de $f$ en $u$ es el funcional lineal:

%$$
%\langle f'(u), v \rangle = \int_\Omega \nabla u \cdot \nabla v , dx, \quad \forall v \in V.
%$$

\section{Criterios de convexidad con diferencial}

Aquí los criterios con derivada de $\mathcal{J}$ y derivada segunda.

\section{Criterios de optimalidad} %bajo restricciones de desigualdad (Euler-Lagrange)}

%\subsection{Condiciones de optimalidad}
Sea $K \subset V$ un conjunto convexo y cerrado, y $\mathcal{J}: V \to \mathbb{R}$ una función convexa y diferenciable.
Un punto $u \in K$ es mínimo de $\mathcal{J}$ en $K$ si y solo si:
%
$$
\langle \mathcal{J}'(u), v - u \rangle \geq 0, \quad \forall v \in K.
$$
%
Esta es la \textbf{desigualdad variacional de Euler-Lagrange}.
%papi te quiero
%\subsection{Multiplicadores de Lagrange}
%Si las restricciones son de la forma $g_i(u) \leq 0$ para $i = 1, \dots, m$, donde $g_i: V \to \mathbb{R}$ son funciones convexas y diferenciables,
%las condiciones de optimalidad de Karush-Kuhn-Tucker (KKT) afirman que existe $u \in V$ y multiplicadores $\lambda_i \geq 0$ tales que:

%\begin{align*}
%\langle f'(u) + \sum_{i=1}^m \lambda_i g_i'(u), v - u \rangle &\geq 0, \quad \forall v \in V, \
%\lambda_i g_i(u) &= 0, \quad \forall i = 1, \dots, m.
%\end{align*}

%\chapter{Condiciones de optimalidad}
%Las condiciones de optimalidad son criterios matemáticos que permiten identificar si un punto es un mínimo local o global de una función. En optimización sin restricciones, un punto $x^*$ es un mínimo local si el gradiente $\nabla f(x^*) = 0$ y la matriz hessiana $\nabla^2 f(x^*)$ es definida positiva. Estas condiciones son fundamentales para desarrollar algoritmos de optimización y analizar su convergencia. Además, en problemas con restricciones, se utilizan condiciones como las de Karush-Kuhn-Tucker (KKT). Su comprensión es clave para resolver problemas prácticos de manera eficiente.

%\textbf{Sugerencia de plan:}
%\begin{itemize}
%    \item Condiciones necesarias y suficientes para óptimos locales y globales.
%    \item Aplicación en problemas con y sin restricciones.
%    \item Ejemplos prácticos.
%\end{itemize}

\chapter{Método de Newton}
El método de Newton es un algoritmo de optimización de segundo orden que utiliza información de la segunda derivada para acelerar la convergencia. A diferencia del descenso de gradiente, este método aproxima la función objetivo localmente por una función cuadrática y encuentra su mínimo en una sola iteración. Aunque requiere el cálculo de la matriz hessiana, su convergencia es típicamente más rápida, especialmente cerca del óptimo. Es ampliamente utilizado en problemas donde la precisión y la velocidad son críticas. Sin embargo, su implementación puede ser costosa en términos computacionales.

%\section{Método de Newton / Newton semi-suave}


\section{Idea del método}

Seguimos los libros de Suli y Allaire.
%
El \emph{método de Newton} utiliza aproximaciones cuadráticas locales para acelerar la convergencia. La iteración se define como:
\[
x_{k+1} = x_k - [\nabla^2 f(x_k)]^{-1} \nabla f(x_k),
\]
donde \( \nabla^2 f(x_k) \) es la matriz Hessiana de \( f \) en \( x_k \). Este método converge cuadráticamente cerca de mínimos locales (Teorema de Kantorovich).
%
%\begin{quotation}
%\textbf{Vínculo con ecuaciones diferenciales:}
El método de Newton puede verse como una aproximación de la dinámica de segundo orden:
\[
\frac{d^2x}{dt^2} + \gamma \frac{dx}{dt} = -\nabla f(x),
\]
donde \( \gamma \) es un parámetro de amortiguamiento. Esta ecuación modela el movimiento de una partícula en un paisaje de energía potencial \( f \).
%\end{quotation}


\section{Mas sobre funciones convexas}

\begin{definition}[Función convexa en dimensión finita]
Una función \[f: \mathbb{R}^n \to (-\infty, +\infty] \] es \textbf{convexa} si su \textbf{dominio extendido} \( \text{dom}(f) = \{x \in \mathbb{R}^n \mid f(x) < +\infty\} \) es un conjunto convexo y, para todo \( x, y \in \text{dom}(f) \) y todo \( \lambda \in [0,1] \), se cumple:
\[
f(\lambda x + (1-\lambda) y) \leq \lambda f(x) + (1-\lambda) f(y).
\]
\end{definition}

%\subsubsection{Explicación del dominio extendido}
El \textbf{dominio extendido} \( \text{dom}(f) \) incluye todos los puntos donde \( f \) no toma el valor \( +\infty \). Este concepto es crucial porque:
\begin{itemize}
    \item Permite manejar funciones definidas solo en subconjuntos específicos de \( \mathbb{R}^n \) (por ejemplo, restricciones en optimización).
    \item Garantiza que la desigualdad de convexidad se aplique \textbf{solo} a pares de puntos donde \( f \) está bien definida.
\end{itemize}

\subsubsection{Ejemplos ilustrativos}

\begin{ex}[Función valor absoluto: \( f(x) = |x| \)] Tiene las c\'aracteristicas siguientes.
\begin{itemize}
    \item \textbf{Dominio extendido:} \( \text{dom}(f) = \mathbb{R} \).
    \item \textbf{Convexidad:} Cumple la desigualdad de convexidad (por desigualdad triangular):
    \[
    |\lambda x + (1-\lambda)y| \leq \lambda |x| + (1-\lambda)|y|.
    \]
    \item \textbf{Visualización:}
\end{itemize}

\begin{center}
\begin{tikzpicture}[scale=0.99, domain=-2:2]
    \draw[-] (-3,0) -- (3,0) node[right] {$x$};
    \draw[-] (0,-0.5) -- (0,2.5) node[above] {$f(x)$};
    \draw[blue, thick] plot (\x, {abs(\x)});
    \node[below] at (2,0) {$f(x) = |x|$};
\end{tikzpicture}
\end{center}
\end{ex}

\begin{ex}[Función parte positiva: \( f(x) = \max(x, 0) \)] Tiene las c\'aracteristicas siguientes.
\begin{itemize}
    \item \textbf{Dominio extendido:} \( \text{dom}(f) = \mathbb{R} \).
    \item \textbf{Convexidad:} %Es lineal por tramos, por lo que es convexa.
    se verifica directamente con la d\'efinici\'on.
    \item \textbf{Visualización:}
\end{itemize}

\begin{center}
\begin{tikzpicture}[scale=0.99, domain=-2:2]
    \draw[-] (-3,0) -- (3,0) node[right] {$x$};
    \draw[-] (0,-0.5) -- (0,2.5) node[above] {$f(x)$};
    \draw[blue, thick] plot (\x, {max(\x,0)});
    \node[below] at (2,0) {$f(x) = \max(x, 0)$};
\end{tikzpicture}
\end{center}
\end{ex}

\begin{ex}[Función indicatriz de un conjunto convexo \( C \)]
% \textbf{Función indicatriz de un conjunto convexo \( C \):}
Se define as\'i:
    \[
    f(x) = \begin{cases}
    0 & \text{si } x \in C, \\
    +\infty & \text{si } x \notin C.
    \end{cases}
    \]
    \begin{itemize}
        \item \textbf{Dominio extendido:} \( \text{dom}(f) = C \), que es convexo por definición.
        \item \textbf{Convexidad:} La desigualdad se cumple trivialmente (ambos lados son \( +\infty \) o \( 0 \)).
    \end{itemize}
\end{ex}

\begin{ex}
[Función norma euclidiana]  

  \textbf{Función norma euclidiana:} Se define por \( f(x) = \|x\|_2 \).
    \begin{itemize}
        \item \textbf{Dominio extendido:} \( \text{dom}(f) = \mathbb{R}^n \).
        \item \textbf{Convexidad:} Desigualdad triangular:
        \[
        \|\lambda x + (1-\lambda)y\|_2 \leq \lambda \|x\|_2 + (1-\lambda)\|y\|_2.
        \]
    \end{itemize}
\end{ex}

%\subsection{Función convexa propia: definición y propiedades}
%\label{subsec:funcion_convexa_propia}

Las funciones convexas propias son un caso especialmente importante en optimización y análisis convexo, ya que combinan la convexidad con condiciones naturales sobre su dominio y valores. A continuación, presentamos su definición formal. % y sus implicaciones clave.

\begin{definition}[Función convexa propia]
Una función \( f: \mathbb{R}^n \to (-\infty, +\infty] \) se denomina \textbf{convexa propia} si cumple las siguientes condiciones:
\begin{enumerate}
    \item \textbf{Convexidad:} \( f \) es convexa en el sentido usual:
    \[
    f(\lambda x + (1-\lambda) y) \leq \lambda f(x) + (1-\lambda) f(y), \quad \forall x, y \in \mathbb{R}^n, \forall \lambda \in [0,1].
    \]
    \item \textbf{Condición de no trivialidad:} \( f(x) < +\infty \) para al menos un \( x \in \mathbb{R}^n \) (es decir, \( \text{dom}(f) \neq \emptyset \)).
    %\item %\textbf{Condición de no infinitud global:} \( f(x) > -\infty \) para todo \( x \in \mathbb{R}^n \) (no existen puntos donde \( f \) tome el valor \( -\infty \)).
\end{enumerate}
\end{definition}

\subsection{Subdiferencial de funciones convexas}

La noción de subdiferencial generaliza el concepto de gradiente a funciones convexas que no son necesariamente diferenciables. Esta herramienta es fundamental en optimización convexa, ya que permite caracterizar óptimos incluso en presencia de no diferenciabilidad.

\begin{definition}[Subgradiente]
Sea \( f: \mathbb{R}^n \to (-\infty, +\infty] \) una función convexa. Un vector \( g \in \mathbb{R}^n \) es un \textbf{subgradiente} de \( f \) en \( x \) si
\[
f(y) \geq f(x) + g^T (y - x) \quad \forall y \in \mathbb{R}^n.
\]
El conjunto de todos los subgradientes de \( f \) en \( x \) se denota \( \partial f(x) \) y se llama \textbf{subdiferencial} de \( f \) en \( x \).
\end{definition}

Intuitivamente, un subgradiente define un hiperplano de soporte al epígrafo de la función en el punto considerado. Cuando la función es diferenciable, el subdiferencial se reduce a un singleton que contiene el gradiente clásico.

A continuación, resumimos algunas propiedades fundamentales del subdiferencial en dimensión finita.

\paragraph{Propiedades clave.}
Sea \( f: \mathbb{R}^n \to (-\infty,+\infty] \) convexa y \( x \in \mathrm{dom}(f) \). Entonces:
\begin{itemize}
    \item El conjunto \( \partial f(x) \) es convexo y cerrado.
    \item Si \( f \) es diferenciable en \( x \), entonces
    \[
    \partial f(x) = \{ \nabla f(x) \}.
    \]
    \item La condición de optimalidad de primer orden se escribe:
    \[
    0 \in \partial f(x) \quad \Longleftrightarrow \quad x \text{ es un minimizador global de } f.
    \]
    \item La aplicación \( x \mapsto \partial f(x) \) es monótona en el sentido de operadores:
    \[
    (g_1 - g_2)^T (x_1 - x_2) \geq 0
    \quad \text{para todo } g_i \in \partial f(x_i).
    \]
\end{itemize}

Para ilustrar esta noción, presentamos ahora algunos ejemplos clásicos de cálculo de subdiferenciales.

\paragraph{Ejemplos.}

\begin{itemize}
    \item \textbf{Valor absoluto.} Sea \( f(x) = |x| \) con \( x \in \mathbb{R} \). Entonces:
    \[
    \partial f(x) =
    \begin{cases}
        \{1\} & \text{si } x > 0, \\
        \{-1\} & \text{si } x < 0, \\
        [-1,1] & \text{si } x = 0.
    \end{cases}
    \]

    \item \textbf{Parte positiva.} Sea \( f(x) = \max(0,x) \). Entonces:
    \[
    \partial f(x) =
    \begin{cases}
        \{1\} & \text{si } x > 0, \\
        \{0\} & \text{si } x < 0, \\
        [0,1] & \text{si } x = 0.
    \end{cases}
    \]

    \item \textbf{Función indicatriz.} Sea \( C \subset \mathbb{R}^n \) un conjunto convexo cerrado y
    \[
    \iota_C(x) =
    \begin{cases}
        0 & \text{si } x \in C, \\
        +\infty & \text{si } x \notin C.
    \end{cases}
    \]
    Entonces, para \( x \in C \), el subdiferencial está dado por el cono normal:
    \[
    \partial \iota_C(x) = N_C(x) = \{ g \in \mathbb{R}^n \mid g^T (y - x) \leq 0, \ \forall y \in C \}.
    \]
\end{itemize}

Estos ejemplos muestran cómo el subdiferencial captura la falta de diferenciabilidad mediante conjuntos de pendientes admisibles, lo cual resulta esencial en el análisis y la resolución de problemas de optimización convexa no suave.

\subsection{Caracterización de mínimos} % mediante el subdiferencial}

Una de las principales ventajas del subdiferencial es que permite formular condiciones de optimalidad de primer orden sin requerir diferenciabilidad. En el contexto de funciones convexas, esta caracterización es particularmente simple y poderosa.
%
El siguiente resultado establece una condición necesaria y suficiente de optimalidad global.

\begin{theorem}[Caracterización de mínimos]
Sea \( f: \mathbb{R}^n \to (-\infty, +\infty] \) una función convexa. Un punto \( x^* \in \mathrm{dom}(f) \) es un mínimo global de \( f \) si y solo si
\[
0 \in \partial f(x^*).
\]
\end{theorem}

Este resultado generaliza la condición clásica \( \nabla f(x^*) = 0 \) del caso diferenciable. En efecto, cuando \( f \) es diferenciable en \( x^* \), se tiene \( \partial f(x^*) = \{\nabla f(x^*)\} \), y la condición anterior se reduce a la conocida anulación del gradiente.

\paragraph{Idea de la demostración.}
La demostración se basa directamente en la definición de subgradiente.
\begin{itemize}
    \item (\emph{Suficiencia}) Si \( 0 \in \partial f(x^*) \), entonces por definición:
    \[
    f(y) \geq f(x^*) \quad \forall y \in \mathbb{R}^n,
    \]
    lo que implica que \( x^* \) es un mínimo global.
    \item (\emph{Necesidad}) Si \( x^* \) es un mínimo global, entonces
    \[
    f(y) \geq f(x^*) \quad \forall y,
    \]
    lo que corresponde exactamente a la desigualdad de subgradiente con \( g = 0 \).
\end{itemize}

Para una demostración más detallada y en un marco más general (espacios vectoriales normados), se puede consultar textos clásicos de optimización convexa como \cite{hiriart2001}. %Rockafellar o Hiriart-Urruty y Lemaréchal.

\paragraph{Interpretación geométrica.}
La condición \( 0 \in \partial f(x^*) \) significa que existe un hiperplano de soporte horizontal al epígrafo de \( f \) en \( x^* \). Intuitivamente, esto indica que no existe ninguna dirección de descenso, lo que caracteriza un mínimo global en el caso convexo.

A continuación, ilustramos esta caracterización con ejemplos clásicos.

\paragraph{Ejemplos.}

\begin{itemize}
    \item \textbf{Valor absoluto.} Sea \( f(x) = |x| \). Entonces:
    \[
    \partial f(0) = [-1,1].
    \]
    Como \( 0 \in [-1,1] \), se concluye que \( x^* = 0 \) es un mínimo global. En cambio, si \( x \neq 0 \), entonces \( \partial f(x) \) no contiene al cero, lo que confirma que no son minimizadores.

    \item \textbf{Parte positiva.} Sea \( f(x) = \max(0,x) \). Se tiene:
    \[
    \partial f(0) = [0,1].
    \]
    Dado que \( 0 \in \partial f(0) \), el punto \( x^* = 0 \) es un mínimo global (de hecho, \( f(x) \geq 0 \) para todo \( x \)).

    \item \textbf{Función indicatriz.} Sea \( C \subset \mathbb{R}^n \) convexo cerrado y \( f = \iota_C \). Entonces:
    \[
    \partial f(x) = N_C(x).
    \]
    La condición de optimalidad se escribe:
    \[
    0 \in N_C(x^*),
    \]
    lo cual equivale a decir que
    \[
    0^T (y - x^*) \leq 0 \quad \forall y \in C,
    \]
    condición que es siempre verdadera. Esto refleja el hecho de que todo punto de \( C \) es minimizador de \( \iota_C \).
\end{itemize}

\paragraph{Comentario.}
Este criterio es fundamental en optimización convexa y sirve como base para numerosos algoritmos (métodos de subgradiente, métodos proximales, etc.). Su simplicidad contrasta con su gran generalidad, ya que sigue siendo válido incluso en ausencia de diferenciabilidad.

\section{Newton semi-suave}

Aquí describimos el algoritmo de Newton semi-suave.

\section{Restricciones}

Para problemas con restricciones, podemos aplicar además las siguiente técnicas.

\subsection{Penalización}

Ver \cite{glowinski80}.

\subsection{Lagrangiano}

Ver \cite{glowinski80} y \cite{allaire2007}.
\subsection{Lagrangiano aumentado}

Ver \cite{glowinski80} y trabajos recientes de Erik Burman.

\chapter{Desigualdades variacionales}

Las desigualdades variacionales son útiles para modelar numerosos problemas de la física o la ingeniería~\cite{duvaut76,glowinski80,kinderlehrer80}. Sin embargo, rara vez se abordan en los cursos de posgrado de análisis numérico. El objetivo es presentar el análisis numérico del problema del obstáculo, cuando se discretiza mediante elementos finitos. El problema del obstáculo es, de hecho, el prototipo ideal de una desigualdad variacional y, quizás, el más sencillo de ellos.

Denotamos por \(H^{s}(\cdot)\), \(s\in \mathbb{R}\), a los espacios de Sobolev. Para un subconjunto abierto \(D\) de \(\mathbb{R}^n\), la norma usual de \(H^{s}(D)\) se denota por \(\|\cdot\|_{s,D}\). El espacio \(H^1_0(D)\) es el subespacio de funciones en \(H^1(D)\) con traza nula en \(\partial D\). La letra \(C\) representa una constante genérica, independiente de los parámetros de discretización.

%\section{El obstáculo}

\section{Ejemplos}

Seguimos~\cite{glowinski80}.
%\section{Desigualdades variacionales}
Las desigualdades variacionales proporcionan un marco general para estudiar problemas con restricciones, como el del obstáculo. Este planteamiento permite modelar fenómenos físicos donde las soluciones deben satisfacer condiciones de desigualdad en lugar de igualdades. La teoría de desigualdades variacionales es fundamental para entender la estructura matemática del problema.

%\section{Formulaciones débiles de problemas variacionales}

%---
\subsection{Problema de Signorini (contacto unilateral sin fricción)}
Sea \( \Omega \subseteq \mathbb{R}^d \) un dominio acotado con frontera \( \Gamma = \Gamma_D \cup \Gamma_C \cup \Gamma_N \) (donde \( \Gamma_D \) es la parte de la frontera con desplazamiento prescrito, \( \Gamma_C \) es la zona de contacto potencial y \( \Gamma_N \) es la parte con carga aplicada).
El problema de Signorini consiste en encontrar \( u \in K \), donde:
\[
K = \{ v \in H^1(\Omega) \colon v = 0 \text{ en } \Gamma_D, \, v \geq 0 \text{ en } \Gamma_C \},
\]
tal que:
\[
a(u, v - u) + j(v) - j(u) \geq \langle f, v - u \rangle \quad \forall v \in K.
\]
Aquí:
\begin{itemize}
    \item \( a(u, v) = \int_\Omega \sigma_{ij}(u) \varepsilon_{ij}(v) \, dx \) es la forma bilineal simétrica asociada al tensor de deformaciones \( \varepsilon \) y tensiones \( \sigma \).
    \item \( j(v) = 0 \) (en este caso, sin fricción).
    \item \( \langle f, v \rangle = \int_\Omega f v \, dx + \int_{\Gamma_N} g v \, ds \) es el funcional lineal de cargas.
    \item \( \sigma_{ij}(u) \) son las componentes del tensor de tensiones, y \( \varepsilon_{ij}(v) \) las del tensor de deformaciones.
\end{itemize}

\subsection{Problema de torsión (Saint-Venant)}
Sea \( \Omega \subseteq \mathbb{R}^2 \) una sección transversal de una barra. El problema de torsión consiste en encontrar \( \phi \in H^1_0(\Omega) \) tal que:
\[
a(\phi, v) = \langle f, v \rangle \quad \forall v \in H^1_0(\Omega),
\]
donde:
\begin{itemize}
    \item \( a(\phi, v) = \int_\Omega \nabla \phi \cdot \nabla v \, dx \) (forma bilineal simétrica).
    \item \( \langle f, v \rangle = \int_\Omega G \theta v \, dx \), con \( G \) el módulo de corte, \( \theta \) el ángulo de torsión por unidad de longitud, y \( f = 2G\theta \) (constante).
    \item El convexo $K$ es relacionado a $|\nabla u| \leq c$.
\end{itemize}

%---
\subsection{Problema de Tresca (fricción)}
Sea \( \Omega \) un dominio con frontera \( \Gamma = \Gamma_D \cup \Gamma_C \cup \Gamma_N \). El problema de Tresca con fricción en \( \Gamma_C \) consiste en encontrar \( u \in K \) tal que:
\[
a(u, v - u) + j(v) - j(u) \geq \langle f, v - u \rangle \quad \forall v \in K,
\]
donde:
\begin{itemize}
    \item \( K = \{ v \in H^1(\Omega) \colon v = 0 \text{ en } \Gamma_D \} \).
    \item \( j(v) = \int_{\Gamma_C} g |v_t| \, ds \), con \( g \) el coeficiente de fricción y \( v_t \) la componente tangencial de \( v \) en \( \Gamma_C \).
    \item \( a(u, v) \) es la forma bilineal simétrica asociada al problema elástico.
\end{itemize}

\subsection{Problema de Bingham (fluido visco-plástico)}
El fluido de Bingham se modela mediante la siguiente desigualdad variacional:
Encontrar \( u \in K \) tal que:
\[
a(u, v - u) + j(v) - j(u) \geq \langle f, v - u \rangle \quad \forall v \in K,
\]
donde:
\begin{itemize}
    \item \( K = \{ v \in H^1_0(\Omega)^d \colon \text{div } v = 0 \} \) (campo de velocidades incompresible).
    \item \( a(u, v) = \int_\Omega 2\mu \varepsilon_{ij}(u) \varepsilon_{ij}(v) \, dx \), con \( \mu \) la viscosidad plástica.
    \item \( j(v) = \int_\Omega \tau_0 \sqrt{2} |\varepsilon(v)| \, dx \), con \( \tau_0 \) el límite de fluencia.
    \item \( \langle f, v \rangle = \int_\Omega f \cdot v \, dx \).
\end{itemize}

\subsection{Formulación con forma bilineal no simétrica}
Consideremos un problema variacional con una forma bilineal no simétrica. Por ejemplo, sea \( \Omega \subseteq \mathbb{R}^d \) y el problema:
Encontrar \( u \in H^1_0(\Omega) \) tal que:
\[
a(u, v) = \langle f, v \rangle \quad \forall v \in H^1_0(\Omega),
\]
donde la forma bilineal \( a(u, v) \) se define como:
\[
a(u, v) = \int_\Omega \nabla u \cdot \nabla v \, dx + \int_\Omega \mathbf{b} \cdot \nabla u \, v \, dx,
\]
con \( \mathbf{b} \in L^\infty(\Omega)^d \) un campo vectorial dado.
Aquí, el término \( \int_\Omega \mathbf{b} \cdot \nabla u \, v \, dx \) introduce la asimetría en la forma bilineal.

\section{Teoría general}

Veamos ahora resultados sobre desigualdades de tipo I y de tipo II.
%papa
%\section{Existencia y unicidad}
El estudio de la existencia y unicidad de soluciones es crucial para garantizar la validez de los modelos matemáticos. En el caso del obstáculo, se demuestra que, bajo condiciones adecuadas, existe una solución única que satisface las restricciones físicas del problema. Este análisis sienta las bases para el desarrollo de métodos numéricos confiables.


%Teoria general. Glowinski.

\subsection{Teorema de Stampacchia}

Aquí detallamos el Teorema de Stampacchia, para tipo I.

\subsection{Teorema de Lions-Stampacchia}

Aquí detallamos del Teorema de Lions-Stampacchia, para tipo II.

\section{Galerkin y elementos finitos}

Aproximaciones de tipo elementos finitos. Convergencia y estimaciones a priori, ver \cite{glowinski80}.

%Ver nota.

%\section{Formulación fuerte}
%Este enfoque clásico del problema del obstáculo presenta la formulación en términos de ecuaciones diferenciales parciales y condiciones de contorno. 
%
%
%
%\section{Aproximación de Galerkin}
La aproximación de Galerkin es un método numérico ampliamente utilizado para resolver problemas descritos por desigualdades variacionales. Este enfoque permite discretizar el problema y transformarlo en un sistema algebraico resoluble. Su simplicidad y eficacia lo convierten en una herramienta estándar en el análisis numérico.

\subsection{Penalización}
El método de penalización reemplaza las restricciones del problema original por términos en la función objetivo. Este enfoque simplifica la implementación numérica y permite aplicar técnicas estándar de resolución. Sin embargo, la elección del parámetro de penalización es clave para la estabilidad y precisión del método.

\subsection{Newton semi-suave}
El método de Newton semi-suave es una técnica avanzada para resolver problemas de desigualdades variacionales, combinando ideas de métodos de Newton y subgradientes. Este enfoque aprovecha la estructura diferenciable a trozos de las funciones involucradas, mejorando la convergencia en problemas no lineales. Su aplicación es especialmente relevante en contextos donde la solución puede presentar comportamientos complejos.

\subsection{Implementación en scikit-fem}
scikit-fem es una biblioteca de Python que facilita la implementación de métodos de elementos finitos para resolver problemas variacionales. Su integración con NumPy y SciPy permite desarrollar soluciones eficientes y modulares. Esta herramienta es ideal para prototipar y validar algoritmos numéricos en entornos académicos y de investigación. Veamos el listado siguiente, escrito por Lucas Curci.

\begin{lstlisting}[caption={Diferenciación automática con scikit-fem.}, label={lst:autodiff-skfem}, style=pythonstyle]
"""Diferenciación automática.
Utiliza diferenciación automática para derivar el sistema tangente
para el método de Newton.
"""

from skfem import *
from skfem.autodiff import NonlinearForm
from skfem.autodiff.helpers import grad, dot
import numpy as np
import jax.numpy as jnp

m = MeshTri().refined(5)

@NonlinearForm(hessian=True)
def J(u, _):
    eps = 1e-2
    return (1/2 * dot(grad(u), grad(u)) - u
            + (1 / eps) * ((u - jnp.minimum(u.value, 0.01))**2))

basis = Basis(m, ElementTriP1())

D = m.boundary_nodes()
x = np.zeros(basis.N)

for itr in range(100):
    x_prev = x.copy()
    x += solve(*condense(*J.assemble(basis, x=x), D=D))
    res = np.linalg.norm(x - x_prev)
    if res < 1e-8:
        break
    if __name__ == "__main__":
        print(res)

if __name__ == "__main__":
    from skfem.visuals.matplotlib import plot3, show
    plot3(m, x)
    show()
\end{lstlisting}


%\begin{figure}[htbp]
%    \centering
%    \includegraphics[width=0.8\linewidth]{download (1).png}
%    \caption{Solución del obstáculo penalizado con scikit-fem}
%    \label{fig:exemple_image}
%\end{figure}

%\begin{verbatim}
%    """
%from skfem import *
%from skfem.autodiff %import NonlinearForm
%from skfem.autodiff.helpers import grad, dot, div
%import numpy as np
%import jax.numpy as jnp

%m = MeshTri().refined(5)

%@NonlinearForm(hessian=True)
%def J(u, _):
%    gamma = 1000
%    return (1/2 * dot(grad(u), grad(u)) - u
%            - (1/2) * (1/gamma) * (0 + 1)**2
%            + (gamma / 2) * (jnp.maximum(u.value - 0.01 - (1/gamma) * (0 + 1),0))**2)

%basis = Basis(m, ElementTriP1())

%D = m.boundary_nodes()
%x = np.zeros(basis.N)

%for itr in range(100):
%    x_prev = x.copy()
%    x += solve(*condense(*J.assemble(basis, x=x), D=D))
%    res = np.linalg.norm(x - x_prev)
%    if res < 1e-8:
%        break
%    if __name__ == "__main__":
%        print(res)
%
%if __name__ == "__main__":
%    from skfem.visuals.matplotlib import plot3, show
%    plot3(m, x)
%    show()
%    
%    ******

\subsection{Comentarios adicionales}

%Aquí unos comentarios para profundizar el tema de las desigualdades variacionales.

%\paragraph{Mas desigualdades variacionales}

% Projets

%\paragraph{Aplicaciones del obstaculo}

%\paragraph{Mas tecnicas de aproximacion}

\begin{itemize}
    \item 
    
    
    \textbf{Método de Nitsche o multiplicadores}: Estas técnicas alternativas permiten incorporar las restricciones del problema de manera explícita en la formulación débil, mejorando la estabilidad numérica en ciertos casos. El método de Nitsche, en particular, ofrece una forma elegante de manejar las condiciones de desigualdad.
    \item \textbf{Aplicaciones}: El problema del obstáculo tiene aplicaciones en mecánica de sólidos, fluidos, finanzas y otras áreas de la ciencia e ingeniería. Su versatilidad lo convierte en un modelo fundamental para estudiar fenómenos con restricciones geométricas o físicas.
\end{itemize}

%\addcontentsline{toc}{chapter}{Bibliograf\'ia}
%\nocite{*}
%{\footnotesize
%\bibliographystyle{siam}
%\bibliography{biblio}
%}

%\addcontentsline{toc}{chapter}{\'Indice alfab\'etico}
%{\scriptsize \printindex}

%\end{document} 


\chapter{Dualidad}
La dualidad en optimización es un concepto fundamental que establece una relación entre un problema de optimización primal y su problema dual asociado. El problema dual proporciona una cota inferior para el valor óptimo del problema primal, lo que permite analizar la solidez de las soluciones encontradas. Las condiciones de optimalidad de Karush-Kuhn-Tucker (KKT) conectan las soluciones primales y duales. Este enfoque es especialmente útil en problemas con restricciones complejas, donde el análisis directo puede ser difícil. La dualidad también tiene interpretaciones geométricas profundas.
%
Seguimos~\cite{allaire2007} (véase también~\cite{BoydVandenberghe2004} y \cite{luenberger1969}).

\section{Problema primal y problema dual}

%Seguimos Allaire.
%%%
% ALLAIRE PAGE 320
%%%

\subsection{Definiciones}

\begin{definition}
Sean \( V \) y \( Q \) dos espacios de Hilbert sobre el cuerpo \( \mathbb{K} \) (donde \( \mathbb{K} = \mathbb{R} \) o \( \mathbb{C} \)), y sea \( \mathcal{L} \colon U \times P \subseteq V \times Q \to \mathbb{K} \) un \emph{Lagrangiano} definido sobre un subconjunto \( U \times P \) de \( V \times Q \).

Un \emph{punto silla} de \( \mathcal{L} \) es un par \( (v^*, q^*) \in U \times P \) tal que:
\[
\mathcal{L}(v^*, q) \leq \mathcal{L}(v^*, q^*) \leq \mathcal{L}(v, q^*) \quad \forall (v, q) \in U \times P.
\]
Es decir, \( v^* \) minimiza \( \mathcal{L}(\cdot, q^*) \) en \( U \) y \( q^* \) maximiza \( \mathcal{L}(v^*, \cdot) \) en \( P \).
\end{definition}

---
\begin{definition}
Dado el Lagrangiano \( \mathcal{L} \), definimos las \emph{funciones primal y dual} como:
\[
P(v) = \sup_{q \in P} \mathcal{L}(v, q) \quad \text{(función primal)},
\]
\[
D(q) = \inf_{v \in U} \mathcal{L}(v, q) \quad \text{(función dual)}.
\]
\end{definition}

---
\begin{definition}
Los \emph{problemas primal y dual} asociados a \( \mathcal{L} \) son:
\[
\text{Problema primal:} \quad \inf_{v \in U} P(v) = \inf_{v \in U} \sup_{q \in P} \mathcal{L}(v, q),
\]
\[
\text{Problema dual:} \quad \sup_{q \in P} D(q) = \sup_{q \in P} \inf_{v \in U} \mathcal{L}(v, q).
\]
\end{definition}

\chapter{Perspectivas}
\label{ch:conclusion}

Para concluir nuestra breve introducción, damos unas sugerencias que permitirán al lector ir más allá en el tema y/o descubrir otros campos relacionados. Se puede referir tambi\'en a \cite{bertsekas1982}.
% References plus generales
% ROCKAFELLAR
% EKELAND et TEMAM

\section{Optimización matemática}
El análisis convexo es fundamental en la teoría de la optimización, especialmente en problemas de programación matemática. Las funciones convexas y los conjuntos convexos permiten modelar y resolver problemas de minimización o maximización con restricciones, garantizando la existencia de soluciones globales óptimas. Por ejemplo, en \textbf{programación convexa}, se minimiza una función convexa sobre un conjunto convexo, lo que asegura convergencia a soluciones globales.
\begin{align*}
    \text{Minimizar } & f(x) \\
    \text{Sujeto a } & g_i(x) \leq 0, \quad i = 1, \dots, m \\
                     & h_j(x) = 0, \quad j = 1, \dots, p,
\end{align*}
donde \( f \) es convexa y \( g_i \), \( h_j \) son funciones convexas.

\section{Control y sistemas dinámicos}
En teoría de control, el análisis convexo se aplica para estabilizar sistemas dinámicos y diseñar controladores robustos. Los \textbf{conos tangentes} y las \textbf{funciones de Lyapunov convexas} son herramientas clave para analizar la estabilidad asintótica y diseñar estrategias de control óptimo.

\section{Aprendizaje automático}
En aprendizaje automático, el análisis convexo es esencial en algoritmos como \textbf{Support Vector Machines (SVM)} y \textbf{regresión logística}. Las SVM, por ejemplo, buscan el hiperplano que maximice el margen entre clases, lo que se formula como un problema de optimización convexa. La convexidad garantiza que el algoritmo encuentre la solución globalmente óptima.
\begin{equation*}
    \min_{w, b} \frac{1}{2} \|w\|^2 \quad \text{sujeto a} \quad y_i (w \cdot x_i + b) \geq 1 \quad \forall i.
\end{equation*}

\section{Economía y finanzas}
En economía, el análisis convexo se utiliza para modelar preferencias de agentes económicos, funciones de producción y problemas de equilibrio general. Por ejemplo, las \textbf{funciones de utilidad cuasi-convexas} permiten analizar decisiones bajo incertidumbre, mientras que la convexidad de las funciones de costo es clave en la teoría de la firma.

\section{Teoría de juegos}
En teoría de juegos, los conceptos de convexidad son centrales para analizar equilibrios de Nash en juegos no cooperativos. Las estrategias mixtas y las funciones de pago convexas permiten demostrar la existencia de equilibrios y estudiar su estabilidad.

\section{Geometría y computación}
% IVR
En geometría computacional, el análisis convexo se aplica en algoritmos de \textbf{envueltas convexas} (como el algoritmo de Graham) y en problemas de localización de puntos. La convexidad permite diseñar algoritmos eficientes para encontrar soluciones aproximadas o exactas en conjuntos de datos de alta dimensionalidad.

\section{Procesamiento de señales e imágenes}
En procesamiento de imágenes, técnicas como la \textbf{reconstrucción convexa} y la \textbf{regularización total variación (TV)} se basan en la convexidad para minimizar artefactos y mejorar la calidad de la imagen. Estas técnicas son fundamentales en tomografía, resonancia magnética y reconstrucción 3D.

\section{Teoría de la información}
En teoría de la información, las funciones convexas como la \textbf{entropía relativa} (o divergencia de Kullback-Leibler) son herramientas esenciales para analizar la capacidad de canales de comunicación y diseñar códigos eficientes.

%\addcontentsline{toc}{chapter}{Bibliograf\'ia}
%\nocite{*}
%{\footnotesize
%\bibliographystyle{siam}
%\bibliography{biblio}
%}

%\addcontentsline{toc}{chapter}{\'Indice alfab\'etico}
%{\scriptsize \printindex}

%\end{document} 

%\chapter{********************}

%\chapter{Algoritmos de minimización}

%Los algoritmos de minimización son herramientas fundamentales en matemática. % aplicada, optimización, aprendizaje automático y ciencia de datos. 
%Este capítulo aborda los %principios y dos 
%algoritmos más relevantes para resolver problemas de minimización de funciones, desde enfoques clásicos hasta técnicas avanzadas basadas en análisis convexo y métodos numéricos. Cada sección incluye las fórmulas clave de los métodos y, cuando es posible, su relación con ecuaciones diferenciales que modelan la dinámica de la minimización.

%\section{Ejemplos}

%Funcional cuadrática. Ver \cite{allaire2007}.

%\section{Gradiente descendente}
%El \textbf{gradiente descendente} es un método iterativo para minimizar una función \( f: \mathbb{R}^n \to \mathbb{R} \) diferenciable. El algoritmo se define como:
%\[
%x_{k+1} = x_k - \alpha_k \nabla f(x_k),
%\]
%donde \( \alpha_k > 0 \) es el tamaño de paso (o tasa de aprendizaje) y \( \nabla f(x_k) \) es el gradiente de \( f \) en \( x_k \).

%\begin{quotation}
%\textbf{Vínculo con ecuaciones diferenciales:}
%El método de gradiente descendente puede interpretarse como una discretización de la ecuación diferencial ordinaria (EDO) asociada al flujo gradiente:
%\[
%\frac{dx}{dt} = -\nabla f(x(t)).
%\]
%La solución de esta EDO, bajo condiciones adecuadas, converge a un mínimo local de \( f \).
%\end{quotation}

\appendix

\chapter{M\'as sobre conjuntos convexos} %Cáscara convexa y aplicaciones en geometría algorítmica}

%***
%Los conjuntos convexos son fundamentales en matemática, especialmente en geometría algorítmica, optimización y aprendizaje automático. Este capítulo explora la teoría de conjuntos convexos, con énfasis en la \textbf{cáscara convexa} (o \emph{convex hull}), sus propiedades, algoritmos para su cálculo y aplicaciones en geometría computacional. 
%***

Los conjuntos convexos son fundamentales en geometría computacional, optimización y aprendizaje automático. Su estudio permite resolver problemas complejos mediante estructuras bien definidas y propiedades robustas. La importancia de estos conjuntos radica en su capacidad para modelar fenómenos reales de manera eficiente, desde la planificación de rutas hasta la clasificación de datos. Además, su definición intuitiva —un conjunto donde el segmento entre dos puntos cualesquiera siempre está contenido— facilita su aplicación en algoritmos y teoremas. En la práctica, los conjuntos convexos aparecen en problemas de minimización, donde garantizan la existencia de soluciones óptimas. Su utilidad se extiende también a la robótica, el procesamiento de imágenes y la inteligencia artificial, donde sirven como herramientas clave para el análisis y la toma de decisiones.

La cáscara convexa, (\emph{convex hull}), es una de las aplicaciones más relevantes de estos conjuntos. Se trata del conjunto convexo más pequeño que contiene a un grupo dado de puntos, actuando como un envoltorio geométrico que preserva la estructura esencial del conjunto original. Este concepto es crucial en geometría computacional, donde permite resolver problemas como la detección de colisiones o la reconstrucción de superficies en tres dimensiones. En el aprendizaje automático, la cáscara convexa se utiliza para definir regiones de decisión en algoritmos de clasificación, como las máquinas de vectores de soporte. Su cálculo eficiente es posible gracias a algoritmos como el de Graham o QuickHull, que optimizan el proceso reduciendo la complejidad computacional. Además, su visualización intuitiva ayuda a comprender mejor la distribución de datos en espacios multidimensionales. 
Este capítulo explora tanto la teoría como las aplicaciones prácticas de los conjuntos convexos y la cáscara convexa. A través de ejemplos concretos y algoritmos, se busca proporcionar al lector una buena comprensión de estos conceptos. 
%Se analizan casos de uso en optimización, donde la convexidad garantiza la convergencia de métodos iterativos como el de Newton. También se discuten aplicaciones en robótica, donde la cáscara convexa facilita la planificación de movimientos en entornos con obstáculos. En el ámbito del aprendizaje automático, se muestra cómo estos conjuntos permiten separar clases de datos de manera eficiente. Finalmente, se destacan las implicaciones de estos conceptos en la resolución de problemas reales, subrayando su relevancia en la ciencia de datos y la ingeniería moderna.

\section{Definición y propiedades}

%*** % Reprendre ici 
Los conjuntos convexos son uno de los conceptos más %versátiles y 
poderosos en matemáticas, con aplicaciones que abarcan desde la optimización hasta el aprendizaje automático. Un conjunto es convexo si, intuitivamente, no presenta huecos ni entrantes: cualquier segmento de recta que una dos puntos del conjunto permanece completamente dentro de él. Esta propiedad, aparentemente simple, tiene implicaciones profundas en teoría y práctica, permitiendo resolver problemas de minimización, clasificación y modelado con garantías de eficiencia y robustez.

\subsection{Para empezar}

En el marco de la geometría y el análisis funcional, los conjuntos convexos emergen como objetos fundamentales, definidos por la propiedad intrínseca de contener, junto con cualquier par de sus puntos, el segmento de recta que los une. %Este concepto, sencillo en apariencia, sienta las bases para explorar estructuras geométricas más complejas y su aplicabilidad en áreas tan diversas como la optimización, la teoría de juegos o incluso la inteligencia artificial.

\begin{definition}[Conjunto convexo]
    Un conjunto \( C \subseteq \mathbb{R}^n \) es \textbf{convexo} si, para todo \( x, y \in C \) y todo \( \lambda \in [0,1] \), se cumple que:
    \[
    \lambda x + (1-\lambda) y \in C.
    \]
    En otras palabras, el conjunto contiene todos los puntos del segmento de recta que une \( x \) e \( y \).
\end{definition}

\begin{ex}
\textbf{\textit{Ejemplos clásicos de conjuntos convexos:}}
\begin{itemize}
    \item Un \textbf{segmento de recta} en \( \mathbb{R}^n \): \([a, b] = \{ \lambda a + (1-\lambda) b \mid \lambda \in [0,1] \}\).
    \item Un \textbf{disco} (o bola cerrada) en \( \mathbb{R}^2 \): \( \{ (x,y) \mid x^2 + y^2 \leq r^2 \} \).
    \item Un \textbf{hiperplano} en \( \mathbb{R}^n \): \( \{ x \mid a^T x = b \} \), donde \( a \in \mathbb{R}^n \) y \( b \in \mathbb{R} \).
    \item El \textbf{espacio vectorial} \( \mathbb{R}^n \) en sí mismo.
\end{itemize}
\end{ex}

\begin{theorem}[Intersección de conjuntos convexos]
    Sea \( \{C_i\}_{i \in I} \) una familia de conjuntos convexos en \( \mathbb{R}^n \). Entonces, su intersección \( \bigcap_{i \in I} C_i \) también es un conjunto convexo.
\end{theorem}

\begin{proof}
    Sea \( x, y \in \bigcap_{i \in I} C_i \) y \( \lambda \in [0,1] \). Por definición de intersección, \( x, y \in C_i \) para todo \( i \in I \). Como cada \( C_i \) es convexo, \[ \lambda x + (1-\lambda) y \in C_i \] para todo \( i \in I \). Por lo tanto, \( \lambda x + (1-\lambda) y \in \bigcap_{i \in I} C_i \), lo que prueba que la intersección es convexa.
\end{proof}

\begin{definition}[Combinación convexa]
    Una \textbf{combinación convexa} de puntos \( x_1, x_2, \dots, x_k \in \mathbb{R}^n \) es una expresión de la forma:
    \[
    \sum_{i=1}^k \lambda_i x_i, \quad \text{donde} \quad \lambda_i \geq 0 \quad \text{y} \quad \sum_{i=1}^k \lambda_i = 1.
    \]
\end{definition}

\subsection{Interior y clausura de un conjunto convexo}

%En esta sección 
Ahora
exploramos propiedades esenciales de los conjuntos convexos: primero, su comportamiento respecto al interior y la clausura. % su relación con los hiperplanos de separación, y sus características topológicas.

\begin{proposition}
    Sea \( C \subseteq \mathbb{R}^n \) un conjunto convexo. Entonces, su interior \( \mathring{C} \) y su clausura \( \overline{C} \) son también conjuntos convexos.
\end{proposition}

\begin{proof} Veamos ambos casos.

%\section*{Demostración}

%\begin{proof}
Sean \( C \subseteq \mathbb{R}^n \) un conjunto convexo y \( \mathring{C} \) su interior. Queremos demostrar que \( \mathring{C} \) es convexo, es decir, que para todo \( x, y \in \mathring{C} \) y todo \( \lambda \in [0,1] \), el punto \( z = \lambda x + (1-\lambda)y \) pertenece a \( \mathring{C} \).

\begin{enumerate}
    \item \textbf{Premisas:}
%    \begin{itemize}
%        \item 
Sean \( x, y \in \mathring{C} \). Por definición de punto interior, existen bolas abiertas \( B(x, r_x) \subseteq C \) y \( B(y, r_y) \subseteq C \), donde \( r_x, r_y > 0 \).
%        \item 
Sea \( \lambda \in [0,1] \).
%        \item 
Definimos \( z = \lambda x + (1-\lambda)y \).
%    \end{itemize}

    \item \textbf{Construcción de la bola centrada en \( z \):}
%    \begin{itemize}
%        \item 
Sea \( r = \min(r_x, r_y) > 0 \).
%        \item 
Consideremos la bola abierta \( B(z, r) \).
%\item 
Para todo \( w \in B(z, r) \), escribimos \( w = z + v \), donde \( \|v\| < r \).
%    \end{itemize}

    \item \textbf{Expresión de \( w \) como combinación convexa:}
%    \begin{itemize}
%        \item 
Observamos que:
        \[
        w = z + v = \lambda x + (1-\lambda)y + v.
        \]
%        \item 
Podemos descomponer \( v \) como:
        \[
        v = \lambda v + (1-\lambda)v,
        \]
        donde \( \|v\| < r \leq r_x \) y \( \|v\| < r \leq r_y \). Esto es posible porque \( r = \min(r_x, r_y) \).
%    \end{itemize}

    \item \textbf{Verificación de que \( w \in C \):}
%    \begin{itemize}
 %       \item 
 Como \( B(x, r_x) \subseteq C \) y \( B(y, r_y) \subseteq C \), entonces:
        \[
        x + v \in B(x, r_x) \subseteq C \quad \text{y} \quad y + v \in B(y, r_y) \subseteq C.
        \]
%        \item 
Por convexidad de \( C \), la combinación convexa:
        \[
        \lambda(x + v) + (1-\lambda)(y + v) = z + v = w
        \]
        también está en \( C \).
%    \end{itemize}

%    \item \textbf{Conclusión:}
%    \begin{itemize}
%        \item 
Dado que para todo \( w \in B(z, r) \) se tiene \( w \in C \), entonces \( B(z, r) \subseteq C \).
%        \item 
Esto implica que \( z \in \mathring{C} \).
%\item 
Como \( z \) es un punto arbitrario en el segmento \( [x, y] \), se concluye que \( [x, y] \subseteq \mathring{C} \).
%    \end{itemize}
\end{enumerate}
Por lo tanto, \( \mathring{C} \) es un conjunto convexo.
%\end{proof}

\medskip

    \textbf{Caso de la clausura \( \overline{C} \):}
    Sean \( x, y \in \overline{C} \) y \( \lambda \in [0,1] \). Existen sucesiones \( \{x_k\} \) e \( \{y_k\} \) en \( C \) tales que \( x_k \to x \) e \( y_k \to y \). Definimos \( z_k = \lambda x_k + (1-\lambda)y_k \). Como \( C \) es convexo, \( z_k \in C \) para todo \( k \). Además, \( z_k \to \lambda x + (1-\lambda)y = z \), por lo que \( z \in \overline{C} \).
\end{proof}

%\subsection{Teorema de separación por un hiperplano}

%\begin{definition}[Hiperplano]
%    Un hiperplano en \( \mathbb{R}^n \) es un conjunto de la forma:
%    \[
%    H = \{ x \in \mathbb{R}^n \mid a^T x = b \},
%    \]
%    donde \( a \in \mathbb{R}^n \setminus \{0\} \) y \( b \in \mathbb{R} \).
%\end{definition}

%\begin{theorem}[Separación por un hiperplano]
%    Sean \( C_1 \) y \( C_2 \) dos conjuntos convexos no vacíos en \( \mathbb{R}^n \), con \( C_1 \) compacto y \( C_2 \) cerrado. Si \( C_1 \cap C_2 = \emptyset \), entonces existe un hiperplano que separa estrictamente \( C_1 \) y \( C_2 \). Es decir, existen \( a \in \mathbb{R}^n \) y \( b \in \mathbb{R} \) tales que:
%    \[
%    \sup_{x \in C_1} a^T x < \inf_{y \in C_2} a^T y.
%    \]
%\end{theorem}

%\begin{proof}
%    Consideremos la función \( f(x, y) = \|x - y\|^2 \), definida para \( x \in C_1 \) e \( y \in C_2 \). Como \( C_1 \) es compacto y \( C_2 \) es cerrado, \( f \) alcanza su mínimo en un punto \( (x_0, y_0) \in C_1 \times C_2 \). Definimos el hiperplano:
%    \[
%    H = \{ z \in \mathbb{R}^n \mid (x_0 - y_0)^T (z - x_0) = 0 \}.
%    \]
%    Como \( C_1 \cap C_2 = \emptyset \), \( x_0 \neq y_0 \). Para todo \( x \in C_1 \), se cumple:
%    \[
%    (x - x_0)^T (x_0 - y_0) \leq 0,
%    \]
%    lo que implica:
%    \[
%    x^T (x_0 - y_0) \leq x_0^T (x_0 - y_0).
%    \]
%    De manera similar, para todo \( y \in C_2 \):
%    \[
%    y^T (x_0 - y_0) \geq y_0^T (x_0 - y_0).
%    \]
%    Por lo tanto, el hiperplano \( H \) separa estrictamente \( C_1 \) y \( C_2 \).
%\end{proof}

%---

%\subsection{Propiedades topológicas de los conjuntos convexos}

%\begin{definition}[Conjunto cerrado y compacto]
%    Un conjunto \( C \subseteq \mathbb{R}^n \) es:
%    \begin{itemize}
%        \item \textbf{Cerrado} si contiene todos sus puntos límite.
%        \item \textbf{Compacto} si es cerrado y acotado.
%    \end{itemize}
%\end{definition}

%\begin{proposition}[Convexidad y topología]
%    Sea \( C \subseteq \mathbb{R}^n \) un conjunto convexo. Entonces:
%    \begin{enumerate}
%        \item \( C \) es \textbf{conexo por arcos}: para todo \( x, y \in C \), existe un arco continuo \( \gamma: [0,1] \to C \) tal que \( \gamma(0) = x \) y \( \gamma(1) = y \).
%        \item \( C \) es \textbf{arcoconexo}: cualquier par de puntos puede unirse por una curva continua contenida en \( C \).
%        \item Si \( C \) es cerrado, su interior \( \mathring{C} \) es denso en \( C \) (es decir, \( \overline{\mathring{C}} = C \)).
%    \end{enumerate}
%\end{proposition}

%\begin{proof}
%    \begin{enumerate}
%        \item Para \( x, y \in C \), el segmento \( \gamma(t) = (1-t)x + ty \) es un arco continuo contenido en \( C \) (por convexidad).
%        \item Consecuencia directa de la convexidad: cualquier curva poligonal que una \( x \) e \( y \) está contenida en \( C \).
%        \item Si \( C \) es cerrado, para todo \( x \in C \) y todo \( \epsilon > 0 \), existe \( y \in \mathring{C} \) tal que \( \|x - y\| < \epsilon \). Esto se debe a que \( C \) es la clausura de su interior.
%    \end{enumerate}
%\end{proof}

%---

%\section{Cono convexo}
%\chapter{Conos Convexos}
\section{Conos convexos}
%\section{Conos Convexos: Definición y Propiedades}

\subsection{Definición}

Un \textbf{cono} en un espacio vectorial real \( V \) es un subconjunto \( K \subseteq V \) que satisface:
\begin{itemize}
    \item Si \( x \in K \) y \( \lambda \geq 0 \), entonces \( \lambda x \in K \).
    \item Si \( x, y \in K \), entonces \( x + y \in K \).
\end{itemize}

Un cono \( K \) es \textbf{convexo} si es convexo como conjunto, es decir, si para todo \( x, y \in K \) y \( \lambda \in [0,1] \), se cumple:
\[
\lambda x + (1 - \lambda) y \in K.
\]
En particular, todo cono convexo es cerrado bajo combinaciones lineales no negativas de sus elementos.

%\begin{definicion}[Cono Convexo]
Sea \( V \) un espacio vectorial real. Un subconjunto \( K \subseteq V \) es un \textbf{cono convexo} si:
\begin{enumerate}
    \item \( K \) es un cono: \( \lambda x \in K \) para todo \( x \in K \) y \( \lambda \geq 0 \).
    \item \( K \) es convexo: \( \lambda x + (1 - \lambda) y \in K \) para todo \( x, y \in K \) y \( \lambda \in [0,1] \).
\end{enumerate}
%\end{definicion}

\subsection{Ejemplos}

%\begin{observacion}[Ejemplos de Conos Convexos]
\begin{enumerate}
    \item \textbf{Cono positivo en \( \mathbb{R}^n \)}:
    \[
    \mathbb{R}^n_+ = \{ x \in \mathbb{R}^n \mid x_i \geq 0 \text{ para todo } i = 1, \ldots, n \}.
    \]
    Es un cono convexo, ya que la suma de vectores no negativos es no negativo y el producto por escalares no negativos preserva la no negatividad.

    \item \textbf{Cono de Lorentz} (o cono de segundo orden):
    \[
    L^n = \{ (x_0, x_1, \ldots, x_n) \in \mathbb{R}^{n+1} \mid x_0 \geq \sqrt{x_1^2 + \cdots + x_n^2} \}.
    \]
    Este cono es convexo y juega un papel importante en optimización convexa.

    \item \textbf{Cono de matrices semidefinidas positivas}:
    \[
    \mathbb{S}^n_+ = \{ X \in \mathbb{R}^{n \times n} \mid X = X^\top \text{ y } x^\top X x \geq 0 \text{ para todo } x \in \mathbb{R}^n \}.
    \]
    Es un cono convexo en el espacio de matrices simétricas.
\end{enumerate}
%\end{observacion}

\subsection{Propiedades principales}

%\begin{teorema}[Propiedades de Conos Convexos]
Sea \( K \subseteq V \) un cono convexo. Entonces:
\begin{enumerate}
    \item \( K \) es cerrado bajo combinaciones lineales no negativas:
    \[
    \sum_{i=1}^m \lambda_i x_i \in K \quad \text{para todo } x_i \in K, \lambda_i \geq 0, m \in \mathbb{N}.
    \]
    \item El \textbf{dual} de \( K \), definido como
    \[
    K^* = \{ y \in V^* \mid \langle y, x \rangle \geq 0 \text{ para todo } x \in K \},
    \]
    también es un cono convexo.
    \item Si \( K \) es cerrado (en la topología apropiada), entonces \( (K^*)^* = K \).
\end{enumerate}
%\end{teorema}

%\begin{observacion}
\begin{itemize}
    \item La propiedad de ser cerrado bajo combinaciones lineales no negativas es equivalente a la convexidad del cono.
    \item El dual \( K^* \) permite caracterizar desigualdades lineales y es fundamental en programación convexa.
\end{itemize}
%\end{observacion}


%\section{Definición y caracterización de la envoltura convexa}
\section{Envoltura convexa}

La \textbf{envoltura convexa} de un conjunto \( S \subseteq \mathbb{R}^n \) es un concepto fundamental en análisis convexo y optimización. Permite caracterizar los conjuntos convexos minimales que contienen a \( S \), y entender sus propiedades geométricas y analíticas.

\subsection{Caracterizaci\'on}

Una definici\'on posible de la envoltura (o c\'ascara) convexa es la siguiente.

\begin{definition}[Caracterización de la envoltura convexa]
Sea \( S \subseteq \mathbb{R}^n \). Entonces,
\[
\envolt(S) = \left\{ \sum_{i=1}^k \lambda_i x_i \ \middle|\ x_i \in S, \lambda_i \geq 0, \sum_{i=1}^k \lambda_i = 1, k \in \mathbb{N} \right\},
\]
es un conjunto convexo que contiene a \( S \), y es el más pequeño en el sentido de la inclusión.
\end{definition}

%Esta definición es en realidad un \textbf{teorema} 
El teorema siguiente ilustra %
% ya 
que existen dos caracterizaciones equivalentes de \( \envolt(S) \): una por combinación convexa, y otra por intersección de conjuntos convexos que contienen a \( S \) (véase
\cite{BoydVandenberghe2004}, 
\cite{hiriart2001}, \cite{Rockafellar1970} por la demonstraci\'on).


\begin{theorem}[Equivalencia por intersección]
Además, la envoltura convexa de \( S \) satisface la siguiente caracterización:
\[
\envolt(S) = \bigcap \{ C \subseteq \mathbb{R}^n \mid \envolt(C) = C \text{ y } S \subseteq C \}.
\]
\end{theorem}

El resultado siguiente, conocido como la \emph{propiedad de los vértices}, también es un teorema (véase \cite{Grünbaum1963}).

\begin{theorem}[Propiedad de los vértices en dimensión 2]
En dimensión 2, si \( S \) es un conjunto finito de puntos en \( \mathbb{R}^2 \), entonces \( \envolt(S) \) es un polígono convexo cuyos vértices son un subconjunto de \( S \). %Este resultado, conocido como la "propiedad de los vértices", también es un teorema (véase \cite{Grünbaum1963}).
\end{theorem}

%\begin{thebibliography}
%\bibitem{Rockafellar1970}
%  Rockafellar, R. T. (1970).
%  *Análisis convexo*.
%  Princeton University Press.
%  \url{https://press.princeton.edu/books/paperback/9780691001081/convex-analysis}

%\bibitem{BoydVandenberghe2004}
%  Boyd, S., & Vandenberghe, L. (2004).
%  *Optimización convexa*.
%  Cambridge University Press.
%  \url{https://web.stanford.edu/~boyd/cvxbook/}

%\bibitem{Grünbaum1963}
%  Grünbaum, B. (1963).
%  *Poliedros convexos*.
%  Tractos de Matemática Pura de Interscience.
%  \url{https://www.ams.org/books/itm/63/}
%\end{thebibliography}

La figura~\ref{fig:pentagone} ilustra la envoltura convexa de un conjunto de puntos en \(\mathbb{R}^2\). En ella, los vértices del pentágono convexo están marcados con cuadrados negros, los puntos interiores con rombos rojos, y los puntos originales del conjunto \( S \) con círculos negros.

\begin{figure}[h!]
\centering
\begin{tikzpicture}[
    scale=1.5,
    point/.style={circle, fill=black, inner sep=0pt, minimum size=6pt},
    sommet/.style={rectangle, fill=black!80, inner sep=0pt, minimum size=8pt},
    interieur/.style={diamond, fill=red, inner sep=0pt, minimum size=8pt, draw=white, thick},
    epaisseur/.style={line width=1.5pt},
    ]

    % Conjunto de puntos S (cercles noirs)
    \foreach \x/\y in {0.5/1, 1.5/2.5, 2.5/2.5, 3.5/1, 2/0, 1.1/1.5, 2.5/1} {
        \node[point] at (\x,\y) {};
    }

    % Sommets du pentagone (carrés noirs)
    \foreach \x/\y in {0.5/1, 1.5/2.5, 2.5/2.5, 3.5/1, 2/0} {
        \node[sommet] at (\x,\y) {};
    }

    % Puntos dentro de la cáscara (losanges rojos avec contour blanc pour mieux voir)
    \foreach \x/\y in {1.1/1.5, 2.5/1, 2/1.5, 3/1.2} {
        \node[interieur] at (\x,\y) {};
    }

    % Pentagone convexe (enveloppe convexa)
    \draw[epaisseur, blue!80!black, %fill=blue!5, 
    rounded corners=2pt] (0.5,1) -- (1.5,2.5) -- (2.5,2.5) -- (3.5,1) -- (2,0) -- cycle;

    % Label
    \node[above, blue!80!black, font=\bfseries] at (2,2.8) {Envoltura convexa de \( S \)};

    % Leyendas
    \node[below] at (2,-0.3) {Puntos en \( S \)};
    \node[red!80!black, right] at (3.5,1.2) {Puntos en \( \mathrm{conv}(S) \)};
\end{tikzpicture}
\caption{Visualización de la envoltura convexa de un conjunto de puntos en \( \mathbb{R}^2 \).} %Los puntos negros son los vértices del pentágono convexo, los rojos están en su interior, y los círculos son puntos del conjunto original \( S \).}
\label{fig:pentagone}
\end{figure}

%\section{Aplicaciones}

%\subsection{Envoltura convexa de un conjunto de puntos}

%La \textbf{envoltura convexa} (o cáscara convexa) de un conjunto finito de puntos en el plano es el polígono convexo más pequeño que contiene a todos los puntos. Este concepto es esencial en:

%\begin{itemize}
%    \item \textbf{Geometría computacional}: Para problemas de colisión, localización de puntos y clustering.
%    \item \textbf{Visión por computadora}: En reconstrucción 3D y reconocimiento de formas.
%    \item \textbf{Robótica}: Para planificación de caminos y evitación de obstáculos.
%    \item \textbf{Aprendizaje automático}: En clustering (por ejemplo, el algoritmo DBSCAN).
%\end{itemize}

%\subsection{Separación de conjuntos convexos}

%\begin{theorem}[Teorema de separación de hiperplanos]
%Si \( C \) y \( D \) son dos conjuntos convexos disjuntos en \( \mathbb{R}^n \), entonces existe un hiperplano \( H \) que separa \( C \) y \( D \), es decir, existe \( a \in \mathbb{R}^n \) y \( b \in \mathbb{R} \) tales que
%\[
%a^T x \leq b \quad \forall x \in C, \quad a^T y \geq b \quad \forall y \in D.
%\]
%\end{theorem}

%Este teorema es la base de algoritmos de separación en optimización convexa y aprendizaje automático (por ejemplo, máquinas de vectores de soporte).

\subsection{Técnicas para calcular la envoltura convexa}
\label{subsec:techniques_enveloppe}

A lo largo de los años, se han desarrollado numerosos métodos para calcular la envoltura convexa de un conjunto de puntos en el plano o en espacios de mayor dimensión. Las técnicas más antiguas, como el \textbf{método de fuerza bruta}, consisten en comparar todos los subconjuntos posibles de puntos, lo que resulta en una complejidad computacional extremadamente alta, del orden de \( O(n^4) \) en dos dimensiones. Una mejora significativa llegó con el algoritmo de \textbf{envuelto de regalo (Jarvis March, 1973)} %\cite{Jarvis1973}, 
que construye la envoltura convexa recorriendo los puntos en orden angular, alcanzando una complejidad de \( O(nh) \), donde \( n \) es el número total de puntos y \( h \) el número de puntos en la envoltura. Con el avance de la computación y la necesidad de mayor eficiencia, surgieron algoritmos más sofisticados, como el \textbf{scan de Graham (1972)}, %\cite{Graham1972}, 
que ordena los puntos por ángulo polar y construye la envoltura en \( O(n \log n) \), y el \textbf{QuickHull}, %\cite{Barber1996}, 
que emplea una estrategia de "divide y vencerás" con una complejidad promedio similar. Entre las técnicas modernas destacan algoritmos híbridos como el de \textbf{Chan} %\cite{Chan1996}, 
que combina ideas de Graham y QuickHull para lograr una complejidad de \( O(n \log h) \), o el \textbf{Andrew's monotone chain} %\cite{Andrew1979}, 
una variante optimizada del scan de Graham para conjuntos de puntos ordenados. Además, en contextos específicos o dimensiones superiores, se recurre a métodos basados en programación lineal o técnicas estocásticas. %Estas soluciones reflejan la evolución desde enfoques intuitivos pero lentos hacia algoritmos altamente optimizados, demostrando los avances en complejidad algorítmica y geometría computacional.

\subsection{Algoritmo de Graham Scan} %para el cálculo de la envoltura convexa}
\label{subsec:graham}

El \textbf{algoritmo de Graham Scan} %\cite{Graham1972} 
es un método eficiente para calcular la \textbf{envoltura convexa} de un conjunto finito de puntos en el plano. Su complejidad es de \( O(n \log n) \), donde \( n \) es el número de puntos en el conjunto \( S \).

\subsubsection{Idea general}
El objetivo es encontrar el polígono convexo más pequeño que contiene todos los puntos de \( S \subseteq \mathbb{R}^2 \). La idea central es:
\begin{itemize}
    \item \textbf{Ordenar los puntos} por su ángulo polar respecto a un \textbf{punto base} (generalmente el de menor coordenada \( y \)).
    \item \textbf{Construir la envoltura} recorriendo los puntos ordenados y eliminando aquellos que no contribuyen a la convexidad.
\end{itemize}

\subsubsection{Etapas del algoritmo}

\begin{enumerate}
    \item \textbf{Selección del punto base \( p_0 \):}
    \begin{itemize}
        \item Elegir el punto con la \textbf{coordenada \( y \) mínima} (y \( x \) mínima en caso de empate).
        \item Este punto \textbf{siempre pertenece} a la envoltura convexa.
    \end{itemize}

    \item \textbf{Ordenamiento por ángulo polar:}
    \begin{itemize}
        \item Calcular el ángulo que forma cada punto \( p_i \) con el punto base \( p_0 \) usando la función \texttt{atan2}.
        \item Ordenar los puntos en \textbf{orden creciente de ángulo polar} respecto a \( p_0 \).
        \item Si dos puntos tienen el mismo ángulo, conservar el más lejano.
    \end{itemize}

    \item \textbf{Construcción de la envoltura con una pila:}
    \begin{itemize}
        \item Inicializar una pila con los primeros tres puntos ordenados (\( p_0 \), \( p_1 \), \( p_2 \)).
        \item Para cada punto \( p_i \) (desde \( p_3 \) hasta \( p_{n-1} \)):
        \begin{enumerate}
            \item Mientras el ángulo formado por los últimos tres puntos en la pila sea \textbf{mayor a 180°} (es decir, hay un "doblez hacia adentro"), eliminar el punto del medio de la pila.
            \item Empujar \( p_i \) en la pila.
        \end{enumerate}
        \item Para verificar la convexidad entre tres puntos consecutivos \( A \), \( B \), \( C \) en la pila, calcular el \textbf{producto cruzado}:
        \[
        \text{Producto cruzado} = (B_x - A_x)(C_y - B_y) - (B_y - A_y)(C_x - B_x)
        \]
        \begin{itemize}
            \item Si el producto cruzado es \textbf{negativo}, el ángulo en \( B \) es mayor a 180° (no convexo).
            \item Si es \textbf{cero}, los puntos son colineales.
            \item Si es \textbf{positivo}, el ángulo es menor a 180° (convexo).
        \end{itemize}
    \end{itemize}
\end{enumerate}

\subsubsection{¿Qué representa \( n \) en la complejidad?}
\begin{itemize}
    \item \( n \) es el \textbf{número total de puntos} en el conjunto \( S \).
    \item La complejidad del algoritmo es \( O(n \log n) \), debido principalmente al \textbf{ordenamiento por ángulo polar}.
    \item La construcción de la pila es \( O(n) \), ya que cada punto se empuja y se saca como máximo una vez.
\end{itemize}

\subsubsection{Ejemplo paso a paso}
Consideremos el conjunto de puntos \[ S = \{(0,0), (1,1), (2,0), (1,-1), (0.5,0.5)\}. \]
\begin{enumerate}
    \item Punto base: \( p_0 = (0,0) \).
    \item Ordenar por ángulo polar: \( (2,0) \), \( (1,-1) \), \( (0.5,0.5) \), \( (1,1) \).
    \item Construcción de la pila:
    \begin{itemize}
        \item Empujar \( (0,0) \), \( (2,0) \), \( (1,-1) \).
        \item Para \( (0.5,0.5) \): 
        Verificar convexidad entre \( (0,0) \), \( (2,0) \), \( (1,-1) \). 
        Eliminar \( (2,0) \).
        \item Empujar \( (0.5,0.5) \).
        \item Para \( (1,1) \): 
        Verificar convexidad entre \( (0,0) \), \( (0.5,0.5) \), \( (1,-1) \). 
        Empujar \( (1,1) \).
    \end{itemize}
    \item Envoltura convexa final: \( [(0,0), (0.5,0.5), (1,1), (1,-1)] \).
\end{enumerate}

%\subsubsection{Pseudocódigo}
%\begin{verbatim}
%Algoritmo Graham-Scan(S):
%    si |S| < 3 entonces
%        devolver S  // La envoltura es el conjunto mismo
%    fin si
%
%    // Paso 1: Encontrar el punto base p0
%    p0 ← punto con y mínima (y x mínima en caso de empate)
%
%    // Paso 2: Ordenar por ángulo polar respecto a p0
%    ordenar S \ {p0} por ángulo polar respecto a p0
%
%    // Paso 3: Construir la envoltura convexa
%    pila ← pila vacía
%    empujar p0 en pila
%    empujar S[1] en pila
%    empujar S[2] en pila
%
%    para i desde 3 hasta |S|-1 hacer
%        mientras el producto cruzado de los últimos tres puntos en pila < 0 hacer
%            sacar de pila
%        fin mientras
%        empujar S[i] en pila
%    fin para
%
%    devolver pila
%\end{verbatim}

%\begin{lstlisting}[style=pythonstyle, caption={Implementación del algoritmo de Graham en Python}, label={lst:graham}]
%import numpy as np
%from math import atan2

%def graham_scan(points):
%    # Ordenar los puntos por coordenada x (y luego por y)
%    points = sorted(points, key=lambda p: (p[0], p[1]))

%    # Función para calcular el ángulo polar
%    def polar_angle(p0, p1=None):
%        if p1 is None:
%            p1 = points[0]
%        y_span = p0[1] - p1[1]
%        x_span = p0[0] - p1[0]
%        return atan2(y_span, x_span)

%    # Ordenar por ángulo polar
%    points = sorted(points[1:], key=polar_angle)

%    # Inicializar la pila
%    stack = [points[0], points[1]]

%    for i in range(2, len(points)):
%        while len(stack) > 1 and cross_product(stack[-2], stack[-1], points[i]) <= 0:
%            stack.pop()
        %stack.append(points[i])

%    return stack

%def cross_product(o, a, b):
%    return (a[0] - o[0]) * (b[1] - o[1]) - (a[1] - o[1]) * (b[0] - o[0])
%\end{lstlisting}

\subsection{Algoritmo QuickHull} %para el cálculo de la envoltura convexa}
\label{subsec:quickhull}

El \textbf{algoritmo QuickHull} %\cite{Barber1996} 
es un método eficiente para calcular la envoltura convexa de un conjunto de puntos en el plano, basado en una estrategia de \textbf{divide y vencerás}. Su nombre proviene de su similitud con el algoritmo de ordenamiento \texttt{Quicksort}, ya que también utiliza una partición recursiva del conjunto de puntos para construir la envoltura convexa.

\subsubsection{Idea general}
La idea central del QuickHull es:
\begin{itemize}
    \item Identificar los puntos que forman los vértices de la envoltura convexa mediante la búsqueda de los puntos más alejados de los segmentos que unen pares de puntos.
    \item Dividir recursivamente el conjunto de puntos en subconjuntos más pequeños, cada uno asociado a un segmento de la envoltura parcial, y repetir el proceso hasta que no queden puntos externos.
\end{itemize}

\subsubsection{Etapas del algoritmo}

\begin{enumerate}
    \item \textbf{Inicialización:}
    \begin{itemize}
        \item Encontrar los puntos con las \textbf{coordenadas \( x \) mínima y máxima} en el conjunto \( S \). Estos puntos \textbf{siempre pertenecen} a la envoltura convexa.
        \item Estos dos puntos definen el \textbf{segmento inicial} de la envoltura convexa.
    \end{itemize}

    \item \textbf{División recursiva:}
    \begin{itemize}
        \item Para cada segmento de la envoltura parcial (inicialmente, el segmento entre los puntos \( x_{\min} \) y \( x_{\max} \)), encontrar el punto \( p \) más alejado del segmento.
        \item Este punto \( p \) es un nuevo vértice de la envoltura convexa.
        \item Dividir el conjunto de puntos en dos subconjuntos:
        \begin{itemize}
            \item Los puntos que están a la \textbf{izquierda} del triángulo formado por el segmento y \( p \).
            \item Los puntos que están a la \textbf{derecha} del mismo triángulo.
        \end{itemize}
        \item Repetir el proceso recursivamente para cada subconjunto.
    \end{itemize}

    \item \textbf{Construcción final:}
    \begin{itemize}
        \item Combinar los resultados de las divisiones recursivas para formar la envoltura convexa completa.
    \end{itemize}
\end{enumerate}

\subsubsection{Complejidad}
\begin{itemize}
    \item En el \textbf{caso promedio}, la complejidad del QuickHull es \( O(n \log n) \), similar al scan de Graham.
    \item En el \textbf{peor caso} (por ejemplo, cuando los puntos están en un círculo), la complejidad puede degradarse a \( O(n^2) \).
    \item Sin embargo, en la práctica, suele ser más rápido que otros algoritmos debido a su naturaleza recursiva y su capacidad para descartar grandes conjuntos de puntos en cada iteración.
\end{itemize}

\subsubsection{Ejemplo}
Consideremos el conjunto de puntos \[ S = \{(0,0), (1,1), (2,0), (1,-1), (0.5,0.5)\}. \]
\begin{enumerate}
    \item Puntos iniciales: \( x_{\min} = (0,0) \) y \( x_{\max} = (2,0) \).
    \item El punto más alejado del segmento \( (0,0)-(2,0) \) es \( (1,1) \).
    \item Dividir \( S \) en dos subconjuntos:
    \begin{itemize}
        \item Izquierda: \( \{(0,0), (0.5,0.5), (1,1)\} \).
        \item Derecha: \( \{(2,0), (1,-1), (0.5,0.5)\} \).
    \end{itemize}
    \item Repetir el proceso para cada subconjunto hasta obtener la envoltura convexa final: \[ [(0,0), (1,1), (2,0), (1,-1)]. \]
\end{enumerate}

%\subsubsection{Pseudocódigo}
%\begin{verbatim}
%Algoritmo QuickHull(S):
%    si |S| < 3 entonces
%        devolver S
%    fin si
%
%    // Paso 1: Encontrar los puntos con x mínima y máxima
%    p_min ← punto con x mínima
%    p_max ← punto con x máxima
%
%    // Conjunto de puntos a la izquierda y derecha del segmento p_min-p_max
%    izquierda ← {p ∈ S | p está a la izquierda de p_min-p_max}
%    derecha ← {p ∈ S | p está a la derecha de p_min-p_max}
%
%    // Paso 2: Llamadas recursivas
%    envoltura_izquierda ← QuickHullRecursivo(p_min, p_max, izquierda)
%    envoltura_derecha ← QuickHullRecursivo(p_max, p_min, derecha)
%
%    // Paso 3: Combinar resultados
%    devolver [p_min] + envoltura_izquierda + [p_max] + envoltura_derecha
%\end{verbatim}

\subsubsection{Ventajas y limitaciones}
\begin{itemize}
    \item \textbf{Ventajas:} Rápido en la práctica, fácil de implementar, y eficiente para conjuntos de puntos dispersos.
    \item \textbf{Limitaciones:} En el peor caso, puede ser lento (\( O(n^2) \)), y su rendimiento depende de la distribución de los puntos.
\end{itemize}

%\subsection{Algoritmo de Jarvis (1973)}

%El algoritmo de Jarvis, también conocido como el "marching ants", tiene una complejidad de \( O(nh) \), donde \( h \) es el número de puntos en la envoltura convexa.

%\begin{lstlisting}[style=pythonstyle, caption={Implementación del algoritmo de Jarvis en Python}, label={lst:jarvis}]
%def jarvis_march(points):
%    # Encontrar el punto más a la izquierda (punto inicial)
%    start = min(points, key=lambda p: (p[0], p[1]))
%
%    hull = []
%    current = start
%    while True:
%        hull.append(current)
%        next_point = points[0]
%        for p in points:
%            if p == current:
%                continue
%            # Calcular el ángulo cruzado
%            cross = (p[0] - current[0]) * (next_point[1] - current[1]) - \
%                    (p[1] - current[1]) * (next_point[0] - current[0])
%            if cross > 0 or (cross == 0 and distance(current, p) > distance(current, next_point)):
%                next_point = p
%
%        if next_point == start:
%            break
%        current = next_point
%
%    return hull
%
%def distance(p1, p2):
%    return (p1[0] - p2[0])**2 + (p1[1] - p2[1])**2
%\end{lstlisting}


%\subsection{Clustering con cáscara convexa}

%En aprendizaje automático, la cáscara convexa se utiliza en algoritmos de clustering como \textbf{DBSCAN} para identificar regiones densas y separar clusters. La idea es que los puntos dentro de una misma cáscara convexa pertenecen al mismo cluster.

%\subsection{Optimización convexa}

%En problemas de optimización, la cáscara convexa de un conjunto de restricciones define la región factible. La optimalidad de un punto \( x^* \) puede caracterizarse utilizando el teorema de separación de hiperplanos y las propiedades de la cáscara convexa.

\subsection{Aplicaciones}

%\subsection{Envoltura convexa de un conjunto de puntos}

La \textbf{envoltura convexa} (o cáscara convexa) %de un conjunto finito de puntos en el plano es el polígono convexo más pequeño que contiene a todos los puntos. Este concepto 
es esencial en:

\begin{itemize}
    \item \textbf{Geometría computacional}: Para problemas de colisión, localización de puntos y clustering.
    \item \textbf{Visión por computadora}: En reconstrucción 3D y reconocimiento de formas.
    \item \textbf{Robótica}: Para planificación de caminos y evitación de obstáculos.
    \item \textbf{Aprendizaje automático}: En clustering (por ejemplo, el algoritmo DBSCAN).
\end{itemize}

% Separacion

\section{Separación de convexos}

En el contexto de la separación de conjuntos convexos, el resultado fundamental establece que, si dos conjuntos convexos son disjuntos, siempre existe un hiperplano que los separa. Este hiperplano actúa como una barrera en el espacio que permite distinguir claramente entre los elementos de un conjunto y los del otro, sin solapamientos. 
%\begin{theorem}
%\dots
%\end{theorem}

% Carateodory -- Minkowski ETC Farkas

\section{Para profundizar}

Recomendamos al lector interesado en profundizar en el tema consultar las siguientes obras: 
\cite{BoydVandenberghe2004}, \cite{Grünbaum1963}, 
\cite{hiriart2001},  
\cite{Rockafellar1970}.

\chapter{Funciones convexas en dimensión finita}

Este capítulo aborda el análisis de funciones convexas en dimensión finita, un concepto central en optimización, economía matemática y aprendizaje automático. Se estudian dos contextos principales: el caso general (no necesariamente continuo) y el caso continuo. Para cada contexto, se presentan definiciones, propiedades fundamentales, caracterizaciones de mínimos locales y globales, y herramientas como el subgradiente y la Hessiana.

\section{Caso general}

Antes de abordar las propiedades específicas de las funciones convexas continuas, comenzaremos analizando el caso más general: funciones convexas que, aunque definidas sobre conjuntos convexos, no necesariamente cumplen con la propiedad de continuidad.

\subsection{
%Funciones convexas en dimensión finita: 
Definición, ejemplos y propiedades}
\label{subsec:funcion_convexa}

\begin{epigraph}{\small
``La convexidad es el lenguaje natural de la optimización y el análisis de funciones en espacios vectoriales.''}
%\source
{Adaptado de \textit{Convex Analysis} (Rockafellar, 1970)}
\end{epigraph}

Las funciones convexas son un pilar fundamental en el análisis matemático, la optimización y la inteligencia artificial. Su definición, aunque aparentemente simple, encierra profundas implicaciones geométricas y algebraicas, especialmente en espacios de dimensión finita. A continuación, exploramos esta noción desde sus bases hasta ejemplos concretos, destacando el papel del dominio extendido y su relación con la convexidad de conjuntos.

\begin{definition}[Función convexa en dimensión finita]
Una función \[f: \mathbb{R}^n \to (-\infty, +\infty] \] es \textbf{convexa} si su \textbf{dominio extendido} \( \text{dom}(f) = \{x \in \mathbb{R}^n \mid f(x) < +\infty\} \) es un conjunto convexo y, para todo \( x, y \in \text{dom}(f) \) y todo \( \lambda \in [0,1] \), se cumple:
\[
f(\lambda x + (1-\lambda) y) \leq \lambda f(x) + (1-\lambda) f(y).
\]
\end{definition}

%\subsubsection{Explicación del dominio extendido}
El \textbf{dominio extendido} \( \text{dom}(f) \) incluye todos los puntos donde \( f \) no toma el valor \( +\infty \). Este concepto es crucial porque:
\begin{itemize}
    \item Permite manejar funciones definidas solo en subconjuntos específicos de \( \mathbb{R}^n \) (por ejemplo, restricciones en optimización).
    \item Garantiza que la desigualdad de convexidad se aplique \textbf{solo} a pares de puntos donde \( f \) está bien definida.
\end{itemize}

\subsubsection{Ejemplos ilustrativos}

\begin{ex}[Función valor absoluto: \( f(x) = |x| \)] Tiene las c\'aracteristicas siguientes.
\begin{itemize}
    \item \textbf{Dominio extendido:} \( \text{dom}(f) = \mathbb{R} \).
    \item \textbf{Convexidad:} Cumple la desigualdad de convexidad (por desigualdad triangular):
    \[
    |\lambda x + (1-\lambda)y| \leq \lambda |x| + (1-\lambda)|y|.
    \]
    \item \textbf{Visualización:}
\end{itemize}

\begin{center}
\begin{tikzpicture}[scale=0.99, domain=-2:2]
    \draw[-] (-3,0) -- (3,0) node[right] {$x$};
    \draw[-] (0,-0.5) -- (0,2.5) node[above] {$f(x)$};
    \draw[blue, thick] plot (\x, {abs(\x)});
    \node[below] at (2,0) {$f(x) = |x|$};
\end{tikzpicture}
\end{center}
\end{ex}

\begin{ex}[Función parte positiva: \( f(x) = \max(x, 0) \)] Tiene las c\'aracteristicas siguientes.
\begin{itemize}
    \item \textbf{Dominio extendido:} \( \text{dom}(f) = \mathbb{R} \).
    \item \textbf{Convexidad:} %Es lineal por tramos, por lo que es convexa.
    se verifica directamente con la d\'efinici\'on.
    \item \textbf{Visualización:}
\end{itemize}

\begin{center}
\begin{tikzpicture}[scale=0.99, domain=-2:2]
    \draw[-] (-3,0) -- (3,0) node[right] {$x$};
    \draw[-] (0,-0.5) -- (0,2.5) node[above] {$f(x)$};
    \draw[blue, thick] plot (\x, {max(\x,0)});
    \node[below] at (2,0) {$f(x) = \max(x, 0)$};
\end{tikzpicture}
\end{center}
\end{ex}

\begin{ex}[Función indicatriz de un conjunto convexo \( C \)]
% \textbf{Función indicatriz de un conjunto convexo \( C \):}
Se define as\'i:
    \[
    f(x) = \begin{cases}
    0 & \text{si } x \in C, \\
    +\infty & \text{si } x \notin C.
    \end{cases}
    \]
    \begin{itemize}
        \item \textbf{Dominio extendido:} \( \text{dom}(f) = C \), que es convexo por definición.
        \item \textbf{Convexidad:} La desigualdad se cumple trivialmente (ambos lados son \( +\infty \) o \( 0 \)).
    \end{itemize}
\end{ex}

\begin{ex}
[Función norma euclidiana]  

  \textbf{Función norma euclidiana:} Se define por \( f(x) = \|x\|_2 \).
    \begin{itemize}
        \item \textbf{Dominio extendido:} \( \text{dom}(f) = \mathbb{R}^n \).
        \item \textbf{Convexidad:} Desigualdad triangular:
        \[
        \|\lambda x + (1-\lambda)y\|_2 \leq \lambda \|x\|_2 + (1-\lambda)\|y\|_2.
        \]
    \end{itemize}
\end{ex}

%\subsection{Función convexa propia: definición y propiedades}
%\label{subsec:funcion_convexa_propia}

Las funciones convexas propias son un caso especialmente importante en optimización y análisis convexo, ya que combinan la convexidad con condiciones naturales sobre su dominio y valores. A continuación, presentamos su definición formal. % y sus implicaciones clave.

\begin{definition}[Función convexa propia]
Una función \( f: \mathbb{R}^n \to (-\infty, +\infty] \) se denomina \textbf{convexa propia} si cumple las siguientes condiciones:
\begin{enumerate}
    \item \textbf{Convexidad:} \( f \) es convexa en el sentido usual:
    \[
    f(\lambda x + (1-\lambda) y) \leq \lambda f(x) + (1-\lambda) f(y), \quad \forall x, y \in \mathbb{R}^n, \forall \lambda \in [0,1].
    \]
    \item \textbf{Condición de no trivialidad:} \( f(x) < +\infty \) para al menos un \( x \in \mathbb{R}^n \) (es decir, \( \text{dom}(f) \neq \emptyset \)).
    %\item %\textbf{Condición de no infinitud global:} \( f(x) > -\infty \) para todo \( x \in \mathbb{R}^n \) (no existen puntos donde \( f \) tome el valor \( -\infty \)).
\end{enumerate}
\end{definition}

%\subsubsection{Interpretación}
%\begin{itemize}
%    \item La condición de convexidad asegura la "rigidez" geométrica necesaria para aplicaciones en optimización.
%    \item La condición de no trivialidad garantiza que la función tenga un interés práctico (no es idénticamente \( +\infty \)).
%    \item La condición de no infinitud global evita comportamientos patológicos como \( f(x) = -\infty \) para todo \( x \).
%\end{itemize}

%\subsubsection{Ejemplos}
%\begin{itemize}
%    \item \textbf{Función propia:} \( f(x) = x^2 \) (convexa, \( \text{dom}(f) = \mathbb{R} \), y \( f(x) > -\infty \)).
%    \item \textbf{No propia:} \( f(x) = +\infty \) para todo \( x \) (trivial y sin interés).
%    \item \textbf{No propia:} \( f(x) = -\infty \) para algún \( x \) (prohibido por definición).
%\end{itemize}

El epígrafe de una función es, en términos sencillos, el 'volumen' o la región en el espacio que se forma por todos los puntos que están por encima de la gráfica de la función.

\begin{definition}[Epígrafe de una función]
El \textbf{epígrafe} de \( f \) se define como:
\[
\text{epi}(f) = \{(x, \alpha) \in \mathbb{R}^n \times \mathbb{R} \mid f(x) \leq \alpha\}.
\]
\end{definition}


\paragraph{Propiedades clave de funciones convexas}
\begin{itemize}
    \item \textbf{Cierre bajo sumas ponderadas:} Si \( f \) y \( g \) son convexas, entonces \( \alpha f + \beta g \) también lo es para \( \alpha, \beta \geq 0 \).
    \item \textbf{Subniveles convexos:} El conjunto \( \{x \mid f(x) \leq \alpha\} \) es convexo para todo \( \alpha \in \mathbb{R} \).
    \item \textbf{Relación con la convexidad de conjuntos:} Una función es convexa si y solo si su \textbf{epígrafe} (conjunto de puntos por encima de su gráfica) es un conjunto convexo en \( \mathbb{R}^{n+1} \).
\end{itemize}

\subsubsection{Transición a propiedades avanzadas}
La convexidad no solo simplifica el análisis de funciones, sino que también garantiza la existencia de mínimos globales en problemas de optimización. En lo siguiente, exploraremos cómo estas propiedades se aplican en algoritmos de optimización y aprendizaje automático, donde la convexidad es una condición suficiente para la eficiencia computacional.

%\textbf{Nota:} Para una implementación práctica, recomendamos consultar librerías como \texttt{CVXPY} o \texttt{Convex.jl}, que permiten modelar y resolver problemas de optimización convexa de manera eficiente.

\subsection{Propiedades de continuidad} % de funciones convexas en dimensión finita}
\label{subsec:continuidad_funciones_convexas}

Las funciones convexas en dimensión finita poseen propiedades de continuidad que las distinguen de otras clases de funciones y son fundamentales en optimización y análisis convexo. A continuación, se presentan los resultados teóricos principales, sus demostraciones y ejemplos ilustrativos.

\begin{theorem}[Continuidad en el interior del dominio]
\label{thm:continuidad_interior}
Sea \[ f: \mathbb{R}^n \to (-\infty, +\infty] \] una función convexa propia. Entonces, \( f \) es \textbf{localmente lipschitziana} y por lo tanto \textbf{continua} en el interior de su dominio extendido \( \text{int}(\text{dom}(f)) \).
\end{theorem}

\begin{proof}
%La convexidad implica que, para todo \( x \in \text{int}(\text{dom}(f)) \), existe un entorno \( B(x, r) \) donde \( f \) está acotada superiormente. Por el teorema de la acotación local de funciones convexas \cite{Rockafellar1970}, \( f \) es localmente lipschitziana en \( \text{int}(\text{dom}(f)) \), lo que garantiza su continuidad. 
La prueba  puede consultarse en \cite[Sección 2.1.2]{BoydVandenberghe2004}. %, Sección 2.1.2.
\end{proof}

\begin{theorem}[Continuidad global en \( \mathbb{R}^n \)]
\label{thm:continuidad_globa}
Si \( f: \mathbb{R}^n \to \mathbb{R} \) es una función convexa y \( \text{dom}(f) = \mathbb{R}^n \), entonces \( f \) es \textbf{globalmente continua} en \( \mathbb{R}^n \).
\end{theorem}

\begin{proof}
%Como \( f \) es convexa y definida en todo \( \mathbb{R}^n \), por el Teorema \ref{thm:continuidad_interior}, \( f \) es localmente lipschitziana en todo punto. La compacidad de los conjuntos cerrados y acotados en \( \mathbb{R}^n \) garantiza que \( f \) es globalmente continua. Para una demostración formal, 
Véase \cite[Teorema 10.1]{Rockafellar1970}.
\end{proof}

\subsubsection{Ejemplos, contraejemplos y casos límite}

\begin{ex}[Función convexa discontinua en la frontera]
Consideremos la función indicatriz de un conjunto convexo \( C \):
\[
f(x) = \begin{cases}
0 & \text{si } x \in C, \\
+\infty & \text{si } x \notin C.
\end{cases}
\]
\begin{itemize}
    \item \( f \) es convexa propia si \( C \) es no vacío y convexo.
    \item \( f \) es discontinua en todo punto \( x \in \partial C \) (frontera de \( C \)), ya que \( \lim_{y \to x, y \notin C} f(y) = +\infty \).
    \item Este ejemplo muestra que la continuidad solo se garantiza en el interior de \( \text{dom}(f) \).
\end{itemize}
\end{ex}

\begin{ex}[Función convexa definida en \( \mathbb{R}^n \)]
La función \( f(x) = \|x\|_2^2 \) es convexa y definida en \( \mathbb{R}^n \). Por el Teorema \ref{thm:continuidad_globa}, \( f \) es continua en todo \( \mathbb{R}^n \).
\end{ex}

\begin{ex}[No convexidad de la función de Heaviside]
\label{ex:heaviside_non_convexe}
%
Consideramos la función de Heaviside, definida como:
\[
H(x) = \begin{cases}
0 & \text{si } x < 0, \\
c & \text{si } x = 0, \quad \text{donde } c \in \{0, 1\}, \\
1 & \text{si } x > 0.
\end{cases}
\]
Demostramos que \( H \) no es convexa, independientemente del valor de \( c \).
\textbf{Caso 1 : \( H(0) = 0 \)}
Tomemos \( x_1 = -1 \), \( x_2 = 2 \), y \( \lambda = 0.5 \):
\[
H(0.5 \cdot (-1) + 0.5 \cdot 2) = H(0.5) = 1,
\]
y
\[
0.5 \cdot H(-1) + 0.5 \cdot H(2) = 0.5 \cdot 0 + 0.5 \cdot 1 = 0.5.
\]
Aquí, \( 1 \not\leq 0.5 \): la desigualdad de convexidad no se cumple.
Por lo tanto, \( H \) no es convexa incluso si \( H(0) = 0 \).

\textbf{Caso 2 : \( H(0) = 1 \)}
Tomemos:
- \( x_1 = -1 \) (donde \( H(-1) = 0 \)),
- \( x_2 = 1 \) (donde \( H(1) = 1 \)),
- \( \lambda = 0.5 \).
Calculamos:
\[
H(\lambda x_1 + (1-\lambda) x_2) = H(0) = 1,
\]
y
\[
0.5 \cdot H(-1) + 0.5 \cdot H(1) = 0.5 \cdot 0 + 0.5 \cdot 1 = 0.5.
\]
La desigualdad de convexidad exige:
\[
1 \leq 0.5,
\]
lo cual es falso. Por lo tanto, \( H \) no es convexa si \( H(0) = 1 \).
Independientemente del valor de \( H(0) \) (\( 0 \) o \( 1 \)), la función de Heaviside no satisface la desigualdad de convexidad para ciertos puntos. Por lo tanto, nunca es convexa.
\end{ex}

\subsubsection{Implicaciones en optimización}

\begin{itemize}
    \item \textbf{Existencia de mínimos:} Si \( f \) es convexa propia y \( C \subset \mathbb{R}^n \) es un conjunto convexo cerrado no vacío, entonces \( f \) alcanza su mínimo sobre \( C \) si \( C \cap \text{dom}(f) \neq \emptyset \) \cite{Rockafellar1970}.
    \item \textbf{Estabilidad numérica:} La continuidad de \( f \) en el interior de su dominio garantiza la convergencia de algoritmos como el de gradiente descendente.
    \item \textbf{Subgradientes:} En puntos donde \( f \) no es diferenciable (por ejemplo, \( f(x) = |x| \) en \( x = 0 \)), la convexidad permite definir subgradientes, que son herramientas clave en optimización.
\end{itemize}

%\subsubsection{Resumen de propiedades}

%\begin{table}[h!]
%\centering
%\begin{tabular}{|l|l|l|}
%\hline
%\textbf{Propiedad} & \textbf{Caso} & \textbf{Consecuencia} \\ \hline
%Continuidad en interior & \( \text{int}(\text{dom}(f)) \neq \emptyset \) & \( f \) es localmente lipschitziana \\ \hline
%Continuidad global & \( \text{dom}(f) = \mathbb{R}^n \) & \( f \) es continua en \( \mathbb{R}^n \) \\ \hline
%Discontinuidad en frontera & \( f(x) = +\infty \) en \( \partial \text{dom}(f) \) & \( f \) puede ser discontinua en \( \partial \text{dom}(f) \) \\ \hline
%\end{tabular}
%\caption{Resumen de propiedades de continuidad para funciones convexas en dimensión finita.}
%\label{tab:continuidad_resumen}
%\end{table}
%
%---
%
%\textbf{Nota:} Para un estudio más detallado de las pruebas y aplicaciones, se recomienda consultar:
%\begin{itemize}
%    \item \cite{Rockafellar1970}: Fundamentos teóricos y demostraciones rigurosas.
%    \item \cite{BoydVandenberghe2004}: Enfoque práctico y ejemplos en optimización convexa.
%    \item \cite{Hiriart-Urruty2001}: Análisis avanzado de funciones convexas y sus propiedades.
%\end{itemize}

\subsection{Subdiferencial de funciones convexas}

La noción de subdiferencial generaliza el concepto de gradiente a funciones convexas que no son necesariamente diferenciables. Esta herramienta es fundamental en optimización convexa, ya que permite caracterizar óptimos incluso en presencia de no diferenciabilidad.

\begin{definition}[Subgradiente]
Sea \( f: \mathbb{R}^n \to (-\infty, +\infty] \) una función convexa. Un vector \( g \in \mathbb{R}^n \) es un \textbf{subgradiente} de \( f \) en \( x \) si
\[
f(y) \geq f(x) + g^T (y - x) \quad \forall y \in \mathbb{R}^n.
\]
El conjunto de todos los subgradientes de \( f \) en \( x \) se denota \( \partial f(x) \) y se llama \textbf{subdiferencial} de \( f \) en \( x \).
\end{definition}

Intuitivamente, un subgradiente define un hiperplano de soporte al epígrafo de la función en el punto considerado. Cuando la función es diferenciable, el subdiferencial se reduce a un singleton que contiene el gradiente clásico.

A continuación, resumimos algunas propiedades fundamentales del subdiferencial en dimensión finita.

\paragraph{Propiedades clave.}
Sea \( f: \mathbb{R}^n \to (-\infty,+\infty] \) convexa y \( x \in \mathrm{dom}(f) \). Entonces:
\begin{itemize}
    \item El conjunto \( \partial f(x) \) es convexo y cerrado.
    \item Si \( f \) es diferenciable en \( x \), entonces
    \[
    \partial f(x) = \{ \nabla f(x) \}.
    \]
    \item La condición de optimalidad de primer orden se escribe:
    \[
    0 \in \partial f(x) \quad \Longleftrightarrow \quad x \text{ es un minimizador global de } f.
    \]
    \item La aplicación \( x \mapsto \partial f(x) \) es monótona en el sentido de operadores:
    \[
    (g_1 - g_2)^T (x_1 - x_2) \geq 0
    \quad \text{para todo } g_i \in \partial f(x_i).
    \]
\end{itemize}

Para ilustrar esta noción, presentamos ahora algunos ejemplos clásicos de cálculo de subdiferenciales.

\paragraph{Ejemplos.}

\begin{itemize}
    \item \textbf{Valor absoluto.} Sea \( f(x) = |x| \) con \( x \in \mathbb{R} \). Entonces:
    \[
    \partial f(x) =
    \begin{cases}
        \{1\} & \text{si } x > 0, \\
        \{-1\} & \text{si } x < 0, \\
        [-1,1] & \text{si } x = 0.
    \end{cases}
    \]

    \item \textbf{Parte positiva.} Sea \( f(x) = \max(0,x) \). Entonces:
    \[
    \partial f(x) =
    \begin{cases}
        \{1\} & \text{si } x > 0, \\
        \{0\} & \text{si } x < 0, \\
        [0,1] & \text{si } x = 0.
    \end{cases}
    \]

    \item \textbf{Función indicatriz.} Sea \( C \subset \mathbb{R}^n \) un conjunto convexo cerrado y
    \[
    \iota_C(x) =
    \begin{cases}
        0 & \text{si } x \in C, \\
        +\infty & \text{si } x \notin C.
    \end{cases}
    \]
    Entonces, para \( x \in C \), el subdiferencial está dado por el cono normal:
    \[
    \partial \iota_C(x) = N_C(x) = \{ g \in \mathbb{R}^n \mid g^T (y - x) \leq 0, \ \forall y \in C \}.
    \]
\end{itemize}

Estos ejemplos muestran cómo el subdiferencial captura la falta de diferenciabilidad mediante conjuntos de pendientes admisibles, lo cual resulta esencial en el análisis y la resolución de problemas de optimización convexa no suave.

\subsection{Caracterización de mínimos} % mediante el subdiferencial}

Una de las principales ventajas del subdiferencial es que permite formular condiciones de optimalidad de primer orden sin requerir diferenciabilidad. En el contexto de funciones convexas, esta caracterización es particularmente simple y poderosa.
%
El siguiente resultado establece una condición necesaria y suficiente de optimalidad global.

\begin{theorem}[Caracterización de mínimos]
Sea \( f: \mathbb{R}^n \to (-\infty, +\infty] \) una función convexa. Un punto \( x^* \in \mathrm{dom}(f) \) es un mínimo global de \( f \) si y solo si
\[
0 \in \partial f(x^*).
\]
\end{theorem}

Este resultado generaliza la condición clásica \( \nabla f(x^*) = 0 \) del caso diferenciable. En efecto, cuando \( f \) es diferenciable en \( x^* \), se tiene \( \partial f(x^*) = \{\nabla f(x^*)\} \), y la condición anterior se reduce a la conocida anulación del gradiente.

\paragraph{Idea de la demostración.}
La demostración se basa directamente en la definición de subgradiente.
\begin{itemize}
    \item (\emph{Suficiencia}) Si \( 0 \in \partial f(x^*) \), entonces por definición:
    \[
    f(y) \geq f(x^*) \quad \forall y \in \mathbb{R}^n,
    \]
    lo que implica que \( x^* \) es un mínimo global.
    \item (\emph{Necesidad}) Si \( x^* \) es un mínimo global, entonces
    \[
    f(y) \geq f(x^*) \quad \forall y,
    \]
    lo que corresponde exactamente a la desigualdad de subgradiente con \( g = 0 \).
\end{itemize}

Para una demostración más detallada y en un marco más general (espacios vectoriales normados), se puede consultar textos clásicos de optimización convexa como \cite{hiriart2001}. %Rockafellar o Hiriart-Urruty y Lemaréchal.

\paragraph{Interpretación geométrica.}
La condición \( 0 \in \partial f(x^*) \) significa que existe un hiperplano de soporte horizontal al epígrafo de \( f \) en \( x^* \). Intuitivamente, esto indica que no existe ninguna dirección de descenso, lo que caracteriza un mínimo global en el caso convexo.

A continuación, ilustramos esta caracterización con ejemplos clásicos.

\paragraph{Ejemplos.}

\begin{itemize}
    \item \textbf{Valor absoluto.} Sea \( f(x) = |x| \). Entonces:
    \[
    \partial f(0) = [-1,1].
    \]
    Como \( 0 \in [-1,1] \), se concluye que \( x^* = 0 \) es un mínimo global. En cambio, si \( x \neq 0 \), entonces \( \partial f(x) \) no contiene al cero, lo que confirma que no son minimizadores.

    \item \textbf{Parte positiva.} Sea \( f(x) = \max(0,x) \). Se tiene:
    \[
    \partial f(0) = [0,1].
    \]
    Dado que \( 0 \in \partial f(0) \), el punto \( x^* = 0 \) es un mínimo global (de hecho, \( f(x) \geq 0 \) para todo \( x \)).

    \item \textbf{Función indicatriz.} Sea \( C \subset \mathbb{R}^n \) convexo cerrado y \( f = \iota_C \). Entonces:
    \[
    \partial f(x) = N_C(x).
    \]
    La condición de optimalidad se escribe:
    \[
    0 \in N_C(x^*),
    \]
    lo cual equivale a decir que
    \[
    0^T (y - x^*) \leq 0 \quad \forall y \in C,
    \]
    condición que es siempre verdadera. Esto refleja el hecho de que todo punto de \( C \) es minimizador de \( \iota_C \).
\end{itemize}

\paragraph{Comentario.}
Este criterio es fundamental en optimización convexa y sirve como base para numerosos algoritmos (métodos de subgradiente, métodos proximales, etc.). Su simplicidad contrasta con su gran generalidad, ya que sigue siendo válido incluso en ausencia de diferenciabilidad.

%\section{Teorema de Existencia de Mínimo para Funciones Convexas en Dimensión Finita}

\section{Existencia de mínimo}

%\subsection{Teorema}

Consideremos una función $f: \mathbb{R}^n \to \mathbb{R} \cup \{+\infty\}$ convexa, propia (es decir, $f(x) < +\infty$ para al menos un $x \in \mathbb{R}^n$) y semicontinua inferiormente. Si el dominio efectivo de $f$, definido como $\text{dom}(f) = \{x \in \mathbb{R}^n : f(x) < +\infty\}$, es no vacío, cerrado y acotado, entonces $f$ alcanza su mínimo global en $\mathbb{R}^n$.

\begin{theorem}
    Sea $f: \mathbb{R}^n \to \mathbb{R} \cup \{+\infty\}$ una función convexa, propia y semicontinua inferiormente. Si $\text{dom}(f)$ es no vacío, cerrado y acotado, entonces existe $x^* \in \text{dom}(f)$ tal que
    \begin{equation*}
        f(x^*) = \inf_{x \in \mathbb{R}^n} f(x).
    \end{equation*}
\end{theorem}

%\subsection{Demostración}

%La demostración de este teorema se basa en el teorema de Weierstrass, que establece que una función continua en un conjunto compacto alcanza su mínimo y máximo. Dado que $f$ es semicontinua inferiormente y $\text{dom}(f)$ es cerrado y acotado (por lo tanto, compacto), se sigue que $f$ alcanza su mínimo en $\text{dom}(f)$. La convexidad de $f$ garantiza la unicidad del mínimo si $f$ es estrictamente convexa.

%\subsection{Ejemplo}

Consideremos la función $f(x) = x^2$, definida en $\mathbb{R}$. Esta función es convexa, propia y semicontinua inferiormente. Su dominio efectivo es $\mathbb{R}$, que no es acotado. Sin embargo, si restringimos el dominio a un intervalo cerrado y acotado, por ejemplo, $[-1, 1]$, entonces $f$ alcanza su mínimo en $x = 0$.

\section{Caso diferenciable}

\begin{definition}[Función convexa diferenciable]
Si \( f: \mathbb{R}^n \to \mathbb{R} \) es convexa y diferenciable en \( x \), entonces
\[
f(y) \geq f(x) + \nabla f(x)^T (y - x) \quad \forall y \in \mathbb{R}^n.
\]
Además, \( \nabla f(x) \) es el único subgradiente de \( f \) en \( x \).
\end{definition}

\begin{definition}[Función convexa dos veces diferenciable]
Si \( f: \mathbb{R}^n \to \mathbb{R} \) es convexa y dos veces diferenciable, entonces su matriz Hessiana \( \nabla^2 f(x) \) es \textbf{semidefinida positiva} para todo \( x \in \mathbb{R}^n \). Recíprocamente, si \( \nabla^2 f(x) \) es semidefinida positiva para todo \( x \), entonces \( f \) es convexa.
\end{definition}

\begin{theorem}[Caracterización de mínimos con gradiente]
Sea \( f: \mathbb{R}^n \to \mathbb{R} \) una función convexa y diferenciable. Un punto \( x^* \in \mathbb{R}^n \) es un mínimo global de \( f \) si y solo si
\[
\nabla f(x^*) = 0.
\]
\end{theorem}

\begin{theorem}[Caracterización de mínimos con Hessiana]
Sea \( f: \mathbb{R}^n \to \mathbb{R} \) una función convexa y dos veces diferenciable. Un punto \( x^* \in \mathbb{R}^n \) es un mínimo global estricto de \( f \) si y solo si
\[
\nabla f(x^*) = 0 \quad \text{y} \quad \nabla^2 f(x^*) \text{ es definida positiva.}
\]
\end{theorem}

\chapter{M\'as métodos}

\section{Método de penalización}
El \textbf{método de penalización} transforma un problema con restricciones \( \min f(x) \) sujeto a \( c_i(x) = 0 \) en un problema sin restricciones:
\[
\min f(x) + \mu \sum_{i=1}^m c_i(x)^2,
\]
donde \( \mu > 0 \) es el parámetro de penalización.

\section{Método de Lagrange}
El \textbf{método de Lagrange} introduce multiplicadores \( \lambda \) para problemas con restricciones \( c_i(x) = 0 \). La función Lagrangiana es:
\[
\mathcal{L}(x, \lambda) = f(x) + \sum_{i=1}^m \lambda_i c_i(x).
\]
Las condiciones de optimalidad son:
\[
\nabla_x \mathcal{L}(x^*, \lambda^*) = 0, \quad c_i(x^*) = 0.
\]

%\begin{quotation}
%\textbf{Vínculo con ecuaciones diferenciales:}
%La dinámica de los multiplicadores de Lagrange puede describirse mediante:
%\[
%\frac{dx}{dt} = -\nabla_x \mathcal{L}(x, \lambda), \quad \frac{d\lambda}{dt} = c(x).
%\]
%Esta ecuación diferencial parcial (EDP) describe la evolución de \( x \) y \( \lambda \) hacia un punto de equilibrio.
%\end{quotation}

\section{Método de Lagrange aumentado} % en Dimensión Finita}

%Se puede consultar Burman et al., 2023.

\subsection{Definición}

El método de Lagrange aumentado, también conocido como método de los multiplicadores de Lagrange con penalización, es una técnica utilizada para resolver problemas de optimización con restricciones. Este método es una generalización del método clásico de los multiplicadores de Lagrange, incorporando una función de penalización que permite manejar restricciones de igualdad de manera más flexible.

\textbf{Problema general:}
Dado un problema de optimización no lineal con restricciones de igualdad:
\begin{equation*}
    \begin{aligned}
        & \min_{x \in \mathbb{R}^n} & & f(x) \\
        & \text{sujeto a} & & h_i(x) = 0, \quad i = 1, \ldots, m,
    \end{aligned}
\end{equation*}
\noindent donde $f: \mathbb{R}^n \to \mathbb{R}$ es la función objetivo y $h_i: \mathbb{R}^n \to \mathbb{R}$ son las funciones de restricción, el método de Lagrange aumentado introduce una función de penalización cuadrática en las restricciones, transformando el problema en uno sin restricciones de la siguiente manera:
\begin{equation*}
    \mathcal{L}_\mu(x, \lambda) = f(x) + \lambda^T h(x) + \frac{\mu}{2} \|h(x)\|^2,
\end{equation*}
\noindent donde $\lambda \in \mathbb{R}^m$ son los multiplicadores de Lagrange asociados a las restricciones, y $\mu > 0$ es un parámetro de penalización.

\subsection{Algoritmo}

El algoritmo básico del método de Lagrange aumentado (Uzawa) es el siguiente:

\begin{enumerate}
    \item Inicializar $\lambda_0$ y $\mu > 0$.
    \item Resolver el problema sin restricciones:
    \begin{equation*}
        x_{k+1} = \arg\min_{x} \mathcal{L}_\mu(x, \lambda_k).
    \end{equation*}
    \item Actualizar los multiplicadores:
    \begin{equation*}
        \lambda_{k+1} = \lambda_k + \mu h(x_{k+1}).
    \end{equation*}
    \item Verificar convergencia. Si no se cumple, repetir desde el paso 2.
\end{enumerate}

\subsection{Ejemplo}

Consideremos el siguiente problema de optimización con restricción de igualdad:
\begin{equation*}
    \begin{aligned}
        & \min_{x,y} & & x^2 + y^2 \\
        & \text{sujeto a} & & x + y = 1.
    \end{aligned}
\end{equation*}
Aplicando el método de Lagrange aumentado, definimos la función de Lagrange aumentada:
\begin{equation*}
    \mathcal{L}_\mu(x, y, \lambda) = x^2 + y^2 + \lambda (x + y - 1) + \frac{\mu}{2} (x + y - 1)^2.
\end{equation*}
Para resolver el problema, minimizamos $\mathcal{L}_\mu$ respecto a $x$ e $y$, y luego actualizamos $\lambda$ iterativamente hasta alcanzar la convergencia.
\noindent Por ejemplo, si tomamos $\mu = 1$ y $\lambda_0 = 0$, el punto óptimo se alcanza en $(x, y) = (0.5, 0.5)$.

\section{Métodos cuasi-Newton (BFGS)}
El método \textbf{BFGS} aproxima la inversa de la Hessiana \( H_k \) mediante actualizaciones iterativas:
\[
H_{k+1} = \left( I - \frac{y_k s_k^T}{y_k^T s_k} \right) H_k \left( I - \frac{s_k y_k^T}{y_k^T s_k} \right) + \frac{s_k s_k^T}{y_k^T s_k},
\]
donde \( s_k = x_{k+1} - x_k \) y \( y_k = \nabla f(x_{k+1}) - \nabla f(x_k) \). La iteración es:
\[
x_{k+1} = x_k - \alpha_k H_k \nabla f(x_k).
\]

\chapter{Medida y análisis funcional}

Este capítulo está dedicado a profundizar en los fundamentos matemáticos esenciales para el análisis avanzado: la \textbf{teoría de medida}, la \textbf{integral de Lebesgue} y los \textbf{espacios funcionales}. Estos conceptos son pilares en el estudio de la probabilidad, la ecuación diferencial parcial, el análisis funcional y la teoría de control. A través de un enfoque riguroso, se presentan las herramientas teóricas necesarias para abordar problemas complejos en matemática, sentando las bases para el desarrollo de métodos modernos en optimización, física matemática y ciencia de datos.

\section{Medida}
La noción de \textbf{medida} generaliza el concepto clásico de longitud, área y volumen a conjuntos abstractos, permitiendo cuantificar el "tamaño" de subconjuntos en espacios arbitrarios. Esta teoría es fundamental para definir integrales generalizadas, como la integral de Lebesgue, y para analizar funciones medibles. Se exploran las propiedades de los espacios de medida, la construcción de medidas a partir de premedidas, y la relación entre la medida y la topología. Además, se estudian los teoremas clásicos de convergencia, como el de la convergencia monótona y el de la convergencia dominada, que son esenciales en el análisis funcional y en la teoría de la probabilidad.

%\section{Propiedades Finas de Funciones en Teoría de la Medida}

\section{Medida y funciones}

\subsection{Funciones BV} %de Variación Acotada (BV)}

\begin{definition}[Variación acotada]
Una función \(f: [a,b] \to \mathbb{R}\) es de \textbf{variación acotada} si su variación total es finita:
\[
V_a^b(f) = \sup \left\{ \sum_{i=1}^n |f(x_i) - f(x_{i-1})| : a = x_0 < x_1 < \cdots < x_n = b \right\} < \infty.
\]
El espacio \(BV([a,b])\) es el conjunto de funciones de variación acotada en \([a,b]\).
\end{definition}

\begin{theorem}[Estructura de funciones BV]
Toda función \(f \in BV([a,b])\) puede escribirse como la diferencia de dos funciones crecientes. Además, \(f\) es derivable en casi todo punto y su derivada es integrable.
\end{theorem}

\subsection{Continuidad absoluta}

\begin{definition}[Continuidad absoluta]
Una función \(f: [a,b] \to \mathbb{R}\) es \textbf{absolutamente continua} si para todo \(\varepsilon > 0\) existe \(\delta > 0\) tal que para cualquier colección finita de intervalos disjuntos \((a_i, b_i) \subset [a,b]\) con \(\sum_{i} (b_i - a_i) < \delta\), se tiene:
\[
\sum_{i} |f(b_i) - f(a_i)| < \varepsilon.
\]
\end{definition}

\begin{theorem}[Caracterización de funciones absolutamente continuas]
Sea \(f: [a,b] \to \mathbb{R}\) una función. Las siguientes afirmaciones son equivalentes:
\begin{enumerate}
    \item \(f\) es absolutamente continua.
    \item \(f\) es la integral indefinida de su derivada: existe \(g \in L^1([a,b])\) tal que
    \[
    f(x) = f(a) + \int_a^x g(t) \, dt \quad \forall x \in [a,b].
    \]
    \item \(f\) es continua, de variación acotada, y la medida de Lebesgue-Stieltjes asociada a \(f\) es absolutamente continua respecto a la medida de Lebesgue.
\end{enumerate}
\end{theorem}

\begin{remark}
Las funciones absolutamente continuas son un subconjunto propio de las funciones continuas de variación acotada. Además, toda función absolutamente continua es derivable en casi todo punto y su derivada coincide con la función \(g\) del teorema anterior.
\end{remark}

\begin{ex}
La función \(f(x) = x^2 \sin(1/x)\) para \(x \neq 0\) y \(f(0) = 0\) es absolutamente continua en \([0,1]\), pero no es Lipschitz continua.
\end{ex}

\begin{ex}
La función de Cantor-Lebesgue es un ejemplo clásico de función en \(BV([0,1])\) pero no absolutamente continua.
\end{ex}

%\subsection{Diferenciabilidad en Casi Todo Punto y Teorema de Lebesgue}

\subsection{Teorema de aproximación de Lebesgue} % y Diferenciabilidad}

\begin{theorem}[Teorema de aproximación de Lebesgue]
Sea \(f: \mathbb{R}^n \to \mathbb{R}\) una función medible. Entonces, existe una sucesión de funciones continuas \(f_k\) tal que \(f_k \to f\) casi en todo punto y en medida.
\end{theorem}

\begin{remark}
Este teorema es fundamental para aproximar funciones no continuas por funciones continuas en el contexto de la teoría de la medida, y está estrechamente relacionado con la diferenciabilidad en casi todo punto.
\end{remark}


\subsection{Teorema de Lebesgue}

\begin{theorem}[Teorema de diferenciación de Lebesgue]
Sea \(f: \mathbb{R}^n \to \mathbb{R}\) una función localmente integrable. Entonces, para casi todo \(x \in \mathbb{R}^n\), se tiene:
\[
\lim_{r \to 0^+} \frac{1}{|B(x,r)|} \int_{B(x,r)} f(y) \, dy = f(x),
\]
donde \(B(x,r)\) es la bola de centro \(x\) y radio \(r\).
\end{theorem}

\begin{theorem}[Teorema de Lebesgue para funciones de variación acotada]
Sea \(f: [a,b] \to \mathbb{R}\) una función de variación acotada. Entonces, \(f\) es derivable en casi todo punto de \([a,b]\), y su derivada \(f'\) es integrable en \([a,b]\).
\end{theorem}

\begin{remark}
El teorema de Lebesgue asegura que, aunque una función no sea continua, su promedio local converge a su valor en casi todo punto. Esto es clave para entender la diferenciabilidad en el contexto de la teoría de la medida.
\end{remark}


\subsection{Teorema de Rademacher} %Derivabilidad de Funciones Lipschitz}

\begin{definition}[Función Lipschitz]
Una función \(f: \mathbb{R}^n \to \mathbb{R}^m\) es \textbf{Lipschitz} si existe \(L \geq 0\) tal que
\[
|f(x) - f(y)| \leq L |x - y| \quad \forall x, y \in \mathbb{R}^n.
\]
El menor \(L\) que satisface esta desigualdad se denomina \textbf{constante de Lipschitz} de \(f\).
\end{definition}

\begin{theorem}[Teorema de Rademacher]
Toda función Lipschitz \(f: \mathbb{R}^n \to \mathbb{R}^m\) es derivable en casi todo punto de \(\mathbb{R}^n\). Además, su diferencial \(Df(x)\) es esencialmente acotada (i.e., \(\|Df(x)\| \leq L\) c.t.p.).
\end{theorem}

\begin{proof}
La demostración clásica usa el teorema de diferenciación de Lebesgue y el hecho de que las funciones Lipschitz son de variación acotada en casi todo punto. %Se puede consultar en Evans y Gariepy \cite{evans1992}.
\end{proof}

\begin{remark}
El teorema de Rademacher es fundamental en el análisis de funciones no diferenciables en el sentido clásico, pero con regularidad Lipschitz. Extiende el resultado de que las funciones de variación acotada son derivables en casi todo punto (Teorema de Lebesgue).
\end{remark}

\begin{ex}
La función \(f(x) = |x|\) es Lipschitz en \(\mathbb{R}\), pero no es derivable en \(x=0\). Sin embargo, es derivable en casi todo punto de \(\mathbb{R}\) (excepto en \(x=0\)).
\end{ex}

\subsection{Resumen de Relaciones entre Propiedades}
{\footnotesize
\begin{center}
\begin{tabular}{|c|c|c|}
\hline
\textbf{Propiedad} & \textbf{Derivable en c.t.p.} & \textbf{Relación con \(L^1\)} \\
\hline
Continuidad absoluta & Sí & \(f(x) = f(a) + \int_a^x g(t) \, dt\) \\
Función Lipschitz & Sí (Rademacher) & \(\|Df\|_\infty \leq L\) \\
Variación acotada & Sí (Lebesgue) & \(f' \in L^1\) \\
Medible & No necesariamente & Aproximable por funciones continuas \\
\hline
\end{tabular}
\end{center}
}

\section{Hahn--Banach}
El \textbf{teorema de Hahn--Banach} es uno de los resultados más profundos y útiles del análisis funcional. Este teorema garantiza la existencia de extensiones lineales y continuas de funcionales definidos en subespacios de un espacio vectorial normado, preservando su norma. Sus consecuencias son vastas: permite demostrar la existencia de funcionales lineales continuos no triviales, la separación de conjuntos convexos disjuntos, y la unicidad de normas equivalentes. El teorema de Hahn--Banach es clave en la dualidad entre espacios de Banach y en el estudio de la convergencia en espacios de funciones.

\section{Espacios de Hilbert}
Los \textbf{espacios de Hilbert} son generalizaciones del espacio euclidiano a dimensiones infinitas, donde la noción de ortogonalidad y proyección juega un papel central. Estos espacios poseen una estructura geométrica rica, basada en el producto interno, y permiten desarrollar teorías de series de Fourier generalizadas, bases ortonormales y operadores compactos. Se analizan propiedades como la completitud, la descomposición ortogonal y la representación de funcionales lineales continuos mediante el teorema de representación de Riesz. Los espacios de Hilbert son fundamentales en mecánica cuántica, teoría de aproximación y procesamiento de señales.

\section{Espacios de Lebesgue}
Los \textbf{espacios de Lebesgue} \( L^p \) son espacios funcionales que generalizan el concepto de integrabilidad en el sentido de Lebesgue. Estos espacios, parametrizados por \( p \in [1, \infty] \), consisten en clases de equivalencia de funciones medibles cuyas potencias \( p \)-ésimas son integrables. Se estudian sus propiedades topológicas, la densidad de funciones continuas, y los teoremas de densidad y aproximación. Además, se exploran las desigualdades fundamentales, como las de Hölder y Minkowski, y su relación con la convergencia en medida y la continuidad de operadores lineales.

\section{Derivada débil}
La \textbf{derivada débil} extiende el concepto clásico de derivada a funciones que no son diferenciables en el sentido usual, pero que poseen una derivada generalizada en un sentido distribucional. Esta noción es esencial en la teoría de distribuciones y en el estudio de ecuaciones diferenciales parciales, donde muchas soluciones no son clásicas. Se definen los espacios de Sobolev y se analizan sus propiedades, incluyendo los teoremas de inmersión y traza, que permiten entender cómo las propiedades de regularidad de una función se transfieren a su frontera. La derivada débil es una herramienta poderosa para formular problemas variacionales y analizar su solubilidad.

\section{Espacios de Sobolev}
Los \textbf{espacios de Sobolev} son espacios funcionales que combinan propiedades de diferenciabilidad y integrabilidad, siendo fundamentales en el estudio de ecuaciones diferenciales parciales y en el análisis numérico. Estos espacios consisten en funciones cuyas derivadas débiles hasta cierto orden pertenecen a espacios de Lebesgue. Se analizan sus propiedades de completitud, densidad, y los teoremas de compactación, que son cruciales para demostrar la existencia de soluciones débiles. Los espacios de Sobolev también permiten formular problemas de minimización en cálculo de variaciones y en la teoría del potencial, proporcionando un marco riguroso para el análisis de regularidad y unicidad de soluciones.

\section{Ejemplos} %de funciones y distribuciones}

\subsection{Generales}

A continuación se dan ejemplos sencillos en una dimensión (dominio habitual \((-1,1)\) o \(\mathbb{R}\)) que ilustran pertenencia o no a los espacios \(L^{2}\) y \(H^{1}\).

\begin{itemize}
    \item \textbf{Función continua que está en \(L^{2}\) pero no en \(H^{1}\):} 
    \[
    f(x)=|x|^{\alpha},\qquad \alpha=\tfrac{1}{4},\quad x\in(-1,1).
    \]
    Aquí \(f\in L^{2}((-1,1))\) porque \(\int_{-1}^{1}|x|^{2\alpha}\,dx<\infty\) (ya que \(2\alpha>-1\)), pero la derivada puntual es \(f'(x)=\alpha\,\operatorname{sgn}(x)|x|^{\alpha-1}\) cuya norma \(L^{2}\) diverge porque \(2(\alpha-1)\le -1\); por tanto no existe derivada débil en \(L^{2}\) y \(f\notin H^{1}((-1,1))\). Este ejemplo es continuo pero no pertenece a \(H^{1}\).

    \item \textbf{Función discontinua que pertenece a \(L^{2}\) pero no a \(H^{1}\):}
    \[
    \chi_{(0,1)}(x)=\begin{cases}1,&x\in(0,1),\\0,&\text{en otro caso},\end{cases}
    \]
    vista como función en \(\mathbb{R}\). Tiene norma \(L^{2}\) finita, pero presenta una discontinuidad por salto; su derivada en sentido de distribuciones contiene deltas de frontera (no representables por una función en \(L^{2}\)), luego \(\chi_{(0,1)}\notin H^{1}(\mathbb{R})\).

    \item \textbf{Función o distribución que no está en \(L^{2}\):}
    \begin{itemize}
        \item La función singular \(g(x)=|x|^{-\beta}\) en \((-1,1)\) con \(\beta\geq\tfrac{1}{2}\) no pertenece a \(L^{2}((-1,1))\) porque \(\int_{-1}^{1}|x|^{-2\beta}\,dx=\infty\).
        \item La delta de Dirac \(\delta_{0}\) (distribución) no es una función en \(L^{2}\) (no es una función cuadrado-integrable).
    \end{itemize}
\end{itemize}

Cada ejemplo indica la razón básica (comportamiento en el origen o saltos) por la que falla la integrabilidad cuadrática o la existencia de derivada débil en \(L^{2}\).

\subsection{Indicador}

La función $\mathrm{1}_{\{0,1\}}(x)$ representa la \textbf{indicador de $\{0,1\}$}, definida como:
\[
\mathrm{1}_{\{0,1\}}(x) =
\begin{cases}
1 & \text{si } x \in \{0,1\}, \\
0 & \text{en otro caso}.
\end{cases}
\]
Esta función es discontinua en $x=0$ y $x=1$, por lo que no pertenece a los espacios clásicos de Sobolev $H^s(\mathbb{R})$ para $s \geq 1/2$ en el sentido estándar. Sin embargo, su regularidad puede analizarse en el marco de los \textbf{espacios de Sobolev fraccionarios} $H^s(\mathbb{R})$, $s \in (0,1)$.

En particular, la indicatriz $\mathrm{1}_{\{0,1\}}$ pertenece al espacio $H^{1/2 - \varepsilon}(\mathbb{R})$ para todo $\varepsilon > 0$, pero \textbf{no} pertenece a $H^{1/2}(\mathbb{R})$. Este resultado se justifica mediante el cálculo de la norma de Sobolev fraccionaria:
\[
\| \mathrm{1}_{\{0,1\}} \|_{H^s(\mathbb{R})}^2
= \int_{\mathbb{R}} |\xi|^{2s} |\widehat{\mathrm{1}_{\{0,1\}}}(\xi)|^2 \, d\xi,
\]
donde $\widehat{\mathrm{1}_{\{0,1\}}}(\xi)$ es la transformada de Fourier de la indicatriz. Dado que la transformada de Fourier de $\mathrm{1}_{\{0,1\}}$ decae como $|\xi|^{-1}$ en el infinito, la integral converge si y solo si $s < 1/2$. Por lo tanto, la regularidad fraccionaria de la indicatriz está limitada por $s = 1/2$.
La \textbf{transformada de Fourier} de la indicatriz \(\mathrm{1}_{\{0,1\}}(x)\) es:
\[
\widehat{\mathrm{1}_{\{0,1\}}}(\xi) = e^{-i\xi} \frac{\sin(\xi)}{\xi},
\]
que también puede escribirse como:
\[
\widehat{1_{\{0,1\}}}(\xi) = e^{-i\xi} \frac{\sin(\xi)}{\xi}.
\]
Esta función decae como \(|\xi|^{-1}\) cuando \(|\xi| \to \infty\).

\section{M\'as espacios funcionales}

\subsection{Espacios de Sobolev \(H^s(\R)\)}
El espacio de Sobolev \(H^s(\R)\) para \(s \in \mathbb{R}\) se define como:
\[
H^s(\R) = \left\{ f \in \mathcal{S}'(\R) : \int_{\R} (1 + |\xi|^2)^s |\widehat{f}(\xi)|^2 \, d\xi < \infty \right\},
\]
donde \(\widehat{f}\) es la transformada de Fourier de \(f\). La norma en \(H^s(\R)\) es:
\[
\|f\|_{H^s(\R)}^2 = \int_{\R} (1 + |\xi|^2)^s |\widehat{f}(\xi)|^2 \, d\xi.
\]

\subsection{Espacios de Besov \(\bes_{p,q}^s(\R)\)}
El espacio de Besov \(\bes_{p,q}^s(\R)\) para \(s \in \mathbb{R}\) y \(p, q \in [1, \infty]\) se define mediante diferencias finitas o mediante la descomposición en frecuencias. Intuitivamente, mide la regularidad de \(f\) en términos de cómo sus diferencias finitas decaen. Para \(0 < s < 1\) y \(1 \leq p, q \leq \infty\), se tiene:
\[
\bes_{p,q}^s(\R) = \left\{ f \in L^p(\R) : \left( \int_0^\infty \left( t^{-s} \omega_p(f, t) \right)^q \frac{dt}{t} \right)^{1/q} < \infty \right\},
\]
donde \(\omega_p(f, t)\) es el módulo de continuidad en \(L^p\).

\subsection{Exponente de Hölder}
El \textbf{exponente de Hölder} \(\alpha\) de una función \(f\) en un punto \(x_0\) es el supremo de todos \(\alpha > 0\) tales que:
\[
|f(x) - f(x_0)| \leq C |x - x_0|^\alpha \quad \text{para } x \text{ cerca de } x_0.
\]

%\section{Ejemplos de Fractales 1D y su Regularidad}
\section{Fractales}

\subsection{Función de Weierstrass}
La \textbf{función de Weierstrass} se define para \(\lambda > 1\) y \(0 < H < 1\) como:
\[
W_H(x) = \sum_{n=0}^{\infty} \lambda^{-nH} \cos(2\pi \lambda^n x).
\]
\begin{itemize}
    \item \textbf{Propiedades:} Es continua pero nowhere differentiable.
    \item \textbf{Regularidad de Sobolev:} \(W_H \in H^s(\R)\) para todo \(s < H\).
    \item \textbf{Regularidad de Besov:} \(W_H \in \bes_{p,q}^H(\R)\) para todo \(p, q \geq 1\).
    \item \textbf{Exponente de Hölder:} \(\alpha = H\) en casi todo punto.
    \item \textbf{Referencia:} Mandelbrot \cite{mandelbrot1982fractal}.
\end{itemize}

\subsection{Función de Takagi} %(Escalera del Diablo)}
La \textbf{función de Takagi} se define como:
\[
T(x) = \sum_{n=0}^{\infty} \frac{1}{2^n} \dist(2^n x, \Z),
\]
donde \(\dist(x, \Z)\) es la distancia de \(x\) al entero más cercano.
\begin{itemize}
    \item \textbf{Propiedades:} Es continua, creciente, y nowhere differentiable.
    \item \textbf{Regularidad de Sobolev:} \(T \in H^s(\R)\) para todo \(s < 1/2\).
    \item \textbf{Regularidad de Besov:} \(T \in \bes_{p,q}^{1/2}(\R)\) para todo \(p, q \geq 1\).
    \item \textbf{Exponente de Hölder:} \(\alpha = 1/2\) en casi todo punto.
    \item \textbf{Referencia:} Takagi \cite{takagi1903}.
\end{itemize}

\subsection{Movimiento browniano fraccionario}
El \textbf{movimiento Browniano fraccionario} \(B_H(t)\) con parámetro de Hurst \(0 < H < 1\) es un proceso estocástico gaussiano con trayectorias fractales.
\begin{itemize}
    \item \textbf{Propiedades:} Tiene trayectorias continuas pero nowhere diferenciables si \(H \neq 1\).
    \item \textbf{Regularidad de Sobolev:} Casi todas las trayectorias satisfacen \(B_H \in H^s(\R)\) para todo \(s < H\).
    \item \textbf{Regularidad de Besov:} Casi todas las trayectorias satisfacen \(B_H \in \bes_{p,q}^H(\R)\) para todo \(p, q \geq 1\).
    \item \textbf{Exponente de Hölder:} \(\alpha = H\) en casi todo punto y trayectoria.
    \item \textbf{Referencia:} Samorodnitsky y Taqqu \cite{samorodnitsky1994}.
\end{itemize}

\subsection{Conjunto de Cantor}
El \textbf{conjunto de Cantor} \(C\) se construye eliminando el tercio central de cada intervalo.
\begin{itemize}
    \item \textbf{Propiedades:} Es un fractal autosimilar con dimensión de Hausdorff \(D = \frac{\log 2}{\log 3}\).
    \item \textbf{Función indicatriz:} La función \(\mathrm{1}_C(x)\) tiene regularidad de Besov \(s < D\).
    \item \textbf{Regularidad de Sobolev:} \(\mathrm{1}_C \notin H^{1/2}(\R)\).
    \item \textbf{Referencia:} Cantor \cite{cantor1883}.
\end{itemize}

\subsection{Función de Cantor-Lebesgue}
La \textbf{función de Cantor-Lebesgue} \(f_C(x)\) es continua, creciente, y singular (constante en los intervalos del complemento de \(C\)).
\begin{itemize}
    \item \textbf{Propiedades:} Es constante en los intervalos de \(\R \setminus C\).
    \item \textbf{Regularidad de Sobolev:} \(f_C \notin H^{1/2}(\R)\).
    \item \textbf{Regularidad de Besov:} \(f_C \in \bes_{p,q}^s(\R)\) para \(s < D\).
    \item \textbf{Referencia:} Lebesgue \cite{lebesgue1904}.
\end{itemize}

\section{Ecuaciones elípticas}
Las ecuaciones elípticas son un tipo de ecuaciones en derivadas parciales de segundo orden cuyo símbolo principal es definido positivo (por ejemplo, la ecuación de Laplace $\Delta u=0$ o la de Poisson $-\Delta u=f$), y modelan estados estacionarios y distribuciones espaciales equilibradas; en problemas de contorno bien planteados suelen garantizar existencia, unicidad y regularidad de soluciones suaves en dominios regulares, y poseen propiedades cualitativas importantes como el principio del máximo, estimaciones a priori y dependencia continua respecto de los datos, lo que las convierte en herramientas centrales en análisis y aplicaciones físicas e ingenieriles.

\chapter{Ejemplos para minimizar con descenso de gradiente}

Este presenta cuatro ejemplos \textbf{concretos y precisos} de funcionales diferenciables que pueden minimizarse utilizando el método de \textbf{descenso de gradiente}. Los ejemplos abarcan:
\begin{itemize}
    \item Una función fuertemente convexa (regresión Ridge).
    \item Un problema físico (energía potencial de un sistema de resortes).
    \item Un problema de optimización de forma (área de un polígono bajo restricción de perímetro).
    \item Un ejemplo de aprendizaje automático (regresión logística).
\end{itemize}
Todos los ejemplos están formulados con un número reducido de variables (entre 2 y 10) para facilitar su comprensión y implementación.

\section{Ridge} 
%Función Fuertemente Convexa: Regresión Lineal con Regularización L2 (Ridge)}
\subsection{Contexto}
En aprendizaje automático, la regresión Ridge busca minimizar el error cuadrático entre las predicciones de un modelo lineal y los datos reales, añadiendo un término de regularización L2 para evitar el sobreajuste.

\subsection{Funcional}
Dada una matriz de datos $ \mathbf{X} \in \mathbb{R}^{n \times d} $, un vector de etiquetas $ \mathbf{y} \in \mathbb{R}^n $, y un vector de pesos $ \mathbf{w} \in \mathbb{R}^d $, la funcional a minimizar es:
$
J(\mathbf{w}) = \frac{1}{2n} \| \mathbf{X}\mathbf{w} - \mathbf{y}\|_2^2 + \lambda \|\mathbf{w}\|_2^2
$
donde $ \lambda > 0 $ es el parámetro de regularización.

\subsection{Propiedades}
\begin{itemize}
    \item \textbf{Fuertemente convexa} si $ \lambda > 0 $.
    \item Diferenciable, con gradiente:
    $
    \nabla J(\mathbf{w}) = \frac{1}{n} \mathbf{X}^T (\mathbf{X}\mathbf{w} - \mathbf{y}) + 2\lambda \mathbf{w}
    $.
\end{itemize}

\subsection{Ejemplo}
Para $ d = 2 $ (dos variables $ w_1, w_2 $) y $ n = 3 $ muestras:
$$
\mathbf{X} = \begin{bmatrix} 1 & 2  3 & 4  5 & 6 \end{bmatrix}, \quad \mathbf{y} = \begin{bmatrix} 3  7  11 \end{bmatrix}, \quad \lambda = 0.1.
$$
La funcional se escribe como: \scriptsize
$$
J(w_1, w_2) = \frac{1}{6} \left( (w_1 + 2w_2 - 3)^2 + (3w_1 + 4w_2 - 7)^2 + (5w_1 + 6w_2 - 11)^2 \right) + 0.1(w_1^2 + w_2^2).
$$
\normalsize


\section{Energía potencial de un sistema de resortes}
\subsection{Contexto}
Minimización de la energía potencial elástica de un sistema de masas conectadas por resortes en 1D. Este es un problema clásico en física de medios continuos.

\subsection{Funcional}
Consideremos $ N $ masas conectadas por resortes. Sea $ x_i $ la posición de la masa $ i $ (con $ x_0 $ y $ x_{N+1} $ fijas). La energía potencial total es:
$
E(\mathbf{x}) = \sum_{i=1}^{N} \frac{k}{2} (x_i - x_{i-1} - L_i)^2
$
donde $ k $ es la constante elástica, $ L_i $ es la longitud en reposo del resorte $ i $, y $ \mathbf{x} = [x_1, \dots, x_N] $.

\subsection{Propiedades}
\begin{itemize}
    \item \textbf{Convexa} (suma de cuadrados).
    \item Diferenciable, con gradiente:
    $
    \frac{\partial E}{\partial x_i} = k \left( 2x_i - x_{i-1} - x_{i+1} - L_i + L_{i+1} \right)
    $
    (con $ x_0 $ y $ x_{N+1} $ fijas).
\end{itemize}

\subsection{Ejemplo}
Para $ N = 2 $ masas (2 variables $ x_1, x_2 $):
$
k = 1, \quad L_1 = L_2 = 1, \quad x_0 = 0, \quad x_3 = 2
$
La funcional se escribe como:
$
E(x_1, x_2) = \frac{1}{2}(x_1 - 1)^2 + \frac{1}{2}(x_2 - x_1 - 1)^2 + \frac{1}{2}(2 - x_2 - 1)^2.
$

\section{Optimización de forma} %: Área de un Triángulo bajo Restricción de Perímetro}
\subsection{Contexto}
Encontrar las coordenadas de los vértices de un triángulo que minimice su área, sujetas a una restricción de perímetro fijo. Este es un problema simplificado de optimización de forma.

\subsection{Funcional}
Para un triángulo con vértices $ \mathbf{p}_1 = (x_1, y_1) $, $ \mathbf{p}_2 = (x_2, y_2) $, $ \mathbf{p}_3 = (x_3, y_3) $, y un perímetro $ P $ fijo:
\begin{itemize}
    \item \textbf{Área} (a minimizar):
    $
    A(\mathbf{p}_1, \mathbf{p}_2, \mathbf{p}_3) = \frac{1}{2}  (\mathbf{p}_2 - \mathbf{p}_1) \times (\mathbf{p}_3 - \mathbf{p}_1) 
    $
    (producto vectorial en 2D).
    \item \textbf{Restricción}:
    $
    \mathbf{p}_2 - \mathbf{p}_1 + \mathbf{p}_3 - \mathbf{p}_2 + \mathbf{p}_1 - \mathbf{p}_3 = P.
    $
\end{itemize}

\subsection{Penalización}
Para transformar el problema restringido en uno no restringido, se usa:
$$
J(\mathbf{p}_1, \mathbf{p}_2, \mathbf{p}_3) = A(\mathbf{p}_1, \mathbf{p}_2, \mathbf{p}_3) + \mu \left( \mathbf{p}_2 - \mathbf{p}_1 + \mathbf{p}_3 - \mathbf{p}_2 + \mathbf{p}_1 - \mathbf{p}_3 - P \right)^2
$$
donde $ \mu $ es un parámetro de penalización.

\subsection{Ejemplo}
Para $ P = 6 $, y fijando $ \mathbf{p}_1 = (0, 0) $, $ \mathbf{p}_2 = (a, 0) $, $ \mathbf{p}_3 = (b, c) $:
\begin{itemize}
    \item Área:
    $
    A(a, b, c) = \frac{1}{2} |a \cdot c|.
    $
    \item Restricción:
    $
    a + \sqrt{(b - a)^2 + c^2} + \sqrt{b^2 + c^2} = 6.
    $
\end{itemize}

\section{Aprendizaje automático} %Log-Verosimilitud para Regresión Logística Binaria}
\subsection{Contexto}
Clasificar puntos en 2D en dos clases (0 o 1) aprendiendo los parámetros de un modelo logístico.

\subsection{Funcional}
Dada una matriz de datos $ \mathbf{X} \in \mathbb{R}^{n \times 2} $, un vector de etiquetas $ \mathbf{y} \in 0, 1^n $, y parámetros $ \mathbf{w} = (w_1, w_2) $, $ b $ (pesos y sesgo), la log-verosimilitud negativa (a minimizar) es:
$$
J(\mathbf{w}, b) = -\sum_{i=1}^n \left[ y_i \log(\sigma(\mathbf{w}^T \mathbf{x}_i + b)) + (1 - y_i) \log(1 - \sigma(\mathbf{w}^T \mathbf{x}_i + b)) \right]
$$
donde $ \sigma(z) = \frac{1}{1 + e^{-z}} $ es la función sigmoide.

\subsection{Propiedades}
\begin{itemize}
    \item \textbf{Convexa} (la log-verosimilitud negativa es convexa para regresión logística).
    \item Diferenciable, con gradiente:
    $$
    \nabla_{w_j} J = -\sum_{i=1}^n (y_i - \sigma(\mathbf{w}^T \mathbf{x}i + b)) x_{ij}, \quad \nabla_b J = -\sum_{i=1}^n (y_i - \sigma(\mathbf{w}^T \mathbf{x}_i + b)).
    $$
\end{itemize}

\subsection{Ejemplo}
Para $ n = 3 $ muestras:
$$
\mathbf{X} = \begin{bmatrix} 1 & 2  2 & 1  3 & 3 \end{bmatrix}, \quad \mathbf{y} = \begin{bmatrix} 1  0  1 \end{bmatrix}.
$$
La funcional se escribe como: \scriptsize
$$
J(w_1, w_2, b) = -\left[ \log(\sigma(w_1 + 2w_2 + b)) + \log(1 - \sigma(2w_1 + w_2 + b)) + \log(\sigma(3w_1 + 3w_2 + b)) \right].
$$
\normalsize




%\addcontentsline{toc}{chapter}{Bibliograf\'ia}
%\nocite{*}
%{\footnotesize
%\bibliographystyle{siam}
%\bibliography{biblio}
%}

%\addcontentsline{toc}{chapter}{\'Indice alfab\'etico}
%{\scriptsize \printindex}

%\end{document} 


%\addcontentsline{toc}{chapter}{Bibliograf\'ia}
%\nocite{*}
{\footnotesize
\bibliographystyle{siam}
\bibliography{biblio}
}

%\addcontentsline{toc}{chapter}{\'Indice alfab\'etico}
{\scriptsize \printindex}

\end{document} 
